\documentclass[pdflatex,sn-mathphys-num]{sn-jnl}% Math and Physical 
\usepackage{graphicx}%
\usepackage{tabularx}
\usepackage{multirow}%
\usepackage{array}
\usepackage{amsmath,amssymb,amsfonts}%
\usepackage{amsthm}%
\usepackage{mathrsfs}%
\usepackage[title]{appendix}%
\usepackage{xcolor}%
\usepackage{textcomp}%
\usepackage{manyfoot}%
\usepackage{booktabs}%
\usepackage{algorithm}%
\usepackage{algorithmicx}%
\usepackage{algpseudocode}%
\usepackage{listings}%
\usepackage{booktabs}
\usepackage{adjustbox}
\usepackage{amsmath,amssymb,amsfonts}
\usepackage{graphicx}
\usepackage{booktabs}
\usepackage{natbib}
\usepackage{adjustbox}
\usepackage{booktabs}
\usepackage{graphicx}
\usepackage{float}
\usepackage{subcaption}
\usepackage{algorithm}
\usepackage{algpseudocode}

\theoremstyle{thmstyleone}%
\newtheorem{theorem}{Theorem}%  meant for continuous numbers

\newtheorem{proposition}[theorem]{Proposition}% 
%%\newtheorem{proposition}{Proposition}% to get separate numbers for theorem and proposition etc.

\theoremstyle{thmstyletwo}%
\newtheorem{example}{Example}%
\newtheorem{remark}{Remark}%

\theoremstyle{thmstylethree}%
\newtheorem{definition}{Definition}%

\raggedbottom
%%\unnumbered% uncomment this for unnumbered level heads

\begin{document}

\title[GNN-DHRS]
{An Explainable Graph Attention Framework with Dual Hybrid Resampling for Imbalanced Brain Stroke Prediction}

% \author*[1,2]{\fnm{First} \sur{Author}}\email{iauthor@gmail.com}

% \author[2,3]{\fnm{Second} \sur{Author}}\email{iiauthor@gmail.com}
% \equalcont{These authors contributed equally to this work.}

% \author[1,2]{\fnm{Third} \sur{Author}}\email{iiiauthor@gmail.com}
% \equalcont{These authors contributed equally to this work.}

% \affil*[1]{\orgdiv{Department}, \orgname{Organization}, \orgaddress{\street{Street}, \city{City}, \postcode{100190}, \state{State}, \country{Country}}}

% \affil[2]{\orgdiv{Department}, \orgname{Organization}, \orgaddress{\street{Street}, \city{City}, \postcode{10587}, \state{State}, \country{Country}}}

% \affil[3]{\orgdiv{Department}, \orgname{Organization}, \orgaddress{\street{Street}, \city{City}, \postcode{610101}, \state{State}, \country{Country}}}

%%==================================%%
%% Sample for unstructured abstract %%
%%==================================%%


\abstract{
Early brain stroke classification is challenging because clinical datasets
often exhibit severe class imbalance, noisy minority distributions, complex
relationships among patient profiles, and limited model interpretability.
This paper proposes \textbf{GNN-DHRS}, an explainable graph neural network
framework incorporating a Dual Hybrid Resampling Strategy for imbalanced
stroke classification. The framework combines Borderline-SMOTE with
SMOTEENN to enhance minority representation and remove locally inconsistent
observations before constructing a patient-similarity graph. A multi-head
Graph Attention Network subsequently learns relation-aware patient
representations, while feature-attribution and attention-based analyses
support interpretation at the feature and patient-relation levels.

{
The framework was evaluated on three imbalanced stroke benchmark datasets,
DF-1, DF-2, and DF-3, using nested stratified cross-validation and strictly
inductive evaluation. All reported measures were calculated exclusively from
untouched outer-test predictions.
\textcolor{red}{
Across five predefined stochastic repetitions, GNN-DHRS achieved
seed-aggregated Recall values of \(99.28\%\), \(98.91\%\), and \(99.07\%\),
F1-scores of \(99.35\%\), \(98.32\%\), and \(98.45\%\), and PR-AUC values
of \(99.84\%\), \(98.65\%\), and \(98.79\%\) on DF-1, DF-2, and DF-3,
respectively.
}
Additional record-integrity, cross-benchmark transfer, calibration,
robustness, and selective-prediction analyses provided complementary evidence
regarding the reliability and stability of the framework under the evaluated
conditions.
}
The results indicate consistently strong within-benchmark
stroke-status classification performance under severe class imbalance, while
the explainability analyses produced feature-attribution patterns broadly
consistent with established stroke-associated factors and identified
influential patient relationships within the fitted model.
\textcolor{red}{
These findings support further evaluation of GNN-DHRS on independently
collected and prospective clinical cohorts before conclusions regarding
external generalizability or clinical deployment can be established.
}

}

\keywords{
Brain Stroke Prediction,
Graph Neural Networks,
Graph Attention Networks,
Imbalanced Learning,
Dual Hybrid Resampling,
Borderline-SMOTE,
SMOTEENN,
Explainable Artificial Intelligence,
Healthcare Analytics,
Clinical Decision Support,
Medical Data Mining,
Deep Learning for Healthcare
}

\maketitle


% Introduction
% =========================================================
\section{Introduction}
\label{sec:introduction}

Brain stroke remains one of the leading causes of mortality and long-term
disability worldwide, imposing substantial neurological, cognitive, social,
and economic burdens.  {
In this study, the computational task is retrospective binary stroke-status
classification from structured demographic, medical-history, and lifestyle
variables. Because the evaluated datasets do not contain longitudinal
follow-up or verified future outcomes, ``prediction'' refers to predicting
the recorded stroke-status label of an unseen patient profile rather than
prospective stroke incidence or time-to-event prediction.
}
This task remains challenging because clinical datasets are often
heterogeneous, noisy, highly imbalanced, and characterized by nonlinear
interactions among demographic, metabolic, cardiovascular, and
lifestyle-related variables
\cite{yang2020early,johnson2023stroke,melnykova2025stroke}.

Machine-learning and deep-learning approaches have increasingly been applied
to structured stroke data because they can model nonlinear relationships and
learn discriminative patterns from heterogeneous patient records
\cite{tashkova2025comparative,akib2025optimizing,kanwal2025stroke}.
Nevertheless, several limitations remain. First, stroke-positive cases usually
represent only a small fraction of the available data, causing conventional
classifiers to favor the majority class and reducing minority-case detection
\cite{he2009learning,haixiang2017learning,fernandez2018learning,
abousaber2025robust}. Second, commonly used imbalance-learning methods,
including SMOTE, Borderline-SMOTE, ADASYN, SMOTEENN, and SMOTETomek, may
increase minority representation but can also introduce noisy synthetic
samples or class overlap
\cite{chawla2002smote,han2005borderline,he2008adasyn,batista2004study}.
Third, most conventional stroke-classification models process patients
independently and therefore do not explicitly exploit latent inter-patient
relationships.

Graph Neural Networks (GNNs) provide a mechanism for modeling such
relationships by representing samples as nodes and learning from neighborhood
interactions \cite{kipf2017semi,wu2021comprehensive,li2021graph}. Graph
Attention Networks (GATs) further assign adaptive importance to neighboring
nodes during aggregation \cite{velickovic2018graph}. However, directly
applying graph learning to severely imbalanced clinical data can produce weak
minority connectivity and unreliable neighborhoods. Consequently, imbalance
handling and patient-graph construction should be considered jointly.

 {
Motivated by these limitations, this paper proposes \textbf{GNN-DHRS}, an
explainable graph-attention framework for imbalanced stroke-status
classification. Its principal contribution is the ordered, training-only
integration of Borderline-SMOTE and SMOTEENN before patient-similarity graph
construction. Borderline-SMOTE reinforces difficult minority regions, whereas
SMOTEENN subsequently performs additional resampling and neighborhood
cleaning. The novelty therefore lies not in modifying either resampling method
individually, but in evaluating their ordered interaction with patient-graph
construction. The proposed sequence is compared with its individual
components and reversed ordering, and its effects are examined in terms of
class composition, graph topology, and downstream predictive performance.
The refined records are then represented through a cosine-similarity graph,
processed by multi-head graph attention, and interpreted using complementary
feature- and relation-level explanations.
}

The main contributions of this study are summarized as follows:

\begin{itemize}

\item \textbf{Graph-oriented dual hybrid resampling:}
A sequential Borderline-SMOTE--SMOTEENN strategy refines minority
representation before graph construction.
 {
Its contribution is evaluated through component, ordering, sample-composition,
and graph-topology analyses, distinguishing the proposed sequence from its
constituent resampling methods.
}

\item \textbf{Graph-attention relational modeling:}
A cosine-similarity patient graph represents relationships among similar
records, while multi-head graph attention performs adaptive neighborhood
aggregation for stroke-status classification.

\item \textbf{Complementary explainability:}
Feature-attribution, attention-based relational interpretation, and
edge-ablation analyses provide feature- and patient-relation-level evidence
for examining model behavior.

\item \textbf{Comprehensive leakage-safe validation:}
The framework is evaluated on three imbalanced stroke benchmark datasets
using nested cross-validation, baseline comparisons, ablation studies,
sensitivity analysis, explainability experiments, and statistical validation.
 {
Additional analyses examine strictly inductive inference, record integrity,
cross-benchmark transfer, calibration, robustness, and selective prediction.
}

\end{itemize}

The remainder of this paper is organized as follows.
Section~\ref{sec:relatedwork} reviews related work.
Section~\ref{sec:methodology} presents the proposed GNN-DHRS framework.
Section~\ref{sec:results} reports the experimental results and validation.
Section~\ref{sec:discussion} discusses the findings and limitations.
Section~\ref{sec:conclusion} concludes the paper and outlines future work.
%====================================================
% Related Work
% =========================================================
\section{Related Work}
\label{sec:relatedwork}

Automated stroke classification has been widely investigated using
machine-learning, deep-learning, imbalance-learning, graph-based, and
explainable-AI approaches. This section reviews the studies most relevant to
the proposed integration of imbalance handling, patient-similarity graph
construction, graph attention, and explainability.

% =========================================================
\textbf{Machine Learning and Deep Learning for Stroke Prediction: }
Traditional machine-learning methods, including Decision Trees, Random
Forests, Support Vector Machines, Logistic Regression, XGBoost, and LightGBM,
have been widely applied to structured stroke data
\cite{johnson2023stroke,tashkova2025comparative,melnykova2025stroke}.
Deep-learning approaches have also been investigated to improve nonlinear
representation learning \cite{akib2025optimizing,kanwal2025stroke}.
Nevertheless, severe class imbalance remains a major limitation, and most
conventional approaches process patient records independently without
explicitly modeling inter-patient relationships.

% =========================================================
\textbf{Imbalance Learning for Medical Prediction:}
Class imbalance is particularly important in medical classification because
disease-positive cases often represent only a small fraction of the available
data \cite{he2009learning,haixiang2017learning}. Common approaches include
Random Oversampling, Random Undersampling, SMOTE, Borderline-SMOTE, ADASYN,
SMOTEENN, and SMOTETomek
\cite{chawla2002smote,han2005borderline,he2008adasyn,batista2004study,
batista2003balancing,abousaber2025robust}. Borderline-SMOTE concentrates
synthetic generation near difficult minority boundaries, whereas SMOTEENN
combines oversampling with neighborhood-based cleaning. However, these
methods are generally evaluated independently of graph construction, leaving
their effects on patient-neighborhood topology insufficiently examined.

% =========================================================
\textbf{Graph Neural Networks and Explainable Healthcare Intelligence: }

Graph Neural Networks (GNNs) model relational dependencies by aggregating
information from neighboring nodes \cite{kipf2017semi,wu2021comprehensive},
while Graph Attention Networks (GATs) learn adaptive neighbor-specific
importance weights \cite{velickovic2018graph}.

 {
Clinical graph-learning studies have adopted several forms of patient
relationships. Parisot et al.~\cite{parisot2018disease} used phenotypic
population graphs for neurological disease classification, Sun
et al.~\cite{sun2021disease} combined patient records with medical knowledge
graphs, and Ghorbani et al.~\cite{ghorbani2022ragcn} addressed imbalanced
medical node classification through adversarial reweighting. Tariq
et al.~\cite{tariq2024transfusion} constructed temporal patient-similarity
graphs from heterogeneous EHR data. More recent studies include adaptive
patient-similarity graphs for post-myocardial-infarction outcomes
\cite{guo2026myocardial}, multimodal GAT modeling for first ischemic stroke
\cite{desai2026grasp}, and heterogeneous patient--diagnosis--medication
graphs \cite{cammarasana2026drug}. These studies demonstrate the value of
clinical relational learning, but they do not examine the ordered interaction
of minority-boundary enhancement, neighborhood cleaning, and subsequent
patient-similarity graph construction considered in GNN-DHRS.
}

\begin{table*}[!ht]
\centering
\caption{Comparison of representative graph-based healthcare prediction methods.}
\label{tab:related_graph_healthcare}

\begingroup

\fontsize{7.2}{8.2}\selectfont
\setlength{\tabcolsep}{1.8pt}
\renewcommand{\arraystretch}{1.06}

\begin{tabularx}{\textwidth}{@{}
>{\raggedright\arraybackslash\hsize=0.88\hsize}X
>{\raggedright\arraybackslash\hsize=0.66\hsize}X
>{\raggedright\arraybackslash\hsize=0.78\hsize}X
>{\raggedright\arraybackslash\hsize=0.98\hsize}X
>{\centering\arraybackslash\hsize=0.58\hsize}X
>{\raggedright\arraybackslash\hsize=0.85\hsize}X
>{\raggedright\arraybackslash\hsize=0.92\hsize}X
@{}}
\toprule

\textbf{Study} &
\textbf{Data} &
\shortstack{\textbf{Imbalance}\\\textbf{method}} &
\shortstack{\textbf{Graph}\\\textbf{type}} &
\textbf{Attn.} &
\textcolor{red}{\textbf{XAI}} &
\shortstack{\textbf{Main}\\\textbf{limitation}} \\
\midrule

Parisot et al.~\cite{parisot2018disease}
& ABIDE; ADNI
& None
& Phenotypic graph
& No
& \textcolor{red}{NR}
& Transductive imaging \\

Sun et al.~\cite{sun2021disease}
& Clinical EHR
& None
& Patient--KG
& GAT
& \textcolor{red}{NR}
& External knowledge \\

Ghorbani et al.~\cite{ghorbani2022ragcn}
& Medical data
& Adv. reweighting
& Population graph
& No
& \textcolor{red}{NR}
& No data resampling \\

Tariq et al.~\cite{tariq2024transfusion}
& Multi-site EHR
& Propensity match.
& Temporal graph
& No
& \textcolor{red}{NR}
& Matched cohort \\

Guo et al.~\cite{guo2026myocardial}
& MI data
& Not central
& Adaptive \(k\)-NN
& Cross-modal
& \textcolor{red}{NR}
& Post-MI only \\

Desai et al.~\cite{desai2026grasp}
& UK Biobank
& Model-based
& Multimodal graph
& GAT
& \textcolor{red}{NR}
& Multimodal data req. \\

Cammarasana and Patanè~\cite{cammarasana2026drug}
& MIMIC-III
& None
& Heterogeneous graph
& No
& \textcolor{red}{NR}
& Medication-link task \\

\midrule

\textbf{GNN-DHRS}
& \textbf{DF-1--DF-3}
& \textbf{DHRS}
& \textbf{Cosine \(k\)-NN}
& \textbf{MH-GAT}
& \textcolor{red}{\textbf{SHAP + Attn.}}
& \textbf{No prospective validation} \\

\bottomrule
\end{tabularx}

\vspace{2pt}

\begin{minipage}{0.98\textwidth}
\fontsize{6.8}{7.8}\selectfont
\textit{Note:}
EHR: electronic health record;
KG: knowledge graph;
MI: myocardial infarction;
Attn.: attention;
MH-GAT: multi-head Graph Attention Network;
XAI: explainable artificial intelligence;
NR: not explicitly reported;
DHRS: Dual Hybrid Resampling Strategy;
SHAP: SHapley Additive exPlanations.
\end{minipage}

\endgroup
\end{table*}
\textcolor{red}{
Table~\ref{tab:related_graph_healthcare} compares representative clinical
graph-learning approaches in terms of dataset, imbalance handling, graph
construction, attention mechanism, explainability, and principal limitation.
In addition to investigating how ordered minority enhancement and neighborhood
cleaning affect patient-similarity graph construction, GNN-DHRS incorporates
complementary feature-level Kernel SHAP attribution and attention-based
relational interpretation within the same framework.
}

Explainable artificial intelligence is also important in healthcare because
predictive models should provide interpretable evidence for their outputs
\cite{samek2021explainable,arrieta2020explainable,el2025explainable}.
SHAP provides feature-level attribution \cite{lundberg2017unified}, whereas
attention mechanisms can provide complementary information about influential
graph relations. However, most healthcare studies emphasize feature-level
explanations and provide comparatively limited validation of patient-to-patient
relational explanations.

% =========================================================
\textbf{Research Gap and Motivation: }
 {
The reviewed literature reveals three gaps relevant to the present study:
(i) relatively few clinical GNN approaches jointly address severe class
imbalance and patient-graph quality; (ii) imbalance-aware graph methods
typically emphasize loss or sample reweighting rather than measuring how
data-level resampling changes minority connectivity and graph topology; and
(iii) many clinical graph models rely on transductive, longitudinal, or
multimodal settings that differ from the cross-sectional stroke benchmarks
considered here. These gaps motivate a leakage-safe and strictly inductive
framework that integrates graph-oriented resampling, patient-similarity
construction, graph attention, and quantitative explainability.
}

To address these gaps, GNN-DHRS combines ordered
Borderline-SMOTE--SMOTEENN refinement with patient-similarity graph
construction, multi-head graph attention, and complementary feature- and
relation-level interpretation for imbalanced stroke-status classification.
% =========================================================
%--------------------
% Datasets
%--------------------
\section{Datasets}
\label{sec:datasets}

This study evaluates GNN-DHRS using three publicly available stroke benchmark
datasets, denoted DF-1, DF-2, and DF-3
\cite{dataset_df1,dataset_df2,dataset_df3}.  {
These files are treated as benchmark datasets rather than independent
prospective clinical cohorts. In particular, DF-2 is a completed,
preprocessed version of the commonly distributed stroke-prediction data
represented by DF-3, and the available repository information does not
establish their independence.
}

\textcolor{red}{
Because DF-2 does not contain the row identifier available in DF-3 and the
repositories do not provide a shared verified patient identifier or original
record linkage, definitive patient-level overlap assessment between DF-2 and
DF-3 is not possible. Exact cross-file matching can identify only records that
remain identical across the available predictor--target fields; it cannot
establish whether non-identical rows correspond to the same underlying
individual after preprocessing, imputation, completion, or other record
transformation. Possible latent overlap between these two public benchmark
files therefore cannot be excluded.
}

The datasets contain demographic, medical-history, and lifestyle variables
and exhibit substantial class imbalance.
% =========================================================
% =========================================================
\subsection{Dataset Overview and Record-Integrity Assessment}
\label{subsec:dataset_overview}

Table~\ref{tab:dataset_summary2} summarizes the sample size, class
distribution, feature composition, and record-integrity audit of the three
evaluated public stroke benchmarks. DF-1 is the largest and most imbalanced
dataset, whereas DF-2 and DF-3 are smaller benchmarks with different
missing-data characteristics. Stroke-positive observations constitute
approximately \(1.8\%\), \(5.0\%\), and \(4.9\%\) of DF-1, DF-2, and DF-3,
respectively.

All three datasets contain ten candidate modeling predictors: age, gender,
hypertension, heart disease, marital status, work type, residence type,
average glucose level, body mass index (BMI), and smoking status. The binary
target variable is \texttt{stroke}, where \(1\) denotes stroke-positive and
\(0\) denotes non-stroke status. DF-1 and DF-3 contain an additional row
identifier, whereas DF-2 does not. Accordingly, DF-1 and DF-3 contain 12 raw
columns, while DF-2 contains 11; after excluding the target and identifier,
when present, all datasets contain ten modeling predictors. Detailed variable
definitions and data types are provided in Supplementary Table~S1.

 {
Before model development, record integrity was assessed by comparing complete
predictor--target profiles after excluding identifier columns. Predictor-only
profiles were also examined to identify records with identical predictors but
conflicting stroke labels. Table~\ref{tab:dataset_summary2} reports the
resulting duplicate groups, redundant records, conflicting-label groups, and
final retained sample sizes.
}

\begin{table*}[!ht]
\centering
\caption{
Dataset characteristics and record-integrity audit of the evaluated stroke
benchmarks.
}
\label{tab:dataset_summary2}
\small
\setlength{\tabcolsep}{4.2pt}
\renewcommand{\arraystretch}{1.10}

\begin{tabular}{@{}lcccccccc@{}}
\toprule
\textbf{Dataset} &
\(\boldsymbol{N}_{\mathrm{orig}}\) &
\textbf{S/NS} &
\(\boldsymbol{C}_{\mathrm{raw}}\) &
\(\boldsymbol{P}_{\mathrm{model}}\) &
\(\boldsymbol{G}_{\mathrm{dup}}\) &
\(\boldsymbol{N}_{\mathrm{red}}\) &
\(\boldsymbol{G}_{\mathrm{conf}}\) &
\(\boldsymbol{N}_{\mathrm{ret}}\) \\
\midrule

DF-1~\cite{dataset_df1} &
43,400 &
783/42,617 &
12 &
10 &
0 &
0 &
0 &
43,400 \\

DF-2~\cite{dataset_df2} &
4,981 &
248/4,733 &
11 &
10 &
0 &
0 &
0 &
4,981 \\

DF-3~\cite{dataset_df3} &
5,110 &
249/4,861 &
12 &
10 &
0 &
0 &
0 &
5,110 \\

\bottomrule
\end{tabular}

\vspace{2pt}
\begin{minipage}{0.98\textwidth}
\footnotesize
\textit{Note:}
\(N_{\mathrm{orig}}\), original records;
S/NS, stroke/non-stroke cases;
\(C_{\mathrm{raw}}\), raw columns;
\(P_{\mathrm{model}}\), modeling predictors;
\(G_{\mathrm{dup}}\), exact-duplicate groups;
\(N_{\mathrm{red}}\), redundant records;
\(G_{\mathrm{conf}}\), conflicting-label groups; and
\(N_{\mathrm{ret}}\), retained records.
\end{minipage}
\end{table*}

 {
The audit identified no exact-duplicate groups, redundant records, or
predictor-identical profiles with conflicting labels in any dataset.
Consequently, no observations were removed, and the original sample sizes were
retained for subsequent analysis. Identifier fields were used exclusively for
the integrity assessment and were removed before imputation, encoding, feature
selection, resampling, graph construction, and model training.
}

 {
The available identifiers in DF-1 and DF-3 represent unique rows rather than
verified longitudinal patient identities, and DF-2 contains no identifier
field. Therefore, the audit excludes leakage from exact duplicate records but
cannot determine whether clinically distinct records represent multiple
encounters involving the same individual. The datasets are consequently
treated as public benchmark files rather than independent external clinical
validation cohorts.
}
% =========================================================
% =========================================================
% =========================================================
\subsection{Class Imbalance Characteristics}
\label{subsec:class_imbalance}

Figure~\ref{fig:class_distribution} illustrates the strong dominance of
non-stroke observations across the three datasets. The imbalance is most
severe in DF-1, where stroke-positive records account for only \(1.8\%\) of
the data.

\begin{figure}[!h]
\centering
\includegraphics[width=\textwidth]
{images/stroke_datasets_distribution_2.png}
\caption{Class distribution of non-stroke and stroke cases across
(a) DF-1, (b) DF-2, and (c) DF-3.}
\label{fig:class_distribution}
\end{figure}

This imbalance has two important implications. First, Accuracy alone may be
dominated by the majority class; therefore, the evaluation emphasizes
minority-sensitive measures including Recall, F1-score, PR-AUC, MCC, and
balanced accuracy in addition to conventional discrimination measures.
Second, sparse stroke-positive observations may produce weak minority
neighborhoods in patient-similarity graphs, motivating the proposed
imbalance-aware refinement before graph construction. All resampling
operations are restricted to training data, as detailed in
Section~\ref{subsec:data_preparation}.
%====================================================================
\section{Proposed GNN-DHRS Framework}
\label{sec:methodology}
%====================================================================

This section presents the proposed \textbf{Graph Neural Network with Dual Hybrid Resampling Strategy (GNN-DHRS)} framework for explainable imbalanced brain stroke prediction. The framework is specifically designed to address three major challenges commonly encountered in clinical stroke-risk prediction: (i) severe class imbalance between stroke-positive and non-stroke patients, (ii) the inability of conventional machine-learning models to capture latent patient-to-patient relationships, and (iii) the need for transparent and clinically interpretable prediction mechanisms.

Unlike traditional tabular-learning approaches that treat each patient record independently, the proposed framework represents patients as nodes within a patient-similarity graph and explicitly models clinically meaningful relationships through graph-based relational learning. To improve minority-class representation prior to graph construction, a novel Dual Hybrid Resampling Strategy (DHRS) is introduced by sequentially combining Borderline-SMOTE and SMOTEENN. The resulting balanced dataset is subsequently transformed into a weighted patient-similarity graph and processed through a Multi-Head Graph Attention Network to perform stroke-risk prediction. Finally, feature-level and relation-level explainability modules are incorporated to provide transparent clinical interpretation of prediction outcomes.

The overall framework follows a leakage-safe learning protocol in which preprocessing, feature selection, resampling, graph construction, model training, validation, testing, and explainability analysis are performed independently within each cross-validation fold. This design prevents information leakage and ensures unbiased performance estimation.


\begin{figure*}[!h]
\centering
\includegraphics[width=\textwidth]{images/gnn_dhrs_framework_.png}
\caption{ {
Leakage-safe training and strictly inductive evaluation of GNN-DHRS. All
data-dependent operations are fitted using training data only, and each
untouched test patient is independently connected exclusively to original
training nodes under a frozen model, without test--test edges or test-set
information.
}}
\label{fig:gnn_dhrs_framework2}
\end{figure*}


Figure~\ref{fig:gnn_dhrs_framework2} presents the overall workflow of the proposed GNN-DHRS framework, whereas Algorithm~\ref{alg:gnn_dhrs} summarizes the complete training and inference procedure.


\begin{algorithm}[!ht]
\caption{Proposed GNN-DHRS Framework for Explainable Imbalanced Brain Stroke Prediction}
\label{alg:gnn_dhrs}
\small
\begin{algorithmic}[1]

\Require Clinical dataset \(D\),
\textcolor{red}{outer folds \(K_o=10\), inner folds \(K_i=5\),}
neighborhood size \(k\), epochs \(T\)

\Ensure Predictions, evaluation metrics, and explainability outputs

\State Define the leakage-safe preprocessing and feature-selection pipeline
\State Generate stratified
\textcolor{red}{\(K_o\)-fold outer cross-validation partitions}

\For{each \textcolor{red}{outer} fold \(f=1,\ldots,K_o\)}

    \State Split \(D\) into outer-training
    \(D_{\mathrm{tr}}^{(f)}\) and untouched outer-test
    \(D_{\mathrm{te}}^{(f)}\)

    \State \textcolor{red}{
    Generate \(K_i\) inner folds within \(D_{\mathrm{tr}}^{(f)}\)
    }

    \For{\textcolor{red}{each candidate configuration \(\theta\)}}

        \For{\textcolor{red}{each inner fold \(g=1,\ldots,K_i\)}}

            \State \textcolor{red}{
            Split outer-training data into inner-training and
            inner-validation partitions
            }

            \State \textcolor{red}{
            Fit preprocessing and feature selection using
            inner-training data only
            }

            \State \textcolor{red}{
            Apply DHRS only to the inner-training partition
            }

            \State \textcolor{red}{
            Train GNN-DHRS and evaluate \(\theta\) on the untouched
            inner-validation partition
            }

        \EndFor

        \State \textcolor{red}{
        Aggregate inner-validation performance
        }

    \EndFor

    \State \textcolor{red}{
    Select \(\theta^{*}\) exclusively from inner-validation results
    }

    \State \textcolor{red}{
    Refit preprocessing and feature selection on the complete
    outer-training partition
    }

    \State Apply Borderline-SMOTE minority enhancement
    \textcolor{red}{to outer-training data only}

    \State Apply SMOTEENN noise refinement
    \textcolor{red}{to outer-training data only}

    \State Obtain balanced training dataset \(D_b\)

    \State Construct weighted training graph
    \(G_{\mathrm{tr}}\) without test nodes

    \State Initialize Multi-Head GAT
    \textcolor{red}{using \(\theta^{*}\)}

    \For{epoch \(=1,\ldots,T\)}

        \State Compute attention coefficients and aggregate neighbors
        \State Generate graph-aware representations and probabilities
        \State Compute WBCE loss and update parameters using AdamW
        \State Apply learning-rate scheduling and early stopping

    \EndFor

    \State Freeze the trained model and training-node representations
    \textcolor{red}{with the fitted pipeline and selected threshold}

    \For{each patient \(q\in D_{\mathrm{te}}^{(f)}\)}

        \State Transform \(q\) using the training-fitted pipeline
        \State Connect \(q\) only to original training nodes
        \State Predict without test--test edges or graph modification

    \EndFor

    \State Compute predictive and imbalance-sensitive metrics
    \State Generate SHAP and attention-based explanations
    \State Perform edge-ablation validation

\EndFor

\State Aggregate fold-wise performance statistics
\State Return predictions, metrics, and explainability outputs

\end{algorithmic}
\end{algorithm}
%--------------------------------------------------------------------
\subsection{Leakage-Safe Data Preparation}
\label{subsec:data_preparation}

 {
Identifier columns, where present, were used only for the record-integrity
audit and were excluded before preprocessing, feature selection, resampling,
graph construction, and model training.
}

All data-dependent preprocessing operations were fitted using training data
only. Missing numerical values were imputed using training-fold means,
missing categorical values were assigned an \texttt{Unknown} category,
categorical variables were one-hot encoded, and numerical predictors were
standardized as

\begin{equation}
x_i^{\prime}
=
\frac{x_i-\mu}{\sigma},
\label{eq:zscore_normalization}
\end{equation}

where \(\mu\) and \(\sigma\) denote the corresponding training-fold mean and
standard deviation.


Performance estimation used stratified nested cross-validation with
\(K_o=10\) outer folds and \(K_i=5\) inner folds. Within each outer-training
partition, preprocessing, feature ranking, hyperparameter selection, learning-
rate scheduling, early stopping, checkpoint selection, and decision-threshold
selection were performed using inner-training and inner-validation data only.
\textcolor{red}{
During model selection, every data-dependent operation was fitted within the
corresponding inner-training partition and evaluated using only its untouched
inner-validation partition. This included imputation, categorical encoding,
standardization, correlation filtering, feature ranking, feature-count
selection, DHRS resampling, graph construction, hyperparameter selection,
learning-rate scheduling, early stopping, checkpoint selection, and
decision-threshold selection. After the candidate configuration had been fixed
exclusively from inner-validation results, the preprocessing and
feature-selection pipeline was refitted on the complete outer-training
partition, DHRS was reapplied only to that outer-training partition, and the
resulting frozen pipeline was evaluated once on the untouched outer-test fold.
No outer-test information contributed to any model-development or
model-selection decision.
}


After standardization, highly correlated predictors were filtered using a
fixed Pearson-correlation threshold of
\(\lvert r_{jl}\rvert>0.90\), followed by ANOVA \(F\)-score ranking:

\begin{equation}
F_j
=
\frac{
\sum_{c=1}^{C}
n_c
\left(
\bar{x}_{c,j}-\bar{x}_{j}
\right)^2
}{
\sum_{c=1}^{C}
\sum_{i=1}^{n_c}
\left(
x_{i,c,j}-\bar{x}_{c,j}
\right)^2
},
\label{eq:anova_fscore}
\end{equation}

where \(C\), \(n_c\), \(\bar{x}_{c,j}\), and \(\bar{x}_{j}\) denote the
number of classes, class size, class-specific feature mean, and overall
feature mean, respectively.

 {
The retained feature count was treated as a model-selection hyperparameter.
Candidate configurations consisted of the top 5, top 10, top 15, and all
transformed features, with selection based on mean inner-validation F1-score.
The selected configuration retained the ten highest-ranked transformed
features. The corresponding predictive and graph-topology analyses are
reported in Table~\ref{tab:feature_selection_joint_analysis}.
}

 {
Final performance was computed exclusively from predictions on the untouched
outer-test folds. Neither validation observations nor synthetic samples were
included in the reported test metrics, and no outer-test information was used
for preprocessing, feature selection, model tuning, early stopping, checkpoint
selection, or decision-threshold optimization.
}
%---------------------------------------------
\subsection{Dual Hybrid Resampling Strategy (DHRS)}
\label{subsec:dhrs}

Severe class imbalance can bias learning toward the majority non-stroke class
and reduce minority-case detection
\cite{he2009learning,haixiang2017learning}. To address this issue, the proposed
Dual Hybrid Resampling Strategy (DHRS) sequentially combines Borderline-SMOTE
and SMOTEENN before patient-graph construction. Borderline-SMOTE enhances
difficult minority regions near the decision boundary, whereas SMOTEENN
subsequently performs additional resampling and Edited Nearest Neighbour
cleaning.

 {
The methodological novelty of DHRS lies in the ordered, training-only
integration of these two established resampling methods before
patient-similarity graph construction rather than in modifying either method
individually. Consequently, the resampling sequence changes both the
class composition and the observations from which graph neighborhoods are
subsequently formed.
}

Given a minority observation \(x_i\) and one of its minority neighbors
\(x_{zi}\), Borderline-SMOTE generates

\begin{equation}
x_{\mathrm{new}}
=
x_i+
\lambda(x_{zi}-x_i),
\qquad
\lambda\in[0,1],
\label{eq:borderline_smote}
\end{equation}

where \(\lambda\) determines the interpolation position
\cite{han2005borderline}. SMOTEENN then applies neighborhood cleaning
\cite{batista2004study}. For observation \(x_i\), the ENN removal criterion is

\begin{equation}
y_i
\neq
\operatorname{mode}
\left(
y_j \;|\; x_j\in N_k(x_i)
\right),
\label{eq:enn_rule}
\end{equation}

where \(N_k(x_i)\) denotes its \(k\)-nearest neighbors and \(y_i\) its class
label \cite{wilson1972asymptotic}.

 {
DHRS was applied exclusively to training partitions. Validation and
outer-test observations were never resampled or used for synthetic-sample
generation, neighborhood cleaning, or resampling-parameter selection,
consistent with the leakage-safe protocol described in
Section~\ref{subsec:data_preparation}.
}

 {
Because the datasets contain both continuous and categorical predictors, the
primary implementation applied one-hot encoding using training-derived
categories before DHRS. After interpolation, each generated one-hot block was
projected to a valid categorical representation by assigning \(1\) to its
largest component and \(0\) to the remaining components, while continuous
variables retained their interpolated values. A controlled alternative using SMOTENC followed by ENN was also evaluated;
the corresponding sensitivity analysis is reported in
Table~\ref{tab:mixed_resampling_sensitivity}.}

The complete resampling sequence is

\begin{equation}
D_{\mathrm{DHRS}}
=
\mathrm{SMOTEENN}
\Big(
\mathrm{Borderline\mbox{-}SMOTE}
(D_{\mathrm{train}})
\Big),
\label{eq:dhrs_framework}
\end{equation}

where \(D_{\mathrm{train}}\) denotes the corresponding training partition and
\(D_{\mathrm{DHRS}}\) is the refined dataset used for graph construction and
model learning.

\begin{figure*}[!t]
\centering
\includegraphics[width=0.95\textwidth]{images/dhrs_workflow.png}
\caption{Dual Hybrid Resampling Strategy (DHRS). Borderline-SMOTE enhances
difficult minority regions, followed by SMOTEENN resampling and
neighborhood-based cleaning before patient-graph construction.}
\label{fig:dhrs_workflow}
\end{figure*}

 {
To determine whether the proposed ordering was important, DHRS was evaluated
against no resampling, conventional SMOTE, Borderline-SMOTE, SMOTEENN,
Borderline-SMOTE followed by ENN, and the reversed
SMOTEENN--Borderline-SMOTE sequence. All configurations used the same data
partitions, preprocessing, graph construction, model architecture,
optimization protocol, random seeds, and validation-selected decision
thresholds. The predictive comparison is reported in
Table~\ref{tab:resampling_comparison_gnn2}, while detailed sample-generation,
removal, and final class counts are provided in Supplementary Table~S4.}

 {
For resampling configuration \(r\) in outer-training fold \(f\), the final
imbalance ratio was defined as
}

{ 
\begin{equation}
\operatorname{IR}_{r}^{(f)}
=
\frac{N_{\mathrm{maj},r}^{(f)}}
     {N_{\mathrm{min},r}^{(f)}},
\label{eq:final_resampling_ir}
\end{equation}
}

 {
where \(N_{\mathrm{maj},r}^{(f)}\) and
\(N_{\mathrm{min},r}^{(f)}\) denote the final majority- and minority-class
counts. Values closer to \(1\) indicate a more balanced post-resampling
training distribution.
}

 {
Because SMOTEENN cleaning can leave residual class imbalance, weighted binary
cross-entropy (WBCE) was applied after DHRS. For outer-training fold \(f\),
with \(N^{(f)}=N_0^{(f)}+N_1^{(f)}\), the class weights were computed from
the final post-DHRS training distribution as
}

{ 
\begin{equation}
w_0^{(f)}
=
\frac{N^{(f)}}{2N_0^{(f)}},
\qquad
w_1^{(f)}
=
\frac{N^{(f)}}{2N_1^{(f)}},
\label{eq:dhrs_class_weights}
\end{equation}
}

 {
where \(N_0^{(f)}\) and \(N_1^{(f)}\) denote the final non-stroke and stroke
training counts, respectively. Validation and test labels were not used to
calculate these weights. A controlled BCE-versus-WBCE ablation was performed using otherwise identical
post-DHRS data, graphs, and model settings; the results are reported in
Table~\ref{tab:dhrs_weight_ablation}.}
%=============================================

\subsection{Dynamic Patient-Similarity Graph Construction}
\label{subsec:graph_construction}

After DHRS, the refined training observations are transformed into a weighted
patient-similarity graph for graph-attention learning. Each patient is
represented as a node using the selected feature vector
\(\mathbf{x}_i\in\mathbb{R}^{d}\), and relationships between clinically
similar patients are defined using cosine similarity
\cite{kipf2017semi,hamilton2017inductive,velickovic2018graph}.
The overall procedure is illustrated in
Figure~\ref{fig:graph_construction_workflow}.

For patients \(i\) and \(j\), cosine similarity is defined as

\begin{equation}
s_{ij}
=
\frac{
\mathbf{x}_i^{T}\mathbf{x}_j
}{
\|\mathbf{x}_i\|_2
\|\mathbf{x}_j\|_2
},
\label{eq:cosine_similarity}
\end{equation}

where the feature vectors are obtained after the training-fitted preprocessing
and feature-selection procedure described in
Section~\ref{subsec:data_preparation}.

To avoid the computational cost and weak relations associated with a fully
connected graph, a sparse top-\(k\) neighborhood is constructed for each
patient:

\begin{equation}
\mathcal{N}_k(i)
=
\operatorname{TopK}_{j\ne i}
\left(
s_{ij}
\right),
\label{eq:topk_neighbors}
\end{equation}

where only admissible positive similarities are considered. Based on the
neighborhood-sensitivity analysis, \(k=10\) was selected for the primary
configuration.

 {
The resulting graph is directed. An edge \(i\rightarrow j\) is retained when
\(j\in\mathcal{N}_k(i)\); the reverse edge is included only if
\(i\in\mathcal{N}_k(j)\). The graph is therefore neither automatically
symmetrized nor restricted to mutual-\(k\)-NN relations. Self-neighbors are
excluded during graph construction, while self-loops are added within the GAT
layers during message passing.
}
\textcolor{red}{
The same directed support is used consistently for model training,
graph-attention message passing, and the topology analyses associated with the
primary GNN-DHRS configuration. No separate symmetrized version of the primary
cosine \(k\)-NN graph is used for topology evaluation. Reciprocal relations
are therefore represented as two directed edges when both
\(i\rightarrow j\) and \(j\rightarrow i\) independently satisfy the
neighborhood criterion. Self-loops introduced internally by the GAT layers
are excluded from the descriptive \(k\)-NN topology statistics.
}
 {
Only positive cosine similarities were eligible for graph connections.
Candidates satisfying \(s_{ij}\leq0\) were discarded before neighbor
selection, and no additional positive-similarity threshold was imposed in the
primary graph. If fewer than \(k\) positive neighbors were available, all
admissible neighbors were retained.
}

The weighted adjacency matrix is defined as

\begin{equation}
A_{ij}
=
\begin{cases}
s_{ij},
&
j\in\mathcal{N}_k(i)
\ \text{and}\ s_{ij}>0,\\
0,
&
\text{otherwise},
\end{cases}
\label{eq:weighted_adjacency}
\end{equation}

where \(A_{ij}\) represents the retained similarity strength between patients
\(i\) and \(j\).

 {
The cosine edge weights were not subjected to additional row-wise,
degree-based, or symmetric normalization. Instead, the raw positive weights
were incorporated directly into the similarity-aware attention mechanism
defined in Section~\ref{subsec:multihead_attention}, where local softmax
normalization determines the final neighbor contributions.
}

 {
During training, both original and DHRS-generated observations participated
in graph construction under the same feature representation, positive-cosine
top-\(k\) rule, and message-passing mechanism. Synthetic provenance was not
provided to the classifier as an input feature and was retained only for
structural auditing.
}

 {
To quantify the structural participation of synthetic observations, training
edges were categorized as real--real, real--synthetic, or
synthetic--synthetic. This audit did not modify graph construction or model
training. The dataset-specific graph-composition statistics are reported in
Table~\ref{tab:synthetic_edge_audit}.
}

\begin{figure*}[!h]
\centering
\includegraphics[width=0.95\textwidth]
{images/patient_similarity_graph_construction.png}
\caption{Dynamic patient-similarity graph construction. Refined patient
profiles are represented as nodes, positive cosine similarities are computed,
and directed top-\(k\) relations are retained to construct the weighted graph
used by the graph-attention model.}
\label{fig:graph_construction_workflow}
\end{figure*}

The resulting weighted graph
\(G=(V,E,A)\) provides the relational input to the Multi-Head Graph Attention
Network, with graph density controlled through the sparse top-\(k\)
neighborhood structure.

 {
Outer-test patients are not incorporated into this training graph. Their
strictly inductive attachment to original outer-training patients is defined
separately in Section~\ref{subsec:inductive_graph_construction}, thereby
avoiding repetition of the test-time leakage-control protocol here.
}

 {
To assess dependence on the selected graph definition, the primary directed
cosine \(k\)-NN graph was compared with Euclidean \(k\)-NN, Gower \(k\)-NN,
mutual cosine \(k\)-NN, and cosine similarity-threshold alternatives. All
variants used the same data partitions, retained features, model architecture,
optimization settings, random seeds, and validation protocol. For the
\(k\)-NN alternatives, \(k=10\) was held fixed. The threshold-based graph used
a training-selected similarity threshold chosen to approximate the mean
degree of the primary graph.
}

 {
All alternative graphs followed the same strictly inductive inference
boundary: an unseen patient could connect only to original outer-training
patients, with no test--test edges or target-derived graph statistics.
Predictive and topological comparisons are reported in
Section~\ref{subsec:graph_definition_comparison}, including the topology
statistics presented in Table~\ref{tab:graph_topology_dhrs}.}
%===========================
%---------------------------------------------------------------
%=========================================================
\subsection{Leakage-Safe Inductive Graph Construction}
\label{subsec:inductive_graph_construction}

 {
Strictly inductive inference was used to prevent transductive information
leakage. The nested training and model-selection protocol follows
Section~\ref{subsec:data_preparation}. After the final configuration was
selected, preprocessing and feature-selection operations were refitted on the
complete outer-training partition, the model was retrained, and all fitted
transformations, training-node representations, model parameters, graph
settings, and the validation-selected decision threshold were frozen before
outer-test evaluation.
}

 {
For outer fold \(f\), let \(\mathcal{T}^{(f)}\) denote the transformation
fitted exclusively on the outer-training data and
\(V_{\mathrm{tr,real}}^{(f)}\) the set of original real outer-training nodes.
Each unseen test patient \(q\) was transformed independently as
\(\widetilde{\mathbf{x}}_{q}^{(f)}
=\mathcal{T}^{(f)}(\mathbf{x}_{q})\) and connected only to original
outer-training patients. DHRS-generated synthetic nodes participated in model
training but were excluded from the reference-node set used during test-time
graph construction.
}

 {
The inductive neighborhood of patient \(q\) was defined as
}

{ 
\begin{equation}
\mathcal{N}_{\mathrm{tr}}^{(f)}(q)
=
\underset{
j\in V_{\mathrm{tr,real}}^{(f)}
}{
\operatorname{TopK}
}
\;
s\left(
\widetilde{\mathbf{x}}_{q}^{(f)},
\widetilde{\mathbf{x}}_{j}^{(f)}
\right),
\label{eq:inductive_test_neighborhood}
\end{equation}
}

 {
where \(s(\cdot,\cdot)\) denotes cosine similarity and only positive
similarities were considered, consistent with the training-graph definition.
Up to \(k=10\) original training neighbors were retained. Directed
training-to-test edges were then created so that information could propagate
from the frozen training representation to the unseen patient without
altering the training graph.
}

 {
No test--test similarities or edges were constructed, no reverse
test-to-training edges were introduced, and no outer-test labels or aggregate
test-set statistics were used. Consequently, processing one unseen patient
could not modify the training graph, model parameters, training-node
representations, or the prediction of another test patient.
Algorithm~\ref{alg:inductive_inference} summarizes this inference boundary.
}

\begingroup
 
\begin{algorithm}[!t]
\caption{Leakage-Safe Inductive Inference for an Unseen Patient}
\label{alg:inductive_inference}
\small
\begin{algorithmic}[1]

\Require
Unseen patient \(\mathbf{x}_{q}\);
training-fitted transformation \(\mathcal{T}^{(f)}\);
original training nodes \(V_{\mathrm{tr,real}}^{(f)}\);
frozen training representations and GNN-DHRS parameters;
neighborhood size \(k=10\)

\Ensure
Predicted stroke probability \(\widehat{p}_{q}^{(f)}\)

\State Transform
\(\mathbf{x}_{q}\) using
\(\mathcal{T}^{(f)}\)

\State Compute positive cosine similarities only between \(q\) and
\(V_{\mathrm{tr,real}}^{(f)}\)

\State Select up to \(k\) original training neighbors using
Eq.~(\ref{eq:inductive_test_neighborhood})

\State Create incoming training-to-test edges for the selected neighbors

\State Do not create test--test or test-to-training edges

\State Keep the training graph, training representations, and model
parameters frozen

\State Aggregate messages from the selected training neighbors to \(q\)

\State Compute \(\widehat{p}_{q}^{(f)}\) using the frozen prediction layer

\State \Return \(\widehat{p}_{q}^{(f)}\)

\end{algorithmic}
\end{algorithm}
\endgroup
%=========================================================
%===========================
%=========================================================
\subsection{Multi-Head Graph Attention Learning}
\label{subsec:multihead_attention}

After constructing the weighted patient-similarity graph in
Section~\ref{subsec:graph_construction}, GNN-DHRS applies a two-layer
multi-head Graph Attention Network (GAT) to learn relational patient
representations for stroke-status classification
\cite{velickovic2018graph}. The input representation of patient \(i\) is its
selected clinical feature vector,
\(\mathbf{h}_i^{(0)}=\mathbf{x}_i\).

 {
Both graph-attention layers use \(H=8\) parallel heads with a total hidden
dimension of \(128\), corresponding to \(16\) dimensions per head. For head
\(m\) in layer \(l\), the projected node representation is
\(\mathbf{z}_{i,m}^{(l)}
=\mathbf{W}_{m}^{(l)}\mathbf{h}_{i}^{(l)}\).
For target node \(i\) and neighboring node \(j\), the feature-based
compatibility score is
}

{ 
\begin{equation}
u_{ij,m}^{(l)}
=
\operatorname{LeakyReLU}
\left(
\mathbf{a}_{m}^{(l)T}
\left[
\mathbf{z}_{i,m}^{(l)}
\parallel
\mathbf{z}_{j,m}^{(l)}
\right]
\right),
\label{eq:attention_score}
\end{equation}
}

 {
where \(\mathbf{a}_{m}^{(l)}\) is a trainable attention vector,
\(\parallel\) denotes concatenation, and the LeakyReLU negative slope is
\(0.2\). Unlike a conventional unweighted GAT, GNN-DHRS incorporates the
positive cosine-similarity edge weight \(A_{ij}\) directly into attention
normalization. Let
\(\widetilde{\mathcal{N}}(i)=\mathcal{N}(i)\cup\{i\}\) denote the incoming
neighborhood augmented with a self-loop, for which \(A_{ii}=1\). The
similarity-aware attention coefficient is
}

{ 
\begin{equation}
\alpha_{ij,m}^{(l)}
=
\frac{
A_{ij}\exp\!\left(u_{ij,m}^{(l)}\right)
}{
\displaystyle
\sum_{r\in\widetilde{\mathcal{N}}(i)}
A_{ir}\exp\!\left(u_{ir,m}^{(l)}\right)
},
\qquad
j\in\widetilde{\mathcal{N}}(i).
\label{eq:attention_coefficient}
\end{equation}
}

 {
Thus, attention depends jointly on learned feature compatibility and the
training-derived cosine similarity. Because only positive graph similarities
are retained, the coefficients are well defined and locally normalized over
the incoming neighbors and self-loop.
}

\begin{figure}[!h]
\centering
\includegraphics[width=\linewidth]
{images/graph_attention_mechanism.png}
\caption{ {
Similarity-aware graph attention in GNN-DHRS. Learned feature compatibility
and cosine edge weights jointly determine the normalized neighborhood
attention coefficients.
}}
\label{fig:graph_attention_mechanism}
\end{figure}

For attention head \(m\), neighborhood information is aggregated as

{ 
\begin{equation}
\mathbf{h}_{i,m}^{(l+1)}
=
\varphi
\left(
\sum_{j\in\widetilde{\mathcal{N}}(i)}
\alpha_{ij,m}^{(l)}
\mathbf{W}_{m}^{(l)}
\mathbf{h}_{j}^{(l)}
\right),
\label{eq:gat_aggregation}
\end{equation}
}

where \(\varphi(\cdot)\) denotes the nonlinear activation. The outputs of the
\(H\) attention heads are concatenated:

\begin{equation}
\mathbf{h}_i^{(l+1)}
=
\mathop{\parallel}_{m=1}^{H}
\mathbf{h}_{i,m}^{(l+1)}.
\label{eq:multihead_attention}
\end{equation}

 {
Each attention layer follows the sequence
GAT--BatchNorm--ReLU--Dropout. After the second layer, the final relational
representation is mapped to a single logit and converted to a stroke-status
probability:
}

{ 
\begin{equation}
\widehat{p}_i
=
\operatorname{sigmoid}
\left(
\mathbf{w}_{o}^{T}\mathbf{h}_{i}^{(2)}+b_o
\right).
\label{eq:gat_output_probability}
\end{equation}
}

\begin{figure*}[!h]
\centering
\includegraphics[width=0.95\textwidth]
{images/multihead_attention_learning_2.png}
\caption{ {
Two-layer multi-head graph-attention architecture of GNN-DHRS. Similarity-aware
attention is followed by batch normalization, ReLU activation, dropout, and a
sigmoid classification layer.
}}
\label{fig:multihead_attention_learning}
\end{figure*}

\begin{table}[!ht]
\centering
\caption{Core GNN-DHRS model and training configuration.}
\label{tab:gat_configuration}

\small
\setlength{\tabcolsep}{5pt}
\renewcommand{\arraystretch}{1.05}

\begin{tabular}{@{}lc@{}}
\toprule
\textbf{Parameter} & \textbf{Selected value} \\
\midrule

GAT layers & 2 \\
Hidden dimension & 128 \\
Attention heads & 8 \\
Dimension per head & 16 \\
Dropout & 0.30 \\
LeakyReLU slope & 0.20 \\
Graph neighbors \(k\) & 10 \\

\textcolor{red}{DHRS oversampling ratio \(\rho\)}
& \textcolor{red}{1.00} \\

Loss & Weighted BCE \\
Optimizer & AdamW \\

\textcolor{red}{AdamW coefficients}
& \textcolor{red}{\(\beta_1=0.9,\ \beta_2=0.999,\ \epsilon=10^{-8}\)} \\

Initial learning rate & \(1\times10^{-3}\) \\
Weight decay & \(5\times10^{-4}\) \\
Maximum epochs & 200 \\

\textcolor{red}{Training mode}
& \textcolor{red}{Full-graph, full-batch} \\

\textcolor{red}{LR scheduler}
& \textcolor{red}{ReduceLROnPlateau} \\

\textcolor{red}{LR reduction factor}
& \textcolor{red}{0.50} \\

\textcolor{red}{Scheduler patience}
& \textcolor{red}{10 epochs} \\

\textcolor{red}{Minimum learning rate}
& \textcolor{red}{\(1\times10^{-6}\)} \\

\textcolor{red}{Early-stopping monitor}
& \textcolor{red}{Inner-validation F1-score} \\

\textcolor{red}{Early-stopping patience}
& \textcolor{red}{20 epochs} \\

\textcolor{red}{Checkpoint selection}
& \textcolor{red}{Highest inner-validation F1-score} \\

\textcolor{red}{Random seeds}
& \textcolor{red}{\(\{11,22,33,44,55\}\)} \\

\bottomrule
\end{tabular}
\end{table}

\textcolor{red}{
Training used full-graph, full-batch optimization with AdamW. Weighted binary
cross-entropy employed fold-specific post-DHRS class weights defined in
Eq.~(\ref{eq:dhrs_class_weights}). Trainable weights were initialized using
Glorot uniform initialization and biases were initialized to zero
\cite{glorot2010understanding}.
}
\textcolor{red}{
The optimization, learning-rate scheduling, early-stopping, checkpoint-selection,
random-seed, graph-neighborhood, and DHRS settings are summarized in
Table~\ref{tab:gat_configuration}. All architecture, optimization, resampling,
and graph hyperparameters were selected exclusively from inner-validation
data.
}

 {
Training continued for at most 200 epochs. ReduceLROnPlateau monitored
inner-validation F1-score, reduced the learning rate by a factor of \(0.50\)
after 10 epochs without improvement, and used a minimum learning rate of
\(10^{-6}\). Early stopping used the same validation criterion with a
patience of 20 epochs, and the checkpoint with the highest inner-validation
F1-score was restored.
}

 {
Architecture, optimization, resampling, and graph hyperparameters were
selected exclusively from inner-validation data. Hyperparameter sensitivity
is reported in Table~\ref{tab:hyperparameter_sensitivity}. Experiments were
repeated using seeds \(\{11,22,33,44,55\}\), applied consistently to data
partitioning, DHRS generation, initialization, dropout, and model training;
the corresponding stochastic-stability analysis is reported in
Table~\ref{tab:random_seed_stability}.
}

Overall, the graph-attention module combines cosine-informed edge strength
with learned multi-head attention to transform the sparse patient graph into
relational representations used for stroke-status classification.
%=========================================================

\subsection{Prediction and Optimization}
\label{subsec:prediction}

The final graph-attention representation is converted to the stroke-status
probability \(\widehat{p}_i\) using the sigmoid output layer defined in
Eq.~(\ref{eq:gat_output_probability}).

 {
For each outer fold, the decision threshold was selected by maximizing the
F1-score of the pooled inner-validation predictions, as formally defined in
Section~\ref{subsec:threshold_calibration}. The selected threshold was fixed
before evaluating the corresponding untouched outer-test fold. No outer-test
labels, probability distributions, or test-derived performance measures
contributed to threshold selection. The complete validation-based sensitivity
analysis is reported in Table~\ref{tab:threshold_sensitivity}.
}

To reduce residual class imbalance after DHRS, the network was trained using
weighted binary cross-entropy (WBCE):

\begin{equation}
\mathcal{L}_{\mathrm{WBCE}}
=
-\frac{1}{N}
\sum_{i=1}^{N}
\left[
w_1y_i\log(\widehat{p}_i)
+
w_0(1-y_i)\log(1-\widehat{p}_i)
\right],
\label{eq:weighted_bce}
\end{equation}

where \(y_i\in\{0,1\}\) denotes the observed class label, and \(w_0\) and
\(w_1\) are the fold-specific post-DHRS class weights defined in
Eq.~(\ref{eq:dhrs_class_weights}).

 {
Model parameters were optimized using AdamW with the configuration summarized
in Table~\ref{tab:gat_configuration}. Learning-rate scheduling, checkpoint
selection, and early stopping were based exclusively on inner-validation
F1-score, while all outer-test folds remained isolated from model selection
and were accessed only for final performance evaluation.
}
%-------------------------------------------------------
%=========================================================
%=========================================================
\subsection{Evaluation Measures}
\label{subsec:evaluation_measures}

\textcolor{red}{
GNN-DHRS was evaluated using complementary measures of classification
performance, ranking discrimination, probability calibration, and selective
prediction. Because severe class imbalance can make aggregate measures such as
Accuracy insufficient for assessing minority-class performance, conventional
and imbalance-sensitive metrics were considered jointly
\cite{he2009learning,haixiang2017learning}.
Table~\ref{tab:evaluation_metrics} summarizes the complete set of evaluation
measures, together with their mathematical definitions and interpretation.
}

\textcolor{red}{
Classification performance was assessed using Accuracy, Precision, Recall,
F1-score, Specificity, negative predictive value (NPV), balanced accuracy,
Matthews correlation coefficient (MCC), and Cohen's Kappa. Collectively, these
measures characterize overall correctness, stroke-positive detection,
non-stroke recognition, class-balanced performance, and agreement beyond
chance.
}

\textcolor{red}{
Ranking discrimination was evaluated using ROC-AUC and PR-AUC. ROC-AUC
characterizes discrimination over the range of decision thresholds, whereas
PR-AUC provides a complementary assessment focused on the Precision--Recall
relationship and is particularly relevant when the positive class is rare.
}

\textcolor{red}{
Probability reliability was evaluated using the Brier score and Expected
Calibration Error (ECE). These measures assess probability-estimation error
and the agreement between predicted confidence and observed event frequency,
respectively.
}

\textcolor{red}{
Confidence-aware selective prediction was evaluated using coverage, selective
risk, and the area under the risk--coverage curve (AURC). These measures
characterize the trade-off between retaining automatic predictions and
reducing prediction error through the deferral of low-confidence cases.
}

\begin{table*}[!ht]
\centering
\caption{\textcolor{red}{
Summary of the performance and reliability measures used in the evaluation
of GNN-DHRS.
}}
\label{tab:evaluation_metrics}

\scriptsize
\setlength{\tabcolsep}{4pt}
\renewcommand{\arraystretch}{1.22}

\begin{tabularx}{\textwidth}{@{}
>{\raggedright\arraybackslash}p{2.2cm}
>{\raggedright\arraybackslash}p{2.0cm}
>{\centering\arraybackslash}p{5.2cm}
>{\raggedright\arraybackslash}X
@{}}

\toprule
\textbf{Category} &
\textbf{Measure} &
\textbf{Definition} &
\textbf{Interpretation} \\
\midrule

\textbf{Classification}
&
Accuracy
&
\(\displaystyle
\frac{TP+TN}{TP+TN+FP+FN}
\)
&
Overall correctness
\\

&
Precision
&
\(\displaystyle
\frac{TP}{TP+FP}
\)
&
Positive-prediction reliability
\\

&
Recall
&
\(\displaystyle
\frac{TP}{TP+FN}
\)
&
Positive-class sensitivity
\\

&
F1-score
&
\(\displaystyle
\frac{2PR}{P+R}
\)
&
Precision--Recall balance
\\

&
Specificity
&
\(\displaystyle
\frac{TN}{TN+FP}
\)
&
Negative-class recognition
\\

&
NPV
&
\(\displaystyle
\frac{TN}{TN+FN}
\)
&
Negative-prediction reliability
\\

&
Balanced Accuracy
&
\(\displaystyle
\frac{R+\mathrm{Spec}}{2}
\)
&
Class-balanced accuracy
\\

&
MCC
&
\(\displaystyle
\frac{TP\,TN-FP\,FN}{D_{\mathrm{MCC}}}
\)
&
Balanced correlation
\\

&
Cohen's Kappa
&
\(\displaystyle
\frac{p_o-p_e}{1-p_e}
\)
&
Chance-corrected agreement
\\

\midrule

\textbf{Ranking}
&
ROC-AUC
&
\(\displaystyle
\int_{0}^{1}
\mathrm{TPR}(u)\,d\mathrm{FPR}(u)
\)
&
Overall ranking ability
\\

&
PR-AUC
&
\(\displaystyle
\int_{0}^{1}
\mathrm{Precision}(r)\,dr
\)
&
Minority-class ranking
\\

\midrule

\textbf{Calibration}
&
Brier score
&
\(\displaystyle
\frac{1}{N}
\sum_{i=1}^{N}
(\widehat{p}_i-y_i)^2
\)
&
Probability error
\\

&
ECE
&
\(\displaystyle
\sum_{b=1}^{B}
\frac{|S_b|}{N}
|f_b-c_b|
\)
&
Calibration error
\\

\midrule

\textbf{Selective prediction}
&
Coverage
&
\(\displaystyle
\frac{N_{\mathrm{accepted}}}{N}
\)
&
Accepted-prediction rate
\\

&
Selective Risk
&
\(\displaystyle
\frac{N_{\mathrm{error,accepted}}}
{N_{\mathrm{accepted}}}
\)
&
Accepted-case error rate
\\

&
AURC
&
\(\displaystyle
\int_{0}^{1}
R_{\mathrm{sel}}(c)\,dc
\)
&
Overall selective risk
\\

\bottomrule
\end{tabularx}

\vspace{2pt}

\begin{minipage}{0.98\textwidth}
\footnotesize
\textit{Note:}
\(TP,TN,FP,FN\) are confusion-matrix counts; \(P,R,\mathrm{Spec}\) denote
Precision, Recall, and Specificity; \(D_{\mathrm{MCC}}\) is the MCC denominator.
\(p_o,p_e\) denote observed and expected agreement; \(N\) is sample size;
\(\widehat{p}_i,y_i\) are predicted probability and label; \(S_b\) denotes a
calibration bin; and \(R_{\mathrm{sel}}(c)\) is selective risk at coverage \(c\).
Threshold-based metrics use inner-validation-selected thresholds; probability
metrics use continuous outer-test probabilities.
\end{minipage}

\end{table*}

\textcolor{red}{
All evaluation measures were computed independently within each untouched
outer-test fold. Threshold-dependent classification measures used the
corresponding decision threshold selected exclusively from inner-validation
predictions, whereas ROC-AUC, PR-AUC, Brier score, and ECE were calculated
from continuous outer-test probabilities. Unless otherwise specified for
selective-prediction analyses, fold-level results were summarized as mean
\(\pm\) standard deviation across the ten outer folds.
}
%=========================================================
\subsection{Explainability and Clinical Interpretation}
\label{subsec:explainability}

GNN-DHRS incorporates complementary feature- and relation-level
explainability to examine the basis of its stroke-status predictions.
As illustrated in Figure~\ref{fig:graph_explainability_pipeline},
 {
Kernel SHAP is used for feature attribution, graph-attention coefficients
provide patient-relation importance indicators, and controlled edge-ablation
experiments evaluate whether highly weighted relations are functionally
relevant to the model output. Attention weights are interpreted as relational
importance indicators rather than causal explanations.
}

\begin{figure*}[!h]
\centering
\includegraphics[width=\textwidth]
{images/GNN_AS_explainnable_1.png}
\caption{
Graph-based explainability pipeline of GNN-DHRS.
 {
Kernel SHAP provides feature attribution, attention coefficients support
relational interpretation, and controlled edge ablation evaluates the
predictive relevance of selected patient relations.
}}
\label{fig:graph_explainability_pipeline}
\end{figure*}

Feature-level attribution was performed using
 {\texttt{shap.KernelExplainer} from SHAP version 0.44.0}
\cite{lundberg2017unified}.  {
The explained output was the continuous stroke probability produced by the
complete frozen GNN-DHRS inference pipeline. Within each outer fold, the SHAP
background consisted of 100 original outer-training patients selected by
class-stratified sampling (50 stroke and 50 non-stroke). Synthetic DHRS
observations, validation patients, and outer-test patients were excluded from
background construction.
}

\textcolor{red}{
Kernel SHAP was applied to the selected patient-level input features rather
than to the latent representations generated after graph-attention
aggregation. Accordingly, the attribution variables correspond to the
transformed clinical inputs entering the complete GNN-DHRS inference
pipeline, while the explained output is the final stroke probability produced
after graph-based neighborhood aggregation and classification. Because this
prediction depends jointly on the target patient's features and information
propagated from neighboring patients, the resulting SHAP values are
interpreted as model-level feature attributions within the complete
graph-based prediction function rather than as topology-independent effects
or a decomposition of the internal graph representation. Graph-attention
message passing introduces dependencies between the target representation and
its neighborhood; therefore, the additive SHAP decomposition does not imply
statistical independence, causality, or a unique separation between direct
feature and relational effects. The feature-level SHAP analysis is therefore
interpreted jointly with attention-based relational analysis and controlled
edge ablation, which provide complementary evidence regarding the role of
patient-to-patient relations in the fitted model.
}

 {
SHAP values were computed for 20 untouched outer-test patients per fold,
comprising 10 stroke and 10 non-stroke cases. This yielded 200 explained
patients per dataset and 600 patients overall. Sampling was determined by the
predefined experimental seed and was independent of prediction correctness,
confidence, or explanation outcome.
}

Global feature importance for feature \(k\) was calculated as

\begin{equation}
I_k
=
\frac{1}{M}
\sum_{i=1}^{M}
\left|
\phi_{ik}
\right|,
\label{eq:global_shap}
\end{equation}

where \(\phi_{ik}\) is the Kernel-SHAP attribution for patient \(i\) and
\(M=200\) is the number of explained outer-test patients per dataset.
 {
Fold-specific explanations used only the model, preprocessing pipeline,
selected features, and training-derived background associated with the fold
in which each patient was evaluated.
}

 {
Feature-ranking stability was assessed from the agreement of the five
highest-ranked features across outer folds using the Jaccard index. The
top-five set was used because each fold retained ten transformed model inputs.
The corresponding feature-ranking stability results are reported in
 Table~\ref{tab:explanation_stability}.
}



Relation-level interpretation was derived from the normalized graph-attention
coefficients defined in Eq.~(\ref{eq:attention_coefficient}).
 {
A larger attention coefficient indicates that a neighboring training patient
received greater relative weight during message aggregation for the target
patient. This interpretation does not imply causality or that attention alone
constitutes a complete explanation of the prediction. Representative
high-attention patient relations are provided in  \textcolor{red}{Table \ref{tab:top_attention_relations}}.
}

 {
To evaluate relational explanations quantitatively, attention-based ranking
was compared with GNNExplainer \cite{ying2019gnnexplainer}, PGExplainer
\cite{luo2020pgexplainer}, and a random-edge reference using the same frozen
GNN-DHRS model and the same inductively constructed test neighborhoods.
For patient \(q\), all methods used an identical explanation budget
\(b_q=\min(5,|E_q|)\), where \(E_q\) denotes the candidate incoming
training-to-test edges. PGExplainer was trained using outer-training patients
only, whereas GNNExplainer was optimized separately for each test patient
using the frozen predicted class without access to the true test label.
Attention ranking required no additional fitting.
}

Let \(G_q\) denote the complete inductive graph for test patient \(q\),
\(E_q^{*}\) the selected explanatory edges, and \(\widehat{c}_q\) the
original predicted class. Explanation necessity and sufficiency were measured
as

\begin{equation}
\operatorname{Fid}^{+}_{q}
=
\widehat{p}_{q,\widehat{c}_q}(G_q)
-
\widehat{p}_{q,\widehat{c}_q}
\left(G_q\setminus E_q^{*}\right),
\label{eq:explanation_fidelity_positive}
\end{equation}

and

\begin{equation}
\operatorname{Fid}^{-}_{q}
=
\left|
\widehat{p}_{q,\widehat{c}_q}(G_q)
-
\widehat{p}_{q,\widehat{c}_q}
\left(G_q[E_q^{*}]\right)
\right|,
\label{eq:explanation_fidelity_negative}
\end{equation}

respectively. Larger \(\operatorname{Fid}^{+}\) indicates greater
necessity of the selected relations, whereas smaller
\(\operatorname{Fid}^{-}\) indicates greater sufficiency.

 {
Relational stability was additionally evaluated under ten small
class-preserving perturbations of each patient's standardized continuous
features, using pairwise Jaccard agreement between the resulting explanatory
edge sets. Categorical inputs and candidate training-node neighborhoods were
held fixed. Fidelity and stability were first averaged separately for stroke
and non-stroke patients and then macro-averaged to prevent majority-class
dominance.
}

 {
Controlled edge ablation further compared removal of the highest-attention,
lowest-attention, and randomly selected incoming relations using the same
patient-specific budget \(b_q\). \textcolor{red}{
Random-edge ablation was repeated ten independent times for each outer-test
patient. Prediction measures were first averaged across the ten random
realizations for that patient and were subsequently aggregated at fold level,
thereby reducing dependence on a single random edge-removal realization.
}
}
%=========================================================
%=========================================================
%-------------------------
%=========================================================
%=========================================================
\section{Experimental Results and Validation}
\label{sec:results}

This section evaluates the proposed \textbf{GNN-DHRS} framework on DF-1,
DF-2, and DF-3. The experiments assess predictive performance,
cross-benchmark transfer, comparisons with existing and controlled baselines,
DHRS and graph-attention contributions, sensitivity and calibration,
explainability, statistical reliability, computational efficiency, and
robustness.

%---------------------------------------------------------
\subsection{Experimental Setup}
\label{subsec:experimental_setup_results}

All experiments followed the leakage-safe nested cross-validation protocol
described in Section~\ref{subsec:data_preparation}, with preprocessing,
feature selection, DHRS, graph construction, model selection, and evaluation
performed independently within the corresponding folds. Resampling was
restricted to training data. GNN-DHRS was implemented using Python, PyTorch,
and PyTorch Geometric, while the conventional baselines used Scikit-learn,
LightGBM, XGBoost, and CatBoost.

{
The complete model configuration and optimization settings are reported in
Table~\ref{tab:gat_configuration}. All hyperparameter selection,
learning-rate scheduling, early stopping, checkpoint selection, and
decision-threshold optimization were performed using inner-validation data
only. The corresponding outer-test folds remained completely untouched until
final evaluation.
}

{
To assess stochastic stability, the experiments were repeated using the
predefined seeds
\(\mathcal{S}=\{11,22,33,44,55\}\). The same outer-fold assignments were
retained across repetitions and compared methods so that variability reflected
stochastic training effects rather than changes in the evaluation partitions.
The complete random-seed stability analysis is reported in
\textcolor{red}{Table~\ref{tab:random_seed_stability}.}
}

\textcolor{red}{
Stochastic robustness was assessed across five predefined seeds while
preserving identical outer-fold assignments to ensure a controlled and directly
comparable evaluation. The fold-level results characterize performance across
the ten outer-test partitions, whereas the five-seed aggregate quantifies the
stability of the selected GNN-DHRS configuration under repeated stochastic
training. The consistently strong performance across seeds indicates that the
reported results are not driven by a particular initialization, DHRS sampling
realization, dropout pattern, or optimization trajectory.
}
 {
All reported test measures were computed exclusively from original,
previously unseen outer-test patients using the strictly inductive inference
protocol described in
Section~\ref{subsec:inductive_graph_construction}. Validation patients and
synthetic observations were excluded from test-metric computation, and no
test--test edges were constructed during inference.
}

 {
Experiments were executed in Kaggle Notebooks using a single NVIDIA Tesla T4
GPU with 16~GB of memory. Detailed hardware, software, and package-version
specifications are provided in \textcolor{red}{Supplementary Table~S2}. Computational timing,
memory utilization, and per-patient inference latency are reported together
with the corresponding efficiency and scalability analysis in
Section~\ref{subsec:computational_efficiency}.
}

\textcolor{red}{
Quantitative figures were generated programmatically in Python using
Matplotlib, while conceptual and methodological schematics were prepared
using diagrams.net (draw.io). This workflow was used consistently throughout
the manuscript to support reproducibility of the reported visual analyses.
}
%=========================================================
%=========================================================
%-----------------------------------------
%-----------------------------------------
\subsection{Overall Performance Across Datasets}
\label{subsec:overall_performance}

This subsection reports the predictive performance of GNN-DHRS on the three
benchmark datasets under the leakage-safe nested cross-validation protocol.
Because the datasets exhibit severe class imbalance, the evaluation combines
conventional classification measures with imbalance-sensitive metrics that
characterize minority-class detection, majority-class recognition, and
agreement beyond chance.

Table~\ref{tab:gnndhrs_overall_results} summarizes the predictive performance
of GNN-DHRS on DF-1, DF-2, and DF-3. The framework attained Accuracy above
\(98.8\%\), F1-score above \(98.3\%\), ROC-AUC above \(99.9\%\), and Cohen's
Kappa above \(96.3\%\).

 {
In addition to these conventional measures, PR-AUC, MCC, balanced accuracy,
specificity, and NPV were computed to characterize minority- and
majority-class performance under severe imbalance. All metrics were calculated
separately for each of the ten outer-test folds and are reported as mean
\(\pm\) standard deviation.
}

 {
All threshold-dependent measures within a given outer fold were computed from
the same set of patient-level predictions using the decision threshold selected
exclusively from the corresponding inner-validation data. Accuracy, Precision,
Recall, F1-score, and specificity were therefore evaluated on the same
untouched outer-test patients. The observed stroke prevalence corresponds to
that same outer-test population, whereas ROC-AUC and PR-AUC were derived from
the continuous prediction scores before thresholding. Complete fold-level
results are provided in \textcolor{red}{Supplementary Table~S3}.
}

 {
During nested model selection, preprocessing and DHRS were applied only to the
corresponding inner-training partition, while the inner-validation patients
remained original and untouched. After hyperparameter selection, the final
model for each outer fold was refitted using the corresponding outer-training
partition, with DHRS again restricted to training data. No synthetic or
resampled observation entered an inner-validation or outer-test fold. Final
evaluation was therefore performed exclusively on original, previously unseen
outer-test patients.
}

\textcolor{blue}{
For clarity, the complete evaluation of GNN-DHRS reports the full set of
conventional and imbalance-sensitive measures, whereas individual comparative
and ablation experiments present the subset of metrics most directly relevant
to the specific objective of each analysis. Accordingly, metric columns are
not required to be identical across all comparison tables. For example, the
resampling-strategy analysis emphasizes Recall, F1-score, PR-AUC, MCC,
balanced accuracy, and specificity to characterize minority-class recovery
and class-specific discrimination, whereas the advanced-baseline comparison
reports Recall, F1-score, PR-AUC, MCC, balanced accuracy, and ROC-AUC to
compare overall discrimination and imbalance-sensitive performance across
heterogeneous learning approaches. All reported metrics were nevertheless
computed from the same leakage-safe outer-test evaluation protocol.
}

\begin{table*}[!ht]
\centering
\caption{
Overall performance of GNN-DHRS across DF-1, DF-2, and DF-3.
 {
Panel A presents conventional performance measures, and Panel B presents
additional imbalance-sensitive measures. Values are mean \(\pm\) standard
deviation across ten untouched outer-test folds.
}}
\label{tab:gnndhrs_overall_results}
\small
\renewcommand{\arraystretch}{1.12}

\textbf{Panel A: Conventional performance measures}

\vspace{3pt}
\setlength{\tabcolsep}{5pt}

\resizebox{\textwidth}{!}{%
\begin{tabular}{@{}lcccccc@{}}
\toprule
\textbf{Dataset} &
\shortstack{\textbf{Accuracy}\\\textbf{(\%)}} &
\shortstack{\textbf{Precision}\\\textbf{(\%)}} &
\shortstack{\textbf{Recall}\\\textbf{(\%)}} &
\shortstack{\textbf{F1-score}\\\textbf{(\%)}} &
\shortstack{\textbf{ROC-AUC}\\\textbf{(\%)}} &
\shortstack{\textbf{Cohen's}\\\textbf{Kappa (\%)}} \\
\midrule

DF-1 &
$99.36\pm0.08$ &
$99.48\pm0.08$ &
$99.31\pm0.15$ &
$99.39\pm0.08$ &
$100.00\pm0.00$ &
$98.72\pm0.16$ \\

DF-2 &
$98.85\pm0.16$ &
$97.84\pm0.18$ &
$98.96\pm0.24$ &
$98.39\pm0.17$ &
$99.92\pm0.08$ &
$96.31\pm0.42$ \\

DF-3 &
$99.43\pm0.42$ &
$97.92\pm0.64$ &
$99.12\pm0.16$ &
$98.52\pm0.41$ &
$99.93\pm0.07$ &
$96.82\pm0.86$ \\

\bottomrule
\end{tabular}%
}

\vspace{7pt}

\begingroup
 

\textbf{Panel B: Imbalance-sensitive performance measures}

\vspace{3pt}
\setlength{\tabcolsep}{7pt}

\begin{tabular}{@{}lccccc@{}}
\toprule
\textbf{Dataset} &
\shortstack{\textbf{PR-AUC}\\\textbf{(\%)}} &
\shortstack{\textbf{MCC}\\\textbf{(\%)}} &
\shortstack{\textbf{Balanced}\\\textbf{Accuracy (\%)}} &
\shortstack{\textbf{Specificity}\\\textbf{(\%)}} &
\shortstack{\textbf{NPV}\\\textbf{(\%)}} \\
\midrule

DF-1 &
$99.87\pm0.04$ &
$98.78\pm0.14$ &
$99.65\pm0.07$ &
$99.99\pm0.01$ &
$99.99\pm0.01$ \\

DF-2 &
$98.71\pm0.19$ &
$96.42\pm0.38$ &
$98.90\pm0.15$ &
$98.84\pm0.17$ &
$99.94\pm0.05$ \\

DF-3 &
$98.84\pm0.22$ &
$96.91\pm0.51$ &
$99.27\pm0.19$ &
$99.42\pm0.28$ &
$99.95\pm0.04$ \\

\bottomrule
\end{tabular}

\endgroup

\vspace{2pt}

\begin{minipage}{0.98\textwidth}
\footnotesize
 {
\textit{Note:} DHRS was restricted to training data within the nested
cross-validation procedure; validation and test folds contained only original
observations. Complete fold-level
results are provided in \textcolor{red}{Supplementary Table~S3}.
}
\end{minipage}

\end{table*}

Among the three datasets, DF-3 attained the highest Accuracy (\(99.43\%\)),
while DF-1 produced \textcolor{red}{
the highest reported combination} of Precision, Recall,
F1-score, and Cohen's Kappa. Recall ranged from \(98.96\%\) to \(99.31\%\),
demonstrating consistently high sensitivity to stroke-positive cases.

 {
Because sensitivity alone does not characterize the false-positive burden,
Recall was interpreted jointly with Precision and specificity. The consistently
high Precision accompanying the high Recall values indicates that increased
stroke-case detection was not achieved through an excessive number of
false-positive predictions.
}

 {
The imbalance-sensitive measures in Panel B further support this
interpretation. PR-AUC ranged from \(98.71\%\) to \(99.87\%\), while MCC
ranged from \(96.42\%\) to \(98.78\%\). Balanced accuracy ranged from
\(98.90\%\) to \(99.65\%\), specificity from \(98.84\%\) to \(99.99\%\),
and NPV exceeded \(99.94\%\) across all three datasets.
}

\textcolor{red}{
The ROC-AUC values of \(100.00\%\), \(99.92\%\), and \(99.93\%\) for DF-1,
DF-2, and DF-3, respectively, demonstrate very high discriminative ability
within the evaluated benchmark settings. These results should nevertheless be
interpreted as benchmark-specific estimates and not as direct evidence of
comparable performance in independently collected clinical cohorts.
}

 {
Taken together, Recall and PR-AUC characterize stroke-positive detection,
specificity and NPV characterize non-stroke recognition, and balanced accuracy
together with MCC summarize performance across both classes. These
complementary measures provide a more informative assessment of performance
under severe class imbalance than Accuracy or ROC-AUC considered in isolation.
}

\textcolor{red}{
The difference between within-dataset and cross-dataset performance highlights
the importance of distributional compatibility between the training and target
benchmarks. Since the patient-similarity graph is constructed from the observed
feature space, variations in predictor distributions, class composition, and
local neighborhood structure can influence the relational patterns learned by
the model. The lower transfer performance therefore indicates that the learned
representation is partly dependent on characteristics of the source benchmark.
Accordingly, the within-dataset results should be interpreted within the scope
of the evaluated data and complemented by validation on independent clinical
cohorts.
}

Figure~\ref{fig:overall_dataset_performance} provides a visual summary of the
aggregate performance of GNN-DHRS across the three evaluated datasets.

\begin{figure}[!h]
\centering
\includegraphics[width=0.90\linewidth]
{images/overall_dataset_performance.png}
\caption{
Overall performance of GNN-DHRS across DF-1, DF-2, and DF-3.
 {
Values summarize predictions from the ten original, untouched outer-test
folds.
}}
\label{fig:overall_dataset_performance}
\end{figure}

{
The broadly consistent performance patterns observed across datasets with
different sample sizes and imbalance ratios indicate stable within-benchmark
behavior under the adopted nested cross-validation framework. However, these
results should be interpreted cautiously because DF-1, DF-2, and DF-3 are
public benchmark datasets rather than independently collected prospective
clinical cohorts. The reported values therefore characterize predictive
performance within the evaluated benchmark distributions and should not be
interpreted as evidence of prospective clinical validity or definitive
external generalizability. Cross-benchmark transportability is examined
separately in Section~\ref{subsec:cross_dataset_results}.
}

%===========================================
%-------------------------------------------
%-------------------------------------------
\subsection{Cross-Benchmark Transfer Analysis}
\label{subsec:cross_dataset_results}

 {
To examine the transportability of GNN-DHRS under differences in dataset
size, class distribution, missingness patterns, and predictor distributions,
directed source-to-target experiments were conducted using the ten predictors
shared by all three benchmarks: age, gender, hypertension, heart disease,
marital status, work type, residence type, average glucose level, body mass
index, and smoking status. Identifier columns were excluded, while variable
names and semantically equivalent categorical values were harmonized using a
predefined mapping. All preprocessing and encoding parameters were estimated
exclusively from the source data. Target-domain categorical levels that were
not represented in the source data were mapped to the source-defined
\texttt{Unknown} category. Because the three public benchmark files are not
documented as independently collected prospective clinical cohorts, this
experiment is interpreted as a cross-benchmark transfer analysis rather than
as independent external clinical validation.
}

 {
For each transfer direction, all model-development operations were performed
using source-domain data only, including preprocessing, feature selection,
DHRS, graph construction, hyperparameter optimization, early stopping, and
decision-threshold selection. DHRS remained restricted to source-side
training data during model development. After training, the fitted
transformations, selected predictors, learned model parameters, source-node
representations, and decision threshold were frozen before any target patient
was evaluated.
}

 {
Each target patient was subsequently processed independently using the
strictly inductive inference procedure described in
Section~\ref{subsec:inductive_graph_construction}. The unseen target patient
was connected only to original source-training nodes, and its prediction was
generated using the frozen source-trained model. No target--target edges,
target labels, target-derived preprocessing parameters, target-domain model
adaptation, or aggregate target-set statistics were used during model
development or inference.
}

 {
All six directed transfer settings were evaluated:
DF-1\(\rightarrow\)DF-2, DF-1\(\rightarrow\)DF-3,
DF-2\(\rightarrow\)DF-1, DF-2\(\rightarrow\)DF-3,
DF-3\(\rightarrow\)DF-1, and DF-3\(\rightarrow\)DF-2.
Evaluating both directions for every dataset pair provides a more stringent
assessment of cross-benchmark transportability than within-dataset
cross-validation alone and also allows asymmetric transfer behavior to be
examined explicitly.
}


\textcolor{red}{
The DF-2\(\leftrightarrow\)DF-3 transfer directions require particular
caution because the two files are related public benchmark representations and
verified patient-level linkage is unavailable. These experiments are
therefore interpreted only as tests of transportability between differently
processed benchmark files and not as independent external replication.
Independent external validation would require a separately collected cohort
with documented patient independence.
}

\begingroup

\begin{table*}[!ht]
\centering
\caption{
Cross-benchmark transfer performance of GNN-DHRS using the common predictor
set. Models were developed exclusively from the source dataset and evaluated
directly on the untouched target dataset without target-domain adaptation.
}
\label{tab:cross_dataset_generalization}
\small
\setlength{\tabcolsep}{5.5pt}
\renewcommand{\arraystretch}{1.12}

\begin{tabular}{@{}llccccc@{}}
\toprule
\textbf{Source} &
\textbf{Target} &
\shortstack{\textbf{Precision}\\\textbf{(\%)}} &
\shortstack{\textbf{Recall}\\\textbf{(\%)}} &
\shortstack{\textbf{F1-score}\\\textbf{(\%)}} &
\shortstack{\textbf{PR-AUC}\\\textbf{(\%)}} &
\shortstack{\textbf{ROC-AUC}\\\textbf{(\%)}} \\
\midrule

DF-1 & DF-2 &
92.47 & 89.36 & 90.89 & 91.18 & 96.82 \\

DF-1 & DF-3 &
91.83 & 88.71 & 90.24 & 90.55 & 96.31 \\

DF-2 & DF-1 &
85.24 & 72.08 & 78.11 & 75.43 & 91.47 \\

DF-2 & DF-3 &
94.12 & 91.28 & 92.68 & 93.01 & 97.54 \\

DF-3 & DF-1 &
84.69 & 70.85 & 77.13 & 74.18 & 90.92 \\

DF-3 & DF-2 &
93.76 & 90.63 & 92.17 & 92.48 & 97.21 \\

\bottomrule
\end{tabular}

\vspace{2pt}

\begin{minipage}{0.96\textwidth}
\footnotesize
\textit{Note:} All data-dependent operations and the decision threshold were
determined using source-domain data only. Target patients were processed
independently using the frozen source-trained pipeline, without target--target
edges, target labels, or target-domain adaptation.
\end{minipage}

\end{table*}

\endgroup

 {
As shown in Table~\ref{tab:cross_dataset_generalization}, cross-benchmark
performance was consistently lower than the corresponding within-dataset
results, revealing a measurable transfer gap under changes in class
prevalence, sample size, missingness patterns, and predictor distributions.
Nevertheless, the source-trained models retained substantial ranking
discrimination and minority-class sensitivity across all six transfer
directions. In particular, the model trained on DF-1 achieved Recall values
of \(89.36\%\) and \(88.71\%\) when transferred to DF-2 and DF-3,
respectively, with corresponding ROC-AUC values of \(96.82\%\) and
\(96.31\%\).
}

 {
Transfer in the opposite direction toward DF-1 was more challenging.
DF-2\(\rightarrow\)DF-1 and DF-3\(\rightarrow\)DF-1 produced Recall values
of \(72.08\%\) and \(70.85\%\), respectively, with PR-AUC values of
\(75.43\%\) and \(74.18\%\). The abovementioned asymmetry may reflect
differences in source-sample size, class imbalance, missingness, and coverage
of the predictor space; however, the present experiment does not establish a
causal explanation for the observed transfer differences.
}

{
By comparison, transfer between DF-2 and DF-3 produced higher performance
than the corresponding transfer directions toward DF-1.
DF-2\(\rightarrow\)DF-3 achieved an F1-score of \(92.68\%\), PR-AUC of
\(93.01\%\), and ROC-AUC of \(97.54\%\), whereas
DF-3\(\rightarrow\)DF-2 achieved corresponding values of \(92.17\%\),
\(92.48\%\), and \(97.21\%\).
}

\textcolor{red}{
Although transfer between DF-2 and DF-3 remained comparatively strong,
these two datasets are related public benchmark representations, and verified
patient-level independence between them cannot be established. Their
bidirectional transfer results should therefore be interpreted as evidence of
cross-benchmark transportability rather than independent external clinical
generalization. More broadly, the reduction from near-ceiling
within-benchmark performance to lower cross-benchmark performance indicates
sensitivity to dataset shift. Because the patient-similarity graph is derived
from the feature distribution of the source dataset, changes in predictor
distributions, missingness patterns, class prevalence, and feature-space
geometry may alter the induced neighborhood structure and, consequently, the
learned relational representation. Thus, part of the strong within-benchmark
performance may depend on benchmark-specific structure. The transfer results
do not establish overfitting in a strict statistical sense, but they show that
the within-benchmark estimates should not be extrapolated directly to unseen
clinical populations.
}

{
Overall, the asymmetric source-to-target results indicate that performance
under dataset shift depends on the characteristics and representativeness of
the source benchmark. GNN-DHRS retained useful predictive capability across
the evaluated public datasets, but its cross-benchmark performance remained
below the within-dataset estimates. The findings therefore provide evidence
of partial transportability across the three harmonized public benchmarks,
while also demonstrating the need for validation on independently collected
clinical cohorts before conclusions regarding external clinical
generalizability can be made.
}

%-------------------------------------------
\subsection{Comparison with State-of-the-Art Methods}
\label{subsec:sota_comparison}

To contextualize the effectiveness of the proposed framework, GNN-DHRS was
compared with representative state-of-the-art stroke-prediction methods,
including conventional machine-learning, ensemble, boosting, deep-learning,
meta-learning, and attention-based approaches.
Table~\ref{tab:sota_comparison2} reports the comparative results across
DF-1, DF-2, and DF-3.

 {
Because the values for existing methods were taken from their respective
published studies, the comparison in Table~\ref{tab:sota_comparison2} should
be interpreted as a literature-level contextual comparison rather than a
strictly controlled head-to-head experiment. The cited studies may differ in
data partitioning, preprocessing, imbalance handling, hyperparameter
optimization, and validation protocol. Accordingly, numerical differences
between GNN-DHRS and previously published methods are reported descriptively,
whereas the controlled same-protocol comparisons are presented separately in
Section~\ref{subsec:advanced_baseline_comparison}.
}

\begin{table*}[!h]
\centering
\caption{
Comparison of GNN-DHRS with representative state-of-the-art
stroke-prediction methods across DF-1, DF-2, and DF-3.
 {
NR denotes a metric not reported in the corresponding cited study.
Published results are reproduced for contextual comparison and were not
necessarily obtained using the same experimental protocol.
}}
\label{tab:sota_comparison2}

\resizebox{\textwidth}{!}{%
\begin{tabular}{llccccc}
\toprule
\textbf{Reference} &
\textbf{Model} &
\textbf{Dataset} &
\textbf{Accuracy (\%)} &
\textbf{Recall (\%)} &
\textbf{F1-score (\%)} &
\textbf{AUC (\%)} \\
\midrule

\multicolumn{7}{c}{\textbf{DF-1}} \\
\midrule

\cite{df1_11}
& XGBoost
& DF-1
& 87.50 & NR & 89.20 & NR \\

\cite{df1_original}
& CatBoost
& DF-1
& 98.90 & NR & 98.00 & NR \\

\textcolor{red}{\cite{abousaber2025metastroke}}
& Meta-Learning
& DF-1
& 99.21 & NR & 99.26 & NR \\

\cite{johnson2023stroke}
& Random Forest
& DF-1
& 96.42 & 78.53 & 82.14 & 95.71 \\

\cite{melnykova2025stroke}
& XGBoost
& DF-1
& 97.11 & 84.72 & 87.23 & 97.12 \\

\textbf{Proposed}
& \textbf{GNN-DHRS}
& \textbf{DF-1}
& \textbf{99.36}
& \textbf{99.31}
& \textbf{99.39}
& \textbf{100.00} \\

\midrule
\multicolumn{7}{c}{\textbf{DF-2}} \\
\midrule

\cite{df2_hybrid_ref29}
& LR, DT, RF, SVM, and NB
& DF-2
& 95.50 & NR & 94.50 & NR \\

\cite{df2_hybrid_ref30}
& Category Boosting Classifier
& DF-2
& 97.00 & NR & 96.00 & NR \\

\cite{df2_hybrid_ref25}
& DT, SVM, and LR
& DF-2
& 95.49 & NR & 96.00 & NR \\

\cite{df2_hybrid_ref31}
& Deep Neural Network
& DF-2
& 95.49 & NR & 96.00 & NR \\

\cite{df2_hybrid}
& Random Forest
& DF-2
& 97.19 & NR & 97.15 & NR \\

\cite{df1_original}
& Stacking Algorithms
& DF-2
& 97.98 & NR & 98.00 & NR \\

\cite{df2_explainable_ref}
& Boosting Algorithms
& DF-2
& 97.97 & NR & 93.00 & NR \\

\textcolor{red}{\cite{abousaber2025metastroke}}
& Meta-Learning
& DF-2
& 98.02 & NR & 98.25 & NR \\

\cite{akib2025optimizing}
& Hybrid Deep Learning
& DF-2
& 98.36 & 82.91 & 85.74 & 97.43 \\

\cite{kanwal2025stroke}
& Attention-Based Model
& DF-2
& 98.74 & 85.61 & 88.02 & 98.06 \\

\textbf{Proposed}
& \textbf{GNN-DHRS}
& \textbf{DF-2}
& \textbf{98.85}
& \textbf{98.96}
& \textbf{98.39}
& \textbf{99.92} \\

\midrule
\multicolumn{7}{c}{\textbf{DF-3}} \\
\midrule

\cite{df3_ref7}
& Random Forest
& DF-3
& 95.30 & NR & NR & NR \\

\cite{df3_ref8}
& BSPE
& DF-3
& 95.30 & NR & NR & NR \\

\cite{df3_ref9}
& LightGBM
& DF-3
& 94.50 & NR & NR & NR \\

\cite{df3_ref10}
& Random Forest
& DF-3
& 98.90 & NR & NR & NR \\

\cite{df3_ref12}
& Random Forest
& DF-3
& 97.20 & NR & NR & NR \\

\cite{df3_ref14}
& Random Forest
& DF-3
& 99.10 & NR & NR & NR \\

\cite{df3_ref16}
& Voting
& DF-3
& 97.00 & NR & NR & NR \\

\cite{df3_ref17}
& Decision Tree
& DF-3
& 93.00 & NR & NR & NR \\

\textcolor{red}{\cite{abousaber2025metastroke}}
& Meta-Learning
& DF-3
& 99.34 & NR & 98.30 & NR \\

\textbf{Proposed}
& \textbf{GNN-DHRS}
& \textbf{DF-3}
& \textbf{99.43}
& \textbf{99.12}
& \textbf{98.52}
& \textbf{99.93} \\

\bottomrule
\end{tabular}%
}
\end{table*}

As shown in Table~\ref{tab:sota_comparison2}, \textcolor{red}{
GNN-DHRS demonstrates competitive reported performance across the three
datasets.
}. On DF-1, the
proposed framework obtains \(99.36\%\) Accuracy, \(99.31\%\) Recall,
\(99.39\%\) F1-score, and \(100.00\%\) AUC, slightly exceeding the recent
meta-learning model in both Accuracy and F1-score. More importantly, its
reported Recall is  higher than that of the Random Forest and
XGBoost methods for which Recall values are available, i\textcolor{red}{
corresponding to higher reported sensitivity to stroke-positive cases in the
literature-level comparison.
}

{
The DF-1 comparison should nevertheless be interpreted with caution because
several published studies did not report Recall or AUC, and the underlying
validation protocols were not identical. Thus, the table establishes the
relative position of GNN-DHRS among reported results but does not by itself
constitute evidence of statistical superiority over every published method.
}

\textcolor{red}{
On DF-2,
} where GNN-DHRS reports
\(98.85\%\) Accuracy, \(98.96\%\) Recall, \(98.39\%\) F1-score, and
\(99.92\%\) AUC. Compared with the reported attention-based model, Recall
increases from \(85.61\%\) to \(98.96\%\), corresponding to a difference of
\(13.35\) percentage points, while F1-score increases from \(88.02\%\) to
\(98.39\%\), corresponding to \(10.37\) percentage points.

 {
These differences are particularly relevant for evaluating severely
imbalanced stroke data because Recall and F1-score provide substantially more
information about minority-class detection than Accuracy alone. However,
because these results originate from separately conducted studies, the
magnitude of the reported differences should not be interpreted as a
controlled treatment effect attributable solely to the proposed framework.
}

 {
On DF-3, GNN-DHRS achieved \(99.43\%\) Accuracy, \(99.12\%\) Recall,
\(98.52\%\) F1-score, and \(99.93\%\) AUC. Its reported Accuracy and
F1-score were slightly higher than those of the available DF-3 comparison
studies in Table~\ref{tab:sota_comparison2}. In particular, the closest
reported meta-learning approach achieved \(99.34\%\) Accuracy and
\(98.30\%\) F1-score. However, because most previous DF-3 studies did not
report Recall or AUC, direct conclusions for these measures are restricted to
methods providing the corresponding metrics.
}

\textcolor{red}{
Overall, the literature-level comparison shows that GNN-DHRS reports
consistently high Recall and F1-score values across the three evaluated
datasets.
}. Many
previous studies report high Accuracy but omit Recall, which limits assessment
of minority stroke-case detection under severe class imbalance. In contrast,
GNN-DHRS reports consistently high Recall, F1-score, and AUC across all three
datasets.

{
The literature comparison therefore provides contextual evidence that the
proposed framework is competitive with recent stroke-prediction approaches,
particularly with respect to minority-class detection. Because the published
methods were not evaluated under a common experimental protocol, attribution
of the observed gains specifically to DHRS, weighted patient-similarity graph
construction, or multi-head graph-attention learning is deferred to the
controlled baseline and ablation experiments presented in the subsequent
sections.
}

\textcolor{red}{
Accordingly, Table~\ref{tab:sota_comparison2} is used solely to contextualize
the reported benchmark performance. Comparative claims regarding the relative
behavior of GNN-DHRS are reserved for the controlled same-protocol
experiments in Section~\ref{subsec:advanced_baseline_comparison}, where the
competing methods use identical data partitions, preprocessing rules,
inner-validation procedures, and outer-test evaluation conditions.
}

%==========================

%==========================
%-------------------------------------------
%-------------------------------------------
\subsection{Classifier-Wise Comparison}
\label{subsec:classifierwise_comparison}

To further evaluate the effectiveness of the proposed framework, GNN-DHRS was
compared with conventional machine-learning, ensemble-learning, neural-network,
and graph-based classifiers under the same preprocessing, DHRS resampling,
feature-selection, and nested cross-validation protocol.

 {
To ensure a fair comparison under severe class imbalance, all classifiers were
additionally evaluated using PR-AUC, MCC, balanced accuracy, specificity, and
NPV. The same outer folds, training-only preprocessing operations, and
validation-selected decision thresholds were used for every method. All
results are reported as mean \(\pm\) standard deviation across the untouched
outer-test folds.
}

%=========================================================
\begin{table*}[!ht]
\centering
\caption{
Classifier-wise performance under DHRS.
 {
Panel A reports conventional metrics and Panel B imbalance-sensitive metrics;
values are mean \(\pm\) standard deviation (\%).
}}
\label{tab:classifierwise_performance2}
\label{tab:classifier_imbalance_metrics}

\scriptsize
\setlength{\tabcolsep}{3.2pt}
\renewcommand{\arraystretch}{0.94}

%=========================================================
% Panel A: Conventional metrics
%=========================================================
\textbf{Panel A: Conventional metrics}

\vspace{2pt}

\resizebox{\textwidth}{!}{%
\begin{tabular}{@{}llcccccc@{}}
\toprule
\textbf{Data} &
\textbf{Classifier} &
\textbf{Acc.} &
\textbf{Prec.} &
\textbf{Rec.} &
\textbf{F1} &
\textbf{AUC} &
\textbf{Kappa} \\
\midrule

\multirow{8}{*}{DF-1}
& DT
& $96.85\pm0.52$
& $89.74\pm1.18$
& $82.18\pm1.45$
& $85.79\pm1.21$
& $95.84\pm0.67$
& $83.91\pm1.09$ \\

& RF
& $97.72\pm0.41$
& $92.83\pm0.89$
& $88.17\pm1.02$
& $90.44\pm0.87$
& $98.16\pm0.38$
& $89.37\pm0.78$ \\

& ET
& $97.84\pm0.38$
& $93.18\pm0.82$
& $88.91\pm0.95$
& $90.99\pm0.81$
& $98.34\pm0.35$
& $90.11\pm0.72$ \\

& GB
& $97.91\pm0.36$
& $93.54\pm0.78$
& $89.33\pm0.91$
& $91.38\pm0.76$
& $98.41\pm0.33$
& $90.52\pm0.69$ \\

& LightGBM
& $98.07\pm0.33$
& $94.22\pm0.71$
& $90.11\pm0.84$
& $92.12\pm0.70$
& $98.58\pm0.29$
& $91.34\pm0.63$ \\

& MLP
& $97.56\pm0.45$
& $91.61\pm0.96$
& $87.42\pm1.11$
& $89.47\pm0.94$
& $97.93\pm0.42$
& $88.26\pm0.85$ \\

& GCN
& $98.16\pm0.31$
& $94.73\pm0.68$
& $91.08\pm0.79$
& $92.87\pm0.65$
& $98.76\pm0.27$
& $92.03\pm0.58$ \\

& \textbf{GNN-DHRS}
& \textbf{$99.36\pm0.08$}
& \textbf{$99.48\pm0.08$}
& \textbf{$99.31\pm0.15$}
& \textbf{$99.39\pm0.08$}
& \textbf{$100.00\pm0.00$}
& \textbf{$98.72\pm0.16$} \\

\midrule

\multirow{8}{*}{DF-2}
& DT
& $97.92\pm0.48$
& $84.33\pm1.25$
& $73.17\pm1.58$
& $78.35\pm1.32$
& $94.26\pm0.71$
& $75.81\pm1.19$ \\

& RF
& $98.61\pm0.37$
& $89.12\pm0.94$
& $82.83\pm1.08$
& $85.86\pm0.91$
& $97.34\pm0.41$
& $84.01\pm0.82$ \\

& ET
& $98.74\pm0.34$
& $89.94\pm0.87$
& $83.72\pm0.99$
& $86.72\pm0.84$
& $97.61\pm0.37$
& $84.93\pm0.76$ \\

& GB
& $98.88\pm0.31$
& $90.57\pm0.81$
& $84.84\pm0.93$
& $87.61\pm0.78$
& $97.89\pm0.34$
& $86.08\pm0.71$ \\

& LightGBM
& $99.02\pm0.28$
& $91.34\pm0.74$
& $85.73\pm0.86$
& $88.44\pm0.72$
& $98.11\pm0.30$
& $87.19\pm0.65$ \\

& MLP
& $98.51\pm0.42$
& $88.25\pm1.01$
& $81.69\pm1.15$
& $84.84\pm0.97$
& $97.08\pm0.45$
& $82.93\pm0.89$ \\

& GCN
& $98.11\pm0.35$
& $92.08\pm0.79$
& $87.16\pm0.88$
& $89.55\pm0.74$
& $98.47\pm0.32$
& $88.81\pm0.68$ \\

& \textbf{GNN-DHRS}
& \textbf{$98.85\pm0.16$}
& \textbf{$97.84\pm0.18$}
& \textbf{$98.96\pm0.24$}
& \textbf{$98.39\pm0.17$}
& \textbf{$99.92\pm0.08$}
& \textbf{$96.31\pm0.42$} \\

\midrule

\multirow{8}{*}{DF-3}
& DT
& $97.58\pm0.51$
& $83.55\pm1.29$
& $74.62\pm1.61$
& $78.83\pm1.35$
& $94.52\pm0.73$
& $76.44\pm1.22$ \\

& RF
& $98.23\pm0.39$
& $87.61\pm0.97$
& $81.47\pm1.12$
& $84.43\pm0.94$
& $96.97\pm0.43$
& $82.51\pm0.85$ \\

& ET
& $98.37\pm0.36$
& $88.26\pm0.90$
& $82.41\pm1.03$
& $85.23\pm0.87$
& $97.24\pm0.39$
& $83.48\pm0.79$ \\

& GB
& $98.51\pm0.33$
& $89.12\pm0.84$
& $83.58\pm0.96$
& $86.26\pm0.81$
& $97.51\pm0.35$
& $84.56\pm0.73$ \\

& LightGBM
& $98.66\pm0.30$
& $89.88\pm0.77$
& $84.25\pm0.89$
& $86.97\pm0.75$
& $97.83\pm0.31$
& $85.48\pm0.67$ \\

& MLP
& $98.16\pm0.44$
& $86.94\pm1.04$
& $80.92\pm1.18$
& $83.82\pm1.00$
& $96.73\pm0.47$
& $81.87\pm0.91$ \\

& GCN
& $98.81\pm0.32$
& $91.02\pm0.76$
& $86.17\pm0.85$
& $88.53\pm0.72$
& $98.06\pm0.29$
& $87.71\pm0.66$ \\

& \textbf{GNN-DHRS}
& \textbf{$99.43\pm0.42$}
& \textbf{$97.92\pm0.64$}
& \textbf{$99.12\pm0.16$}
& \textbf{$98.52\pm0.41$}
& \textbf{$99.93\pm0.07$}
& \textbf{$96.82\pm0.86$} \\

\bottomrule
\end{tabular}%
}

\vspace{4pt}

%=========================================================
% Panel B: Imbalance-sensitive metrics
%=========================================================
\begingroup
 

\textbf{Panel B: Imbalance-sensitive metrics}

\vspace{2pt}

\begin{tabular}{@{}llccccc@{}}
\toprule
\textbf{Data} &
\textbf{Classifier} &
\textbf{PR-AUC} &
\textbf{MCC} &
\textbf{BAcc.} &
\textbf{Spec.} &
\textbf{NPV} \\
\midrule

\multirow{8}{*}{DF-1}
& DT
& $84.12\pm1.28$
& $84.05\pm1.15$
& $89.48\pm0.92$
& $96.78\pm0.41$
& $99.65\pm0.12$ \\

& RF
& $90.87\pm0.87$
& $89.51\pm0.79$
& $92.95\pm0.68$
& $97.73\pm0.29$
& $99.76\pm0.09$ \\

& ET
& $91.64\pm0.81$
& $90.28\pm0.74$
& $93.38\pm0.63$
& $97.85\pm0.27$
& $99.78\pm0.08$ \\

& GB
& $92.21\pm0.76$
& $90.69\pm0.71$
& $93.62\pm0.59$
& $97.91\pm0.25$
& $99.80\pm0.08$ \\

& LightGBM
& $93.08\pm0.69$
& $91.52\pm0.65$
& $94.09\pm0.54$
& $98.07\pm0.23$
& $99.82\pm0.07$ \\

& MLP
& $89.73\pm0.94$
& $88.41\pm0.88$
& $92.49\pm0.73$
& $97.57\pm0.32$
& $99.74\pm0.10$ \\

& GCN
& $93.85\pm0.61$
& $92.18\pm0.58$
& $94.62\pm0.49$
& $98.17\pm0.21$
& $99.84\pm0.06$ \\

& \textbf{GNN-DHRS}
& \textbf{$99.87\pm0.04$}
& \textbf{$98.78\pm0.14$}
& \textbf{$99.65\pm0.07$}
& \textbf{$99.99\pm0.01$}
& \textbf{$99.99\pm0.01$} \\

\midrule

\multirow{8}{*}{DF-2}
& DT
& $76.48\pm1.41$
& $76.12\pm1.33$
& $85.55\pm1.08$
& $97.93\pm0.38$
& $98.72\pm0.21$ \\

& RF
& $85.21\pm0.95$
& $84.28\pm0.89$
& $90.72\pm0.76$
& $98.63\pm0.27$
& $99.18\pm0.14$ \\

& ET
& $86.34\pm0.89$
& $85.19\pm0.83$
& $91.23\pm0.71$
& $98.74\pm0.25$
& $99.23\pm0.13$ \\

& GB
& $87.52\pm0.82$
& $86.35\pm0.77$
& $91.86\pm0.65$
& $98.88\pm0.23$
& $99.28\pm0.12$ \\

& LightGBM
& $88.41\pm0.76$
& $87.47\pm0.71$
& $92.38\pm0.59$
& $99.02\pm0.21$
& $99.33\pm0.11$ \\

& MLP
& $84.05\pm1.02$
& $83.17\pm0.96$
& $90.10\pm0.81$
& $98.52\pm0.29$
& $99.12\pm0.15$ \\

& GCN
& $89.67\pm0.68$
& $88.95\pm0.64$
& $92.64\pm0.53$
& $98.12\pm0.24$
& $99.41\pm0.10$ \\

& \textbf{GNN-DHRS}
& \textbf{$98.71\pm0.19$}
& \textbf{$96.42\pm0.38$}
& \textbf{$98.90\pm0.15$}
& \textbf{$98.84\pm0.17$}
& \textbf{$99.94\pm0.05$} \\

\midrule

\multirow{8}{*}{DF-3}
& DT
& $77.19\pm1.37$
& $76.68\pm1.29$
& $86.10\pm1.05$
& $97.60\pm0.39$
& $98.78\pm0.20$ \\

& RF
& $83.92\pm0.98$
& $82.79\pm0.92$
& $89.85\pm0.78$
& $98.25\pm0.28$
& $99.11\pm0.14$ \\

& ET
& $84.87\pm0.91$
& $83.76\pm0.85$
& $90.39\pm0.72$
& $98.37\pm0.26$
& $99.16\pm0.13$ \\

& GB
& $86.03\pm0.84$
& $84.84\pm0.79$
& $91.05\pm0.66$
& $98.51\pm0.24$
& $99.21\pm0.12$ \\

& LightGBM
& $86.94\pm0.78$
& $85.76\pm0.73$
& $91.46\pm0.61$
& $98.66\pm0.22$
& $99.25\pm0.11$ \\

& MLP
& $83.15\pm1.05$
& $82.13\pm0.99$
& $89.54\pm0.83$
& $98.17\pm0.30$
& $99.07\pm0.15$ \\

& GCN
& $88.41\pm0.71$
& $87.95\pm0.67$
& $92.49\pm0.55$
& $98.82\pm0.23$
& $99.36\pm0.10$ \\

& \textbf{GNN-DHRS}
& \textbf{$98.84\pm0.22$}
& \textbf{$96.91\pm0.51$}
& \textbf{$99.27\pm0.19$}
& \textbf{$99.42\pm0.28$}
& \textbf{$99.95\pm0.04$} \\

\bottomrule
\end{tabular}

\endgroup

\vspace{1pt}

\begin{minipage}{0.98\textwidth}
\footnotesize
 {
\textit{Note:} BAcc.: balanced accuracy; Spec.: specificity. All classifiers
used identical outer folds, training-only preprocessing, DHRS, and
validation-selected thresholds.
}
\end{minipage}

\end{table*}
%=========================================================

As shown in Table~\ref{tab:classifierwise_performance2}, GNN-DHRS consistently
\textcolor{red}{
achieves consistently high performance
} across the evaluated classifier
families, with particularly large gains in Recall and F1-score. The advantage
is evident across all three datasets despite the differences in dataset size
and imbalance ratio.

 {
The imbalance-sensitive measures in Panel B reinforce the same pattern.
Relative to GCN, the strongest graph-based comparator in this experiment,
GNN-DHRS improved PR-AUC by approximately \(6.02\), \(9.04\), and
\(10.43\) percentage points on DF-1, DF-2, and DF-3, respectively. The
corresponding MCC improvements were \(6.60\), \(7.47\), and \(8.96\)
percentage points, while balanced accuracy increased by \(5.03\), \(6.26\),
and \(6.78\) percentage points.
}

 {
GNN-DHRS also achieved the highest NPV on all three datasets and the highest
specificity on DF-1 and DF-3. On DF-2, LightGBM obtained a slightly higher
specificity (\(99.02\pm0.21\%\)) than GNN-DHRS
(\(98.84\pm0.17\%\)); however, GNN-DHRS achieved substantially higher
Recall, F1-score, PR-AUC, MCC, balanced accuracy, and NPV. This result
illustrates the trade-off between maximizing majority-class specificity and
maintaining strong stroke-positive sensitivity.
}

The performance advantage is particularly evident when compared with GCN.
On DF-1, Recall improves from \(91.08\pm0.79\%\) to
\(99.31\pm0.15\%\), while F1-score increases from
\(92.87\pm0.65\%\) to \(99.39\pm0.08\%\). Similar improvements are observed
on DF-2, where Recall increases from \(87.16\pm0.88\%\) to
\(98.96\pm0.24\%\) and F1-score from \(89.55\pm0.74\%\) to
\(98.39\pm0.17\%\), and on DF-3, where Recall improves from
\(86.17\pm0.85\%\) to \(99.12\pm0.16\%\) and F1-score from
\(88.53\pm0.72\%\) to \(98.52\pm0.41\%\).

A clear performance trend can also be observed across classifier families.
Tree-based and neural-network models achieve competitive Accuracy but exhibit
noticeably lower Recall. Ensemble and boosting approaches improve overall
discrimination; however, their sensitivity to minority stroke cases remains
substantially below that of GNN-DHRS. Although GCN benefits from relational
patient modelling, its performance remains consistently below that of the
complete proposed framework.

\textcolor{red}{
These findings suggest that high Accuracy alone does not adequately
characterize predictive quality under severe class imbalance. More stringent
controlled comparisons with class-weighted tabular models,
transformer-based tabular learning, and alternative graph-learning
architectures are presented in
Section~\ref{subsec:advanced_baseline_comparison}.
}

%-------------------------------------------
%=========================================================
\subsection{
Controlled Evaluation of Baseline Models and Imbalance-Handling Strategies}
\label{subsec:advanced_baseline_comparison}
\label{subsec:controlled_imbalance_analysis}

 {
This subsection presents a coordinated set of controlled experiments designed
to distinguish three potential sources of predictive improvement: classifier
architecture, loss-level imbalance correction, and training-data refinement.
\textcolor{red}{
First, GNN-DHRS is compared with conventional class-weighted tabular models,
a modern transformer-based tabular architecture, alternative graph
architectures, and cost-sensitive GAT variants.
} Second, the
graph-attention classifier is held fixed while alternative resampling
strategies and DHRS component-ordering ablations are evaluated. Third, the
implementation of DHRS is examined through stage-wise class distributions,
categorical-variable treatment, and post-resampling class weighting.
}

 {
Across all comparisons, the outer folds, inner-validation partitions,
training-only preprocessing, selected features, graph settings, optimization
parameters, random seeds, and validation-selected decision thresholds were
held constant unless explicitly identified as the experimental factor.
Resampling was restricted to the corresponding training partition, and all
predictive measures were calculated from original, untouched outer-test
observations.
}

%=========================================================
\subsubsection{
Architectural and Loss-Level Baselines}
\label{subsubsec:advanced_baselines}


The first experiment evaluates whether the performance of GNN-DHRS can be
reproduced by alternative tabular and graph-learning architectures under the
same leakage-safe evaluation protocol. In addition to class-weighted
XGBoost~\cite{chen2016xgboost}, LightGBM~\cite{ke2017lightgbm}, and
CatBoost~\cite{prokhorenkova2018catboost}, \textcolor{red}{FT-Transformer was included as a
modern attention-based architecture specifically developed for tabular
learning~\cite{gorishniy2021revisiting}}. The graph baselines comprise
GCN~\cite{kipf2017semi}, GraphSAGE~\cite{hamilton2017inductive}, and
GAT~\cite{velivckovic2018graph}.


\textcolor{red}{
FT-Transformer was evaluated under the same leakage-safe protocol as the
controlled baselines, using identical outer-test folds, inner-validation
partitions, training-only preprocessing, selected predictors, random seeds,
and validation-based model selection. Two FT-Transformer configurations were
considered: FT-Transformer-W, using class-weighted learning without
resampling, and FT-Transformer-DHRS, using the training-only DHRS procedure
adopted in GNN-DHRS. This comparison isolates the effect of DHRS within a
tabular self-attention architecture and provides a controlled contrast with
graph-based relational learning}. Among the remaining baselines, XGBoost,
LightGBM, and CatBoost used class weighting; GCN and GraphSAGE used weighted
binary cross-entropy (WBCE); and GAT was evaluated with BCE, WBCE, focal loss
(FL), class-balanced loss (CB), and class-balanced focal loss (CBF). All
competing graph configurations excluded DHRS, whereas the complete GNN-DHRS
framework combined training-only DHRS, cosine-weighted patient-graph
construction, multi-head graph attention, and WBCE.


\begingroup

\begin{table*}[!ht]
\centering
\caption{
Controlled comparison of advanced tabular, transformer-based, graph, and
imbalance-aware baselines with GNN-DHRS using a common set of
minority-sensitive, class-balanced, and discrimination measures. Values are
mean \(\pm\) standard deviation (\%) across ten untouched outer-test folds.
}
\label{tab:advanced_imbalance_baselines}

\scriptsize
\setlength{\tabcolsep}{1.8pt}
\renewcommand{\arraystretch}{0.98}

\resizebox{\textwidth}{!}{%
\begin{tabular}{@{}llcccccccc@{}}
\toprule
\textbf{Dataset} &
\textbf{Configuration} &
\textbf{Recall} &
\textbf{F1-score} &
\textbf{PR-AUC} &
\textbf{MCC} &
\textbf{BAcc.} &
\textbf{Specificity} &
\textbf{NPV} &
\textbf{ROC-AUC} \\
\midrule

% =========================================================
% DF-1
% =========================================================
\multirow{13}{*}{DF-1}

& XGB-CW
& $91.24\pm0.78$
& $93.18\pm0.65$
& $94.05\pm0.58$
& $92.41\pm0.71$
& $95.37\pm0.49$
& \textcolor{red}{$99.50\pm1.25$}
& \textcolor{red}{$99.84\pm0.01$}
& $98.62\pm0.31$ \\

& LGBM-CW
& $91.87\pm0.71$
& $93.64\pm0.59$
& $94.51\pm0.52$
& $92.89\pm0.64$
& $95.68\pm0.45$
& \textcolor{red}{$99.49\pm1.15$}
& \textcolor{red}{$99.85\pm0.01$}
& $98.79\pm0.27$ \\

& CatB-CW
& $92.35\pm0.67$
& $94.02\pm0.55$
& $94.88\pm0.48$
& $93.27\pm0.59$
& $95.92\pm0.42$
& \textcolor{red}{$99.49\pm1.07$}
& \textcolor{red}{$99.86\pm0.01$}
& $98.91\pm0.24$ \\

% -------- NEW TRANSFORMER BASELINES --------

& \textcolor{red}{FT-T-W}
& \textcolor{red}{$94.10\pm0.60$}
& \textcolor{red}{$95.02\pm0.51$}
& \textcolor{red}{$95.85\pm0.45$}
& \textcolor{red}{$94.41\pm0.54$}
& \textcolor{red}{$96.67\pm0.39$}
& \textcolor{red}{$99.24\pm0.88$}
& \textcolor{red}{$99.90\pm0.01$}
& \textcolor{red}{$99.20\pm0.18$} \\

& \textcolor{red}{FT-T-DHRS}
& \textcolor{red}{$96.85\pm0.42$}
& \textcolor{red}{$97.10\pm0.36$}
& \textcolor{red}{$97.55\pm0.31$}
& \textcolor{red}{$96.38\pm0.40$}
& \textcolor{red}{$98.05\pm0.28$}
& \textcolor{red}{$99.25\pm0.67$}
& \textcolor{red}{$99.95\pm0.01$}
& \textcolor{red}{$99.55\pm0.12$} \\

\addlinespace[1pt]

& GCN-W
& $91.08\pm0.79$
& $92.87\pm0.65$
& $93.85\pm0.61$
& $92.18\pm0.58$
& $94.62\pm0.49$
& \textcolor{red}{$98.16\pm1.26$}
& \textcolor{red}{$99.83\pm0.01$}
& $98.76\pm0.27$ \\

& SAGE-W
& $92.14\pm0.73$
& $93.51\pm0.61$
& $94.32\pm0.54$
& $92.76\pm0.62$
& $95.41\pm0.46$
& \textcolor{red}{$98.68\pm1.17$}
& \textcolor{red}{$99.85\pm0.01$}
& $98.84\pm0.25$ \\

\addlinespace[1pt]

& GAT-BCE
& $90.27\pm0.84$
& $92.15\pm0.71$
& $93.18\pm0.63$
& $91.42\pm0.76$
& $94.81\pm0.54$
& \textcolor{red}{$99.35\pm1.37$}
& \textcolor{red}{$99.82\pm0.02$}
& $98.41\pm0.33$ \\

& GAT-W
& $92.68\pm0.69$
& $94.03\pm0.57$
& $94.87\pm0.49$
& $93.21\pm0.61$
& $95.89\pm0.43$
& \textcolor{red}{$99.10\pm1.10$}
& \textcolor{red}{$99.86\pm0.01$}
& $98.93\pm0.25$ \\

& GAT-FL
& $93.15\pm0.65$
& $94.41\pm0.53$
& $95.22\pm0.46$
& $93.64\pm0.56$
& $96.17\pm0.40$
& \textcolor{red}{$99.19\pm1.03$}
& \textcolor{red}{$99.87\pm0.01$}
& $99.01\pm0.23$ \\

& GAT-CB
& $93.42\pm0.61$
& $94.58\pm0.50$
& $95.39\pm0.43$
& $93.87\pm0.52$
& $96.31\pm0.38$
& \textcolor{red}{$99.20\pm0.97$}
& \textcolor{red}{$99.88\pm0.01$}
& $99.06\pm0.21$ \\

& GAT-CBF
& $93.89\pm0.58$
& $94.92\pm0.47$
& $95.71\pm0.41$
& $94.28\pm0.49$
& $96.54\pm0.36$
& \textcolor{red}{$99.19\pm0.92$}
& \textcolor{red}{$99.89\pm0.01$}
& $99.12\pm0.19$ \\

& \textbf{GNN-DHRS}
& \textbf{$99.31\pm0.15$}
& \textbf{$99.39\pm0.08$}
& \textbf{$99.87\pm0.04$}
& \textbf{$98.78\pm0.14$}
& \textbf{$99.65\pm0.07$}
& \textcolor{red}{\textbf{$99.99\pm0.01$}}
& \textcolor{red}{\textbf{$99.99\pm0.01$}}
& \textbf{$100.00\pm0.00$} \\

\midrule

% =========================================================
% DF-2
% =========================================================
\multirow{13}{*}{DF-2}

& XGB-CW
& $86.42\pm0.91$
& $88.15\pm0.78$
& $89.27\pm0.69$
& $86.94\pm0.82$
& $92.18\pm0.61$
& \textcolor{red}{$97.94\pm1.52$}
& \textcolor{red}{$99.28\pm0.05$}
& $97.85\pm0.34$ \\

& LGBM-CW
& $87.19\pm0.84$
& $88.73\pm0.72$
& $89.84\pm0.63$
& $87.58\pm0.75$
& $92.61\pm0.56$
& \textcolor{red}{$98.03\pm1.40$}
& \textcolor{red}{$99.32\pm0.05$}
& $98.07\pm0.30$ \\

& CatB-CW
& $87.85\pm0.79$
& $89.21\pm0.68$
& $90.31\pm0.58$
& $88.12\pm0.71$
& $92.94\pm0.52$
& \textcolor{red}{$98.03\pm1.31$}
& \textcolor{red}{$99.35\pm0.04$}
& $98.24\pm0.27$ \\

% -------- NEW TRANSFORMER BASELINES --------

& \textcolor{red}{FT-T-W}
& \textcolor{red}{$90.10\pm0.74$}
& \textcolor{red}{$91.42\pm0.63$}
& \textcolor{red}{$92.18\pm0.57$}
& \textcolor{red}{$90.76\pm0.66$}
& \textcolor{red}{$94.12\pm0.47$}
& \textcolor{red}{$98.14\pm1.08$}
& \textcolor{red}{$99.47\pm0.04$}
& \textcolor{red}{$98.91\pm0.21$} \\

& \textcolor{red}{FT-T-DHRS}
& \textcolor{red}{$95.42\pm0.47$}
& \textcolor{red}{$95.11\pm0.40$}
& \textcolor{red}{$95.86\pm0.36$}
& \textcolor{red}{$93.81\pm0.49$}
& \textcolor{red}{$96.72\pm0.33$}
& \textcolor{red}{$98.02\pm0.79$}
& \textcolor{red}{$99.76\pm0.03$}
& \textcolor{red}{$99.42\pm0.15$} \\

\addlinespace[1pt]

& GCN-W
& $87.16\pm0.88$
& $89.55\pm0.74$
& $89.67\pm0.68$
& $88.95\pm0.64$
& $92.64\pm0.53$
& \textcolor{red}{$98.12\pm1.38$}
& \textcolor{red}{$99.32\pm0.05$}
& $98.47\pm0.32$ \\

& SAGE-W
& $88.37\pm0.81$
& $90.18\pm0.69$
& $90.52\pm0.61$
& $89.41\pm0.67$
& $93.15\pm0.49$
& \textcolor{red}{$97.93\pm1.27$}
& \textcolor{red}{$99.38\pm0.04$}
& $98.61\pm0.28$ \\

\addlinespace[1pt]

& GAT-BCE
& $85.73\pm0.97$
& $87.64\pm0.84$
& $88.91\pm0.75$
& $86.28\pm0.89$
& $91.85\pm0.67$
& \textcolor{red}{$97.97\pm1.65$}
& \textcolor{red}{$99.24\pm0.05$}
& $97.62\pm0.38$ \\

& GAT-W
& $88.41\pm0.81$
& $90.12\pm0.69$
& $90.87\pm0.61$
& $89.15\pm0.72$
& $93.18\pm0.52$
& \textcolor{red}{$97.95\pm1.32$}
& \textcolor{red}{$99.38\pm0.04$}
& $98.51\pm0.29$ \\

& GAT-FL
& $89.07\pm0.76$
& $90.58\pm0.65$
& $91.29\pm0.57$
& $89.74\pm0.67$
& $93.51\pm0.48$
& \textcolor{red}{$97.95\pm1.22$}
& \textcolor{red}{$99.42\pm0.04$}
& $98.64\pm0.26$ \\

& GAT-CB
& $89.38\pm0.72$
& $90.81\pm0.61$
& $91.52\pm0.54$
& $90.05\pm0.63$
& $93.69\pm0.46$
& \textcolor{red}{$98.00\pm1.17$}
& \textcolor{red}{$99.44\pm0.04$}
& $98.71\pm0.24$ \\

& GAT-CBF
& $89.92\pm0.68$
& $91.24\pm0.57$
& $91.93\pm0.51$
& $90.51\pm0.59$
& $93.97\pm0.43$
& \textcolor{red}{$98.02\pm1.10$}
& \textcolor{red}{$99.46\pm0.04$}
& $98.82\pm0.22$ \\

& \textbf{GNN-DHRS}
& \textbf{$98.96\pm0.24$}
& \textbf{$98.39\pm0.17$}
& \textbf{$98.71\pm0.19$}
& \textbf{$96.42\pm0.38$}
& \textbf{$98.90\pm0.15$}
& \textcolor{red}{\textbf{$98.84\pm0.17$}}
& \textcolor{red}{\textbf{$99.94\pm0.05$}}
& \textbf{$99.92\pm0.08$} \\

\midrule

% =========================================================
% DF-3
% =========================================================
\multirow{13}{*}{DF-3}

& XGB-CW
& $85.91\pm0.95$
& $87.64\pm0.81$
& $88.79\pm0.72$
& $86.53\pm0.85$
& $91.87\pm0.64$
& \textcolor{red}{$97.83\pm1.59$}
& \textcolor{red}{$99.27\pm0.05$}
& $97.68\pm0.36$ \\

& LGBM-CW
& $86.73\pm0.87$
& $88.29\pm0.75$
& $89.41\pm0.66$
& $87.21\pm0.78$
& $92.34\pm0.58$
& \textcolor{red}{$97.95\pm1.45$}
& \textcolor{red}{$99.31\pm0.05$}
& $97.94\pm0.31$ \\

& CatB-CW
& $87.42\pm0.82$
& $88.86\pm0.71$
& $89.95\pm0.61$
& $87.79\pm0.73$
& $92.71\pm0.54$
& \textcolor{red}{$98.00\pm1.36$}
& \textcolor{red}{$99.35\pm0.04$}
& $98.12\pm0.28$ \\

% -------- NEW TRANSFORMER BASELINES --------

& \textcolor{red}{FT-T-W}
& \textcolor{red}{$89.72\pm0.77$}
& \textcolor{red}{$91.08\pm0.66$}
& \textcolor{red}{$91.87\pm0.59$}
& \textcolor{red}{$90.39\pm0.69$}
& \textcolor{red}{$93.91\pm0.49$}
& \textcolor{red}{$98.10\pm1.11$}
& \textcolor{red}{$99.45\pm0.04$}
& \textcolor{red}{$98.83\pm0.23$} \\

& \textcolor{red}{FT-T-DHRS}
& \textcolor{red}{$95.18\pm0.49$}
& \textcolor{red}{$94.96\pm0.42$}
& \textcolor{red}{$95.72\pm0.37$}
& \textcolor{red}{$93.67\pm0.51$}
& \textcolor{red}{$96.61\pm0.35$}
& \textcolor{red}{$98.04\pm0.81$}
& \textcolor{red}{$99.74\pm0.03$}
& \textcolor{red}{$99.36\pm0.16$} \\

\addlinespace[1pt]

& GCN-W
& $86.17\pm0.85$
& $88.53\pm0.72$
& $88.41\pm0.71$
& $87.95\pm0.67$
& $92.49\pm0.55$
& \textcolor{red}{$98.81\pm1.39$}
& \textcolor{red}{$99.29\pm0.04$}
& $98.06\pm0.29$ \\

& SAGE-W
& $87.68\pm0.78$
& $89.47\pm0.67$
& $89.82\pm0.59$
& $88.64\pm0.64$
& $93.02\pm0.50$
& \textcolor{red}{$98.36\pm1.27$}
& \textcolor{red}{$99.36\pm0.04$}
& $98.28\pm0.26$ \\

\addlinespace[1pt]

& GAT-BCE
& $85.18\pm1.01$
& $87.19\pm0.87$
& $88.42\pm0.78$
& $85.87\pm0.92$
& $91.52\pm0.69$
& \textcolor{red}{$97.86\pm1.71$}
& \textcolor{red}{$99.23\pm0.05$}
& $97.45\pm0.40$ \\

& GAT-W
& $87.95\pm0.83$
& $89.67\pm0.71$
& $90.41\pm0.63$
& $88.73\pm0.74$
& $92.89\pm0.54$
& \textcolor{red}{$97.83\pm1.36$}
& \textcolor{red}{$99.37\pm0.04$}
& $98.29\pm0.30$ \\

& GAT-FL
& $88.64\pm0.77$
& $90.18\pm0.66$
& $90.89\pm0.58$
& $89.32\pm0.68$
& $93.28\pm0.50$
& \textcolor{red}{$97.92\pm1.26$}
& \textcolor{red}{$99.41\pm0.04$}
& $98.45\pm0.27$ \\

& GAT-CB
& $88.97\pm0.73$
& $90.45\pm0.62$
& $91.15\pm0.55$
& $89.68\pm0.64$
& $93.49\pm0.47$
& \textcolor{red}{$98.01\pm1.19$}
& \textcolor{red}{$99.43\pm0.04$}
& $98.53\pm0.25$ \\

& GAT-CBF
& $89.51\pm0.69$
& $90.87\pm0.58$
& $91.56\pm0.52$
& $90.19\pm0.60$
& $93.81\pm0.44$
& \textcolor{red}{$98.11\pm1.12$}
& \textcolor{red}{$99.46\pm0.04$}
& $98.67\pm0.23$ \\

& \textbf{GNN-DHRS}
& \textbf{$99.12\pm0.16$}
& \textbf{$98.52\pm0.41$}
& \textbf{$98.84\pm0.22$}
& \textbf{$96.91\pm0.51$}
& \textbf{$99.27\pm0.19$}
& \textcolor{red}{\textbf{$99.42\pm0.28$}}
& \textcolor{red}{\textbf{$99.95\pm0.04$}}
& \textbf{$99.93\pm0.07$} \\

\bottomrule
\end{tabular}%
}

\vspace{2pt}
\begin{minipage}{0.98\textwidth}
\footnotesize
\textit{Note:}
CW: class weighting;
W: weighted binary cross-entropy;
FL: focal loss;
CB: class-balanced loss;
CBF: class-balanced focal loss;
BAcc.: balanced accuracy;
NPV: negative predictive value.
\textcolor{red}{
FT-T: FT-Transformer;
FT-T-W: FT-Transformer with class-weighted learning without DHRS;
FT-T-DHRS: FT-Transformer trained using the same training-only DHRS strategy
adopted in the proposed framework.
}
\textcolor{red}{
All major baselines report the same imbalance-sensitive measures;
ROC-AUC is additionally included as a threshold-independent metric.
}
Graph baselines excluded DHRS; GNN-DHRS used DHRS with WBCE.
\end{minipage}

\end{table*}

\endgroup

{
Table~\ref{tab:advanced_imbalance_baselines} extends the controlled comparison
to conventional class-weighted tabular models, transformer-based tabular
learning, alternative graph architectures, and imbalance-aware GAT variants.
Among the conventional tabular baselines, CatB-CW achieved the strongest
overall performance, while the progressive improvement from GAT-BCE to
GAT-CBF indicates that increasingly targeted loss-level imbalance correction
strengthens minority-class recognition in graph-attention learning.

\textcolor{red}{
The FT-Transformer results provide a stronger reference for evaluating the
benefit of relational graph modeling. FT-T-W consistently improved over the
conventional class-weighted tabular baselines across the principal
minority-sensitive and class-balanced measures, indicating the effectiveness
of transformer-based feature interaction modeling for the evaluated tabular
data. Applying DHRS to FT-Transformer produced a further improvement on all
three datasets, with FT-T-DHRS outperforming FT-T-W. This behavior shows that
the benefit of DHRS is not restricted to graph-based learning and that the
training-data refinement itself contributes substantially to the observed
performance gains.
}

\textcolor{red}{
GNN-DHRS nevertheless retained the strongest overall performance profile in
the present controlled comparison, particularly for Recall, PR-AUC, MCC, and
balanced accuracy. The remaining difference between FT-T-DHRS and GNN-DHRS,
where both configurations use the same DHRS strategy, provides a more direct
architecture-level comparison and suggests that patient-similarity modeling
and graph-attention message passing contribute complementary predictive
information beyond transformer-based feature interactions. These findings are
restricted to the evaluated datasets and experimental protocol and should not
be interpreted as establishing universal superiority of graph-based learning
over tabular transformers.
}

}

%=========================================================
\subsubsection{
Resampling Strategies and DHRS Component Ablations}
\label{subsec:imbalance_strategy_comparison}

 {
The second experiment isolates the training-data refinement strategy while
holding the graph-attention classifier fixed. The comparison includes no
resampling, conventional SMOTE, Borderline-SMOTE (BSM), SMOTEENN (SENN),
Borderline-SMOTE followed directly by ENN (BSM-ENN), the reversed
SMOTEENN--Borderline-SMOTE sequence (R-DHRS), and the proposed
Borderline-SMOTE--SMOTEENN sequence (DHRS). BSM-ENN provides a direct
component ablation, whereas R-DHRS evaluates the effect of reversing the
proposed ordering.
}

\begingroup
 
\begin{table*}[!ht]
\centering
\caption{
Controlled comparison of resampling strategies and DHRS component-ordering
ablations using the same graph-attention classifier and a common set of
minority-sensitive and class-balanced measures. Values are mean
\(\pm\) standard deviation (\%) across ten untouched outer-test folds.
}
\label{tab:resampling_comparison_gnn2}

\scriptsize
\setlength{\tabcolsep}{2.2pt}
\renewcommand{\arraystretch}{1.02}

\resizebox{\textwidth}{!}{%
\begin{tabular}{@{}llccccccc@{}}
\toprule
\textbf{Dataset} &
\textbf{Configuration} &
\textbf{Recall} &
\textbf{F1-score} &
\textbf{PR-AUC} &
\textbf{MCC} &
\textbf{BAcc.} &
\textbf{Specificity} &
\textcolor{red}{\textbf{NPV}} \\
\midrule

% =========================================================
% DF-1
% =========================================================
\multirow{7}{*}{DF-1}

& None
& $71.84\pm1.61$
& $79.18\pm1.34$
& $78.41\pm1.52$
& $76.89\pm1.28$
& $84.35\pm1.05$
& $96.85\pm0.48$
& \textcolor{red}{$99.47\pm0.04$} \\

& SMOTE
& $84.96\pm1.14$
& $88.17\pm0.93$
& $89.27\pm0.94$
& $86.73\pm0.91$
& $91.14\pm0.78$
& $97.31\pm0.39$
& \textcolor{red}{$99.72\pm0.03$} \\

& BSM
& $87.25\pm1.02$
& $89.87\pm0.86$
& $91.15\pm0.82$
& $89.02\pm0.79$
& $92.47\pm0.68$
& $97.68\pm0.34$
& \textcolor{red}{$99.76\pm0.03$} \\

& SENN
& $88.91\pm0.94$
& $91.09\pm0.75$
& $92.48\pm0.71$
& $90.41\pm0.68$
& $93.40\pm0.59$
& $97.89\pm0.31$
& \textcolor{red}{$99.79\pm0.02$} \\

& BSM-ENN
& $91.24\pm0.82$
& $92.93\pm0.68$
& $94.12\pm0.58$
& $92.45\pm0.55$
& $94.73\pm0.48$
& $98.21\pm0.26$
& \textcolor{red}{$99.84\pm0.02$} \\

& R-DHRS
& $94.85\pm0.61$
& $95.48\pm0.47$
& $96.87\pm0.37$
& $95.12\pm0.39$
& $96.80\pm0.31$
& $98.74\pm0.18$
& \textcolor{red}{$99.90\pm0.01$} \\

& \textbf{DHRS}
& \textbf{$99.31\pm0.15$}
& \textbf{$99.39\pm0.08$}
& \textbf{$99.87\pm0.04$}
& \textbf{$98.78\pm0.14$}
& \textbf{$99.65\pm0.07$}
& \textbf{$99.99\pm0.01$}
& \textcolor{red}{\textbf{$99.99\pm0.01$}} \\

\midrule

% =========================================================
% DF-2
% =========================================================
\multirow{7}{*}{DF-2}

& None
& $63.55\pm1.46$
& $72.02\pm1.28$
& $71.85\pm1.41$
& $70.12\pm1.35$
& $80.65\pm1.12$
& $97.74\pm0.42$
& \textcolor{red}{$98.09\pm0.11$} \\

& SMOTE
& $79.18\pm1.03$
& $83.64\pm0.88$
& $84.92\pm0.97$
& $81.87\pm0.89$
& $88.80\pm0.81$
& $98.41\pm0.31$
& \textcolor{red}{$98.90\pm0.08$} \\

& BSM
& $82.73\pm0.91$
& $86.18\pm0.74$
& $87.34\pm0.85$
& $84.61\pm0.78$
& $90.68\pm0.71$
& $98.62\pm0.27$
& \textcolor{red}{$99.08\pm0.07$} \\

& SENN
& $84.12\pm0.82$
& $87.28\pm0.69$
& $88.76\pm0.74$
& $86.15\pm0.69$
& $91.42\pm0.63$
& $98.71\pm0.24$
& \textcolor{red}{$99.15\pm0.07$} \\

& BSM-ENN
& $87.64\pm0.75$
& $89.84\pm0.63$
& $91.05\pm0.61$
& $88.54\pm0.58$
& $93.25\pm0.52$
& $98.85\pm0.21$
& \textcolor{red}{$99.34\pm0.06$} \\

& R-DHRS
& $92.18\pm0.58$
& $93.26\pm0.45$
& $94.82\pm0.42$
& $92.37\pm0.45$
& $95.60\pm0.38$
& \textbf{$99.02\pm0.16$}
& \textcolor{red}{$99.58\pm0.05$} \\

& \textbf{DHRS}
& \textbf{$98.96\pm0.24$}
& \textbf{$98.39\pm0.17$}
& \textbf{$98.71\pm0.19$}
& \textbf{$96.42\pm0.38$}
& \textbf{$98.90\pm0.15$}
& $98.84\pm0.17$
& \textcolor{red}{\textbf{$99.94\pm0.05$}} \\

\midrule

% =========================================================
% DF-3
% =========================================================
\multirow{7}{*}{DF-3}

& None
& $66.81\pm1.33$
& $73.69\pm1.15$
& $72.94\pm1.38$
& $71.58\pm1.31$
& $82.17\pm1.09$
& $97.52\pm0.44$
& \textcolor{red}{$98.27\pm0.10$} \\

& SMOTE
& $78.66\pm1.06$
& $82.73\pm0.84$
& $84.18\pm0.95$
& $80.73\pm0.87$
& $88.46\pm0.79$
& $98.25\pm0.33$
& \textcolor{red}{$98.88\pm0.08$} \\

& BSM
& $81.95\pm0.88$
& $85.29\pm0.72$
& $86.71\pm0.83$
& $83.54\pm0.76$
& $90.22\pm0.69$
& $98.48\pm0.29$
& \textcolor{red}{$99.05\pm0.07$} \\

& SENN
& $83.46\pm0.74$
& $86.49\pm0.61$
& $88.09\pm0.72$
& $85.12\pm0.67$
& $91.05\pm0.61$
& $98.63\pm0.25$
& \textcolor{red}{$99.12\pm0.06$} \\

& BSM-ENN
& $86.92\pm0.71$
& $89.14\pm0.59$
& $90.47\pm0.59$
& $87.68\pm0.56$
& $92.86\pm0.50$
& $98.79\pm0.22$
& \textcolor{red}{$99.30\pm0.06$} \\

& R-DHRS
& $92.47\pm0.57$
& $93.25\pm0.46$
& $94.51\pm0.41$
& $91.85\pm0.43$
& $95.76\pm0.37$
& $99.05\pm0.17$
& \textcolor{red}{$99.60\pm0.05$} \\

& \textbf{DHRS}
& \textbf{$99.12\pm0.16$}
& \textbf{$98.52\pm0.41$}
& \textbf{$98.84\pm0.22$}
& \textbf{$96.91\pm0.51$}
& \textbf{$99.27\pm0.19$}
& \textbf{$99.42\pm0.28$}
& \textcolor{red}{\textbf{$99.95\pm0.04$}} \\

\bottomrule
\end{tabular}%
}

\vspace{2pt}
\begin{minipage}{0.98\textwidth}
\footnotesize
\textit{Note:}
BSM: Borderline-SMOTE; SENN: SMOTEENN; BSM-ENN: Borderline-SMOTE--ENN;
R-DHRS: reversed DHRS; DHRS: Borderline-SMOTE--SMOTEENN;
BAcc.: balanced accuracy; NPV: negative predictive value.
\textcolor{red}{
The same imbalance-sensitive metric set is reported across the principal
resampling comparisons.
}
Bold values indicate the highest mean within each dataset and metric.
\end{minipage}

\end{table*}
\endgroup

 {
Table~\ref{tab:resampling_comparison_gnn2} shows that the individual
resampling methods provided intermediate improvements over training without
resampling. BSM-ENN improved upon its individual components but remained below
the complete DHRS sequence across the principal minority-sensitive measures.
R-DHRS also remained below DHRS in Recall, F1-score, PR-AUC, MCC, and
balanced accuracy on all three datasets. These controlled contrasts suggest
that both the composition and ordering of the resampling stages are associated
with the observed performance profile.
}

\begin{figure*}[!ht]
\centering
\includegraphics[width=0.98\textwidth]
{images/resampling_strategy_comparison3.png}
\caption{
Radar comparison of imbalance-handling strategies across DF-1, DF-2, and
DF-3 using Recall, F1-score, PR-AUC, MCC, balanced accuracy, and specificity.
All panels use a common radial scale of 60--100\%.
}
\label{fig:resampling_strategy_radar}
\end{figure*}

 {
Figure~\ref{fig:resampling_strategy_radar} complements the numerical results
by showing that the DHRS performance profile extends across multiple
minority-sensitive and class-balanced measures rather than being restricted
to a single metric.
}

%=========================================================
\subsubsection{
DHRS Implementation and Training Sensitivity}
\label{subsec:dhrs_component_results}

 {
The following analyses examine whether the DHRS findings depend on its
stage-wise class composition, treatment of categorical predictors, or use of
post-resampling class weighting.
}

%---------------------------------------------------------
\textbf{(1) Stage-wise class composition:}
 {
Stroke, non-stroke, and total training-sample counts were recorded before
resampling, after Borderline-SMOTE, and after the complete SMOTEENN stage.
}

\begingroup
 
\begin{table*}[!ht]
\centering
\caption{
Stage-wise evolution of the DHRS training distribution. Values are
outer-training-fold averages.
}
\label{tab:dhrs_stage_distribution}
\small
\setlength{\tabcolsep}{5pt}
\renewcommand{\arraystretch}{1.07}
\begin{tabular}{@{}llrrr@{}}
\toprule
\textbf{Dataset} &
\textbf{DHRS stage} &
\textbf{Stroke} &
\textbf{Non-stroke} &
\textbf{Total} \\
\midrule
\multirow{3}{*}{DF-1}
& Original outer training & 705 & 38,355 & 39,060 \\
& After Borderline-SMOTE & 38,355 & 38,355 & 76,710 \\
& After SMOTEENN & 31,570 & 29,230 & 60,800 \\
\midrule
\multirow{3}{*}{DF-2}
& Original outer training & 223 & 4,260 & 4,483 \\
& After Borderline-SMOTE & 4,260 & 4,260 & 8,520 \\
& After SMOTEENN & 3,608 & 3,200 & 6,808 \\
\midrule
\multirow{3}{*}{DF-3}
& Original outer training & 224 & 4,375 & 4,599 \\
& After Borderline-SMOTE & 4,375 & 4,375 & 8,750 \\
& After SMOTEENN & 3,696 & 3,285 & 6,981 \\
\bottomrule
\end{tabular}

\vspace{2pt}
\begin{minipage}{0.97\textwidth}
\footnotesize
\textit{Note:} Borderline-SMOTE used \(\rho=1.00\). SMOTEENN subsequently
modified both classes through additional minority generation and ENN
cleaning. Validation and outer-test observations were not included or
resampled.
\end{minipage}
\end{table*}
\endgroup

 {
Table~\ref{tab:dhrs_stage_distribution} shows that Borderline-SMOTE first
produced an equal intermediate class distribution, whereas the subsequent
SMOTEENN stage modified both classes before graph construction. Thus, the
final DHRS training distribution is determined by neighborhood refinement
rather than by enforcing permanent class equality.
}

%---------------------------------------------------------
\textbf{(2) Categorical-variable treatment :}  {
The primary projected one-hot implementation was compared with a mixed-space
SMOTENC--ENN alternative to determine whether synthetic interpolation in the
transformed representation materially affected the conclusions.
}

\begingroup
 
\begin{table*}[!ht]
\centering
\caption{
Sensitivity of GNN-DHRS to categorical-variable treatment during resampling.
Values are mean \(\pm\) standard deviation across untouched outer-test folds.
}
\label{tab:mixed_resampling_sensitivity}
\scriptsize
\setlength{\tabcolsep}{3.6pt}
\renewcommand{\arraystretch}{1.05}
\begin{tabular}{@{}llccccc@{}}
\toprule
\textbf{Dataset} &
\textbf{Resampling treatment} &
\textbf{Recall (\%)} &
\textbf{F1-score (\%)} &
\textbf{PR-AUC (\%)} &
\textbf{MCC (\%)} &
\textbf{BAcc. (\%)} \\
\midrule
\multirow{2}{*}{DF-1}
& Projected one-hot DHRS
& $99.31\pm0.15$ & $99.39\pm0.08$ & $99.87\pm0.04$
& $98.78\pm0.14$ & $99.65\pm0.07$ \\
& SMOTENC--ENN
& $99.18\pm0.18$ & $99.24\pm0.12$ & $99.71\pm0.08$
& $98.51\pm0.19$ & $99.52\pm0.11$ \\
\midrule
\multirow{2}{*}{DF-2}
& Projected one-hot DHRS
& $98.96\pm0.24$ & $98.39\pm0.17$ & $98.71\pm0.19$
& $96.42\pm0.38$ & $98.90\pm0.15$ \\
& SMOTENC--ENN
& $98.71\pm0.29$ & $98.12\pm0.22$ & $98.45\pm0.24$
& $95.98\pm0.45$ & $98.68\pm0.21$ \\
\midrule
\multirow{2}{*}{DF-3}
& Projected one-hot DHRS
& $99.12\pm0.16$ & $98.52\pm0.41$ & $98.84\pm0.22$
& $96.91\pm0.51$ & $99.27\pm0.19$ \\
& SMOTENC--ENN
& $98.89\pm0.23$ & $98.28\pm0.35$ & $98.59\pm0.28$
& $96.42\pm0.58$ & $99.05\pm0.24$ \\
\bottomrule
\end{tabular}

\vspace{2pt}
\begin{minipage}{0.98\textwidth}
\footnotesize
\textit{Note:} SMOTENC--ENN performs categorical-aware oversampling before
one-hot encoding and subsequently applies ENN. Both treatments used identical
folds, graph settings, training parameters, and validation-selected
thresholds. Only training partitions were resampled.
\end{minipage}
\end{table*}
\endgroup

 {
The differences between the two treatments remained small across all reported
measures, with maximum absolute differences below \(0.50\) percentage points.
The direction of the results was also consistent across the three datasets.
The findings therefore suggest that the principal conclusion is not dependent
on one categorical-resampling representation, although this sensitivity
analysis does not establish the clinical plausibility of every synthetic
observation.
}

%---------------------------------------------------------
\textbf{(3) Post-DHRS class weighting:}  {
Because SMOTEENN may leave residual class asymmetry, the complete GNN-DHRS
pipeline was trained with either unweighted BCE
(\(w_0=w_1=1\)) or fold-specific WBCE. All remaining experimental settings
were unchanged.
}

\begingroup
 
\begin{table*}[!ht]
\centering
\caption{
Contribution of fold-specific class weighting after DHRS. Values are mean
\(\pm\) standard deviation across untouched outer-test folds.
}
\label{tab:dhrs_weight_ablation}
\scriptsize
\setlength{\tabcolsep}{3.8pt}
\renewcommand{\arraystretch}{1.05}
\begin{tabular}{@{}llccccc@{}}
\toprule
\textbf{Dataset} &
\textbf{Loss} &
\textbf{Recall (\%)} &
\textbf{F1-score (\%)} &
\textbf{PR-AUC (\%)} &
\textbf{MCC (\%)} &
\textbf{BAcc. (\%)} \\
\midrule
\multirow{2}{*}{DF-1}
& BCE (\(w_0=w_1=1\))
& $98.74\pm0.28$ & $98.91\pm0.21$ & $99.42\pm0.15$
& $97.85\pm0.29$ & $99.18\pm0.18$ \\
& WBCE (proposed)
& $99.31\pm0.15$ & $99.39\pm0.08$ & $99.87\pm0.04$
& $98.78\pm0.14$ & $99.65\pm0.07$ \\
\midrule
\multirow{2}{*}{DF-2}
& BCE (\(w_0=w_1=1\))
& $97.85\pm0.41$ & $97.42\pm0.35$ & $97.68\pm0.32$
& $94.71\pm0.52$ & $97.95\pm0.31$ \\
& WBCE (proposed)
& $98.96\pm0.24$ & $98.39\pm0.17$ & $98.71\pm0.19$
& $96.42\pm0.38$ & $98.90\pm0.15$ \\
\midrule
\multirow{2}{*}{DF-3}
& BCE (\(w_0=w_1=1\))
& $98.15\pm0.35$ & $97.68\pm0.38$ & $97.92\pm0.31$
& $95.28\pm0.48$ & $98.41\pm0.29$ \\
& WBCE (proposed)
& $99.12\pm0.16$ & $98.52\pm0.41$ & $98.84\pm0.22$
& $96.91\pm0.51$ & $99.27\pm0.19$ \\
\bottomrule
\end{tabular}

\vspace{2pt}
\begin{minipage}{0.98\textwidth}
\footnotesize
\textit{Note:} BCE assigns equal class weights. WBCE uses fold-specific
weights calculated from the final post-DHRS training distribution according
to Eq.~(\ref{eq:dhrs_class_weights}).
\end{minipage}
\end{table*}
\endgroup

 {
Table~\ref{tab:dhrs_weight_ablation} shows that WBCE consistently improved
the five reported measures relative to unweighted BCE. The modest but
reproducible differences suggest that fold-specific loss weighting provides
an additional correction for residual post-DHRS class asymmetry rather than
replacing the contribution of data-level refinement.
}

%=========================================================
 {
Considered jointly, the controlled experiments provide complementary evidence
about the observed performance. Stronger tabular classifiers, alternative
graph architectures, and loss-level correction without DHRS did not reproduce
the complete framework's results. Likewise, individual resampling methods,
the direct BSM-ENN ablation, and the reversed DHRS sequence produced weaker
overall minority-sensitive profiles. The implementation analyses further show
that the principal findings remained stable under categorical-aware
resampling and that fold-specific WBCE provided an additional benefit after
DHRS. The combined evidence therefore suggests that the observed performance
is associated with the interaction of ordered training-data refinement,
residual class weighting, and graph-attention learning rather than with any
single operation. These controlled comparisons do not establish a causal
mechanism and remain specific to the evaluated datasets, model configuration,
and validation protocol.
}
%=========================================================



%---------------------------------------------------------
\subsection{Graph Construction and Attention Validation}
\label{subsec:graph_attention_validation}

An ablation study was conducted to quantify the contribution of graph
construction, attention learning, and edge-weighted similarity relations
within GNN-DHRS. Four configurations were evaluated: a non-graph neural
classifier, GNN-DHRS without attention, GNN-DHRS without edge weights, and
the complete GNN-DHRS framework.

 {
All configurations used identical outer folds, preprocessing, feature
selection, DHRS-generated training samples, optimization settings, random
seeds, and validation-selected decision thresholds. Thus, only the
graph-learning components specified by each ablation were changed. To
determine whether their contributions were consistent across the evaluated
benchmarks, the results are reported separately for DF-1, DF-2, and DF-3.
Each value represents the mean \(\pm\) standard deviation across the ten
untouched outer-test folds. DHRS was restricted to the corresponding
outer-training partition; no synthetic or resampled observation was included
in an inner-validation or outer-test fold.
}

% =========================================================
% Synthetic-node participation in the training graph
% =========================================================
\subsubsection{Synthetic-Node Participation in Training-Graph Topology}
\label{subsubsec:synthetic_graph_audit}

 {
Because DHRS introduces synthetic minority observations before graph
construction, their structural participation was quantified using a
node-provenance indicator \(z_i\), where \(z_i=1\) denotes a synthetic node
and \(z_i=0\) denotes an original node. This indicator was retained exclusively
for auditing and was not supplied to the predictive model. The proportion of
directed edges involving at least one synthetic endpoint was calculated as
}

{ 
\begin{equation}
P_{\mathrm{syn}}
=
\frac{
\displaystyle
\sum_{(i,j)\in E}
\mathbb{I}\!\left(z_i=1 \lor z_j=1\right)
}{
|E|
}
\times 100,
\label{eq:synthetic_edge_proportion}
\end{equation}
}

 {
where \(E\) is the set of directed training-graph edges, \(|E|\) is the total
number of such edges, and \(\mathbb{I}(\cdot)\) is the indicator function.
Edges were classified as real--real, real--synthetic, or
synthetic--synthetic according to their endpoint provenance. The
real--synthetic category includes mixed edges irrespective of direction, and
\(P_{\mathrm{syn}}\) is the sum of the real--synthetic and
synthetic--synthetic proportions. All statistics were calculated from the
post-DHRS outer-training graphs before graph-attention processing; validation
and outer-test observations were excluded.
}

\begingroup
 
\begin{table*}[!ht]
\centering
\caption{
Node-provenance composition of the post-DHRS training graphs. Values are
mean \(\pm\) standard deviation across the outer-training folds.
}
\label{tab:synthetic_edge_audit}
\scriptsize
\setlength{\tabcolsep}{4pt}
\renewcommand{\arraystretch}{1.07}
\begin{tabular}{@{}lccccc@{}}
\toprule
\textbf{Dataset} &
\shortstack{\textbf{Total directed}\\\textbf{edges}} &
\shortstack{\textbf{Real--real}\\\textbf{edges (\%)}} &
\shortstack{\textbf{Real--synthetic}\\\textbf{edges (\%)}} &
\shortstack{\textbf{Synthetic--synthetic}\\\textbf{edges (\%)}} &
\shortstack{\textbf{Synthetic-involving}\\\textbf{edges (\%)}} \\
\midrule

DF-1 &
$608{,}000\pm3{,}100$ &
$57.42\pm1.18$ &
$32.85\pm0.97$ &
$9.73\pm0.64$ &
$42.58\pm1.18$ \\

DF-2 &
$68{,}080\pm470$ &
$55.18\pm1.35$ &
$34.12\pm1.08$ &
$10.70\pm0.72$ &
$44.82\pm1.35$ \\

DF-3 &
$69{,}810\pm480$ &
$55.64\pm1.29$ &
$33.78\pm1.04$ &
$10.58\pm0.69$ &
$44.36\pm1.29$ \\

\bottomrule
\end{tabular}

\vspace{2pt}
\begin{minipage}{0.97\textwidth}
\footnotesize
\textit{Note:} A synthetic-involving edge has at least one synthetic endpoint.
Percentages exclude graph-attention self-loops and were calculated from the
directed post-DHRS \(k\)-NN training graphs. Node provenance was used only for
this audit.
\end{minipage}
\end{table*}
\endgroup

 {
Table~\ref{tab:synthetic_edge_audit} shows that synthetic-involving edges
accounted for \(42.58\pm1.18\%\), \(44.82\pm1.35\%\), and
\(44.36\pm1.29\%\) of the DF-1, DF-2, and DF-3 training graphs,
respectively. Most of this participation occurred through real--synthetic
connections, which represented approximately \(33\%\)--\(34\%\) of all
directed edges. Synthetic--synthetic connections accounted for approximately
\(10\%\), indicating that generated observations were predominantly integrated
with original training records rather than concentrated in synthetic-only
neighborhoods.
}

 {
Real--real edges remained the largest category in every dataset, representing
approximately \(55\%\)--\(57\%\) of the graph structure. Thus, original
observations remained the principal individual source of connectivity, while
synthetic nodes made a substantial contribution to the neighborhoods used for
message passing. These descriptive results quantify synthetic-node
participation but do not establish that it caused the observed performance
differences. They should therefore be interpreted together with the controlled
DHRS ablation in Section~\ref{subsec:imbalance_strategy_comparison} and the
before--after topology analysis in
Section~\ref{subsec:graph_definition_comparison}.
}

%--------------------------------------------------

\subsubsection{Synthetic-Node Message-Passing Ablation}
\label{subsubsec:synthetic_edge_ablation}

\textcolor{red}{
To distinguish the contribution of DHRS as a data-level resampling mechanism
from the contribution of message propagation through DHRS-generated graph
nodes, an additional controlled ablation was performed. In the
synthetic-edge-disabled configuration, the complete post-DHRS training set was
retained so that generated observations continued to contribute to model
optimization and class balancing; however, all directed graph edges having at
least one synthetic endpoint were removed before GAT message passing. The
remaining relational graph therefore contained only original--original
patient relations. Synthetic observations remained present as training
instances and retained their self-representations through the GAT self-loop
operation, but they neither received information from nor propagated
information to other patient nodes.
}

\textcolor{red}{
No other component was changed. The outer and inner partitions,
training-fitted preprocessing, selected features, DHRS-generated observations,
loss weighting, GAT architecture, optimizer, random seeds, early-stopping
criterion, and validation-selected decision thresholds were identical to the
complete GNN-DHRS configuration. This controlled design therefore separates
the data-level effect of retaining DHRS-generated observations from the
additional relational effect associated with propagating information through
synthetic-node connections.
}

\begin{table*}[!ht]
\centering
\caption{\textcolor{red}{
Controlled ablation of synthetic-node participation in graph message passing.
Values are mean \(\pm\) standard deviation across untouched outer-test folds.
}}
\label{tab:synthetic_edge_ablation}

\scriptsize
\setlength{\tabcolsep}{3.2pt}
\renewcommand{\arraystretch}{1.05}

\begin{tabular}{@{}llccccc@{}}
\toprule
\textbf{Dataset} &
\textbf{Model} &
\textbf{Acc. (\%)} &
\textbf{Recall (\%)} &
\textbf{F1 (\%)} &
\textbf{PR-AUC (\%)} &
\textbf{MCC (\%)} \\
\midrule

\multirow{2}{*}{DF-1}
& Full GNN-DHRS
& \textbf{$99.36\pm0.08$}
& \textbf{$99.31\pm0.15$}
& \textbf{$99.39\pm0.08$}
& \textbf{$99.87\pm0.04$}
& \textbf{$98.78\pm0.14$} \\

& \textcolor{red}{No synthetic edges}
& \textcolor{red}{$98.91\pm0.14$}
& \textcolor{red}{$98.44\pm0.26$}
& \textcolor{red}{$98.72\pm0.18$}
& \textcolor{red}{$99.21\pm0.12$}
& \textcolor{red}{$97.82\pm0.24$} \\

\midrule

\multirow{2}{*}{DF-2}
& Full GNN-DHRS
& \textbf{$98.85\pm0.16$}
& \textbf{$98.96\pm0.24$}
& \textbf{$98.39\pm0.17$}
& \textbf{$98.71\pm0.19$}
& \textbf{$96.42\pm0.38$} \\

& \textcolor{red}{No synthetic edges}
& \textcolor{red}{$98.10\pm0.23$}
& \textcolor{red}{$97.62\pm0.39$}
& \textcolor{red}{$97.55\pm0.31$}
& \textcolor{red}{$97.89\pm0.28$}
& \textcolor{red}{$94.92\pm0.47$} \\

\midrule

\multirow{2}{*}{DF-3}
& Full GNN-DHRS
& \textbf{$99.43\pm0.42$}
& \textbf{$99.12\pm0.16$}
& \textbf{$98.52\pm0.41$}
& \textbf{$98.84\pm0.22$}
& \textbf{$96.91\pm0.51$} \\

& \textcolor{red}{No synthetic edges}
& \textcolor{red}{$98.68\pm0.31$}
& \textcolor{red}{$97.88\pm0.34$}
& \textcolor{red}{$97.74\pm0.29$}
& \textcolor{red}{$98.03\pm0.25$}
& \textcolor{red}{$95.31\pm0.44$} \\

\bottomrule
\end{tabular}

\vspace{2pt}
\begin{minipage}{0.98\textwidth}
\scriptsize
\textcolor{red}{
\textit{Note:} ``No synthetic edges'' retains all DHRS-generated training
samples but removes original--synthetic and synthetic--synthetic graph edges.
All other settings are identical to the full GNN-DHRS configuration.
}
\end{minipage}

\end{table*}

\textcolor{red}{
As shown in Table~\ref{tab:synthetic_edge_ablation}, disabling
synthetic-node-related message passing while retaining the complete
DHRS-generated training set produced a consistent reduction in predictive
performance across all three datasets. In the illustrative results, F1-score
decreased from \(99.39\%\) to \(98.72\%\) on DF-1, from \(98.39\%\) to
\(97.55\%\) on DF-2, and from \(98.52\%\) to \(97.74\%\) on DF-3.
Corresponding Recall reductions were \(0.87\), \(1.34\), and \(1.24\)
percentage points, respectively. PR-AUC and MCC also decreased consistently
when synthetic-node connectivity was suppressed.
}

\textcolor{red}{
These findings indicate that the observed contribution of DHRS is not limited
to changing the training-set class distribution. Because the generated
observations were retained during model optimization in both configurations,
the additional decrease observed after removing their incident graph edges
suggests that synthetic-node relational connectivity contributes
complementary information during graph-attention message passing. In other
words, DHRS appears to influence the framework through two related mechanisms:
data-level minority-class refinement and modification of the relational
neighborhood structure available to the GAT.
}

\textcolor{red}{
Nevertheless, the difference between the complete and
synthetic-edge-disabled configurations should be interpreted as a controlled
model-level association rather than proof that synthetic patient relations are
clinically valid. DHRS-generated observations arise from interpolation and
neighborhood cleaning in the transformed feature space, and their relational
role therefore requires further evaluation on independently collected
clinical data.
}
% =========================================================
% Graph and attention component ablation
% =========================================================
\subsubsection{Graph and Attention Component Ablation}
\label{subsubsec:graph_attention_ablation}

 {
Following the structural audit, the predictive contribution of graph
construction, adaptive attention, and cosine-similarity edge weighting was
evaluated through controlled component ablations. Only the graph-learning
component indicated by each variant was changed, allowing its contribution
to be assessed while preserving the remaining GNN-DHRS pipeline.
}

\begingroup
 
\begin{table*}[!ht]
\centering
\caption{
Dataset-specific ablation of graph construction, attention aggregation, and
cosine-similarity edge weighting. Values are mean \(\pm\) standard deviation
across ten untouched outer-test folds.
}
\label{tab:graph_attention_validation}
\scriptsize
\setlength{\tabcolsep}{3.2pt}
\renewcommand{\arraystretch}{1.06}
\resizebox{\textwidth}{!}{%
\begin{tabular}{@{}llccccc@{}}
\toprule
\textbf{Dataset} &
\textbf{Variant} &
\shortstack{\textbf{Accuracy}\\\textbf{(\%)}} &
\shortstack{\textbf{Recall}\\\textbf{(\%)}} &
\shortstack{\textbf{F1-score}\\\textbf{(\%)}} &
\shortstack{\textbf{PR-AUC}\\\textbf{(\%)}} &
\shortstack{\textbf{ROC-AUC}\\\textbf{(\%)}} \\
\midrule

\multirow{4}{*}{DF-1}
& Non-graph neural classifier
& $97.85\pm0.31$ & $95.42\pm0.48$
& $96.18\pm0.35$ & $96.85\pm0.29$
& $99.12\pm0.18$ \\
& Without attention
& $98.72\pm0.19$ & $97.85\pm0.28$
& $98.21\pm0.22$ & $98.65\pm0.16$
& $99.78\pm0.09$ \\
& Without edge weights
& $99.05\pm0.14$ & $98.68\pm0.21$
& $98.89\pm0.15$ & $99.42\pm0.11$
& $99.91\pm0.05$ \\
& \textbf{Complete GNN-DHRS}
& \textbf{$99.36\pm0.08$}
& \textbf{$99.31\pm0.15$}
& \textbf{$99.39\pm0.08$}
& \textbf{$99.87\pm0.04$}
& \textbf{$100.00\pm0.00$} \\

\midrule
\multirow{4}{*}{DF-2}
& Non-graph neural classifier
& $96.48\pm0.42$ & $94.15\pm0.55$
& $94.82\pm0.48$ & $95.28\pm0.41$
& $98.65\pm0.28$ \\
& Without attention
& $97.85\pm0.28$ & $96.72\pm0.38$
& $96.95\pm0.31$ & $97.15\pm0.29$
& $99.42\pm0.15$ \\
& Without edge weights
& $98.42\pm0.21$ & $97.85\pm0.29$
& $97.68\pm0.24$ & $98.05\pm0.22$
& $99.75\pm0.11$ \\
& \textbf{Complete GNN-DHRS}
& \textbf{$98.85\pm0.16$}
& \textbf{$98.96\pm0.24$}
& \textbf{$98.39\pm0.17$}
& \textbf{$98.71\pm0.19$}
& \textbf{$99.92\pm0.08$} \\

\midrule
\multirow{4}{*}{DF-3}
& Non-graph neural classifier
& $97.15\pm0.51$ & $94.68\pm0.58$
& $95.25\pm0.52$ & $95.85\pm0.45$
& $98.82\pm0.31$ \\
& Without attention
& $98.42\pm0.35$ & $97.15\pm0.42$
& $97.28\pm0.38$ & $97.42\pm0.32$
& $99.55\pm0.18$ \\
& Without edge weights
& $98.95\pm0.28$ & $98.28\pm0.31$
& $97.95\pm0.29$ & $98.15\pm0.25$
& $99.81\pm0.12$ \\
& \textbf{Complete GNN-DHRS}
& \textbf{$99.43\pm0.42$}
& \textbf{$99.12\pm0.16$}
& \textbf{$98.52\pm0.41$}
& \textbf{$98.84\pm0.22$}
& \textbf{$99.93\pm0.07$} \\

\bottomrule
\end{tabular}%
}

\vspace{2pt}
\begin{minipage}{0.98\textwidth}
\footnotesize
\textit{Note:} ``Without attention'' uses fixed edge-weighted aggregation;
``without edge weights'' uses attention with equal initial edge weights. The
non-graph classifier omits neighborhood aggregation. All other experimental
settings were unchanged.
\end{minipage}
\end{table*}
\endgroup

 {
Table~\ref{tab:graph_attention_validation} presents the controlled ablation
results separately for each dataset. Relative to the non-graph classifier,
the complete framework improved Accuracy by \(1.51\), \(2.37\), and \(2.28\)
percentage points and F1-score by \(3.21\), \(3.57\), and \(3.27\) percentage
points on DF-1, DF-2, and DF-3, respectively. These consistently positive
differences quantify the combined contribution of patient-similarity graph
construction, attention-based aggregation, and similarity-weighted message
passing relative to independent feature-based classification.
}

 {
The contribution of attention was assessed by comparing the complete
framework with the otherwise equivalent configuration without attention.
Introducing learned attention increased Recall by \(1.46\), \(2.24\), and
\(1.97\) percentage points and PR-AUC by \(1.22\), \(1.56\), and \(1.42\)
percentage points on DF-1, DF-2, and DF-3, respectively. The ablation results
suggest that adaptive neighbor-specific coefficients provide additional
predictive information beyond fixed edge-weighted aggregation.
}

 {
The contribution of cosine-similarity weighting was evaluated by comparing
the complete framework with the unweighted-attention configuration. Restoring
the edge weights increased F1-score by \(0.50\), \(0.71\), and \(0.57\)
percentage points and PR-AUC by \(0.45\), \(0.66\), and \(0.69\) percentage
points on DF-1, DF-2, and DF-3, respectively. Although these differences were
smaller than the attention-related gains, their consistent direction suggests
that preserving patient-similarity strength contributes to neighborhood
aggregation across the three datasets.
}

 {
Overall, every component ablation reduced the reported performance on each
dataset. The direction of the differences was therefore consistent across
DF-1, DF-2, and DF-3, indicating that the aggregate improvement was not
driven solely by DF-1, the largest benchmark. Together with the
synthetic-edge audit, these findings suggest complementary roles for
DHRS-informed graph construction, graph-based patient representation,
adaptive attention, and cosine-similarity weighting within GNN-DHRS.
Nevertheless, these controlled ablations establish associations under the
evaluated experimental conditions rather than causal or prospective clinical
effects.
}
%---------------------------------------------------------

%=========================================================
\subsection{Feature-Selection and Graph-Topology Analysis}
\label{subsec:feature_selection_analysis}

 {
Feature selection determines both the model input and the feature space used
to construct patient-similarity neighborhoods. Its influence was therefore
evaluated jointly at the predictive and graph-topology levels using the top 5,
top 10, and top 15 ANOVA-ranked transformed features and all candidate
features without ANOVA selection. Only the retained feature set was varied;
the data partitions, training-only preprocessing, DHRS configuration, cosine
\(k\)-NN graph construction, model architecture, optimization settings, and
threshold-selection procedure remained unchanged.
}

 {
Within each outer fold, correlation filtering and ANOVA ranking were fitted
using the inner-training data. The feature-count configuration was selected
according to inner-validation F1-score and subsequently fixed before
outer-test evaluation. After selection, feature ranking was refitted on the
complete outer-training partition, and the retained columns were applied
unchanged to the corresponding outer-test fold.
}

\begingroup
 
\begin{table*}[!ht]
\centering
\caption{
Joint predictive and graph-topology analysis of feature-count configurations.
Predictive measures are calculated from untouched outer-test folds, whereas
topology measures are calculated from the corresponding post-DHRS
outer-training graphs. Values are mean \(\pm\) standard deviation.
}
\label{tab:feature_selection_joint_analysis}
\label{tab:feature_count_performance}
\label{tab:feature_count_topology}
\scriptsize
\setlength{\tabcolsep}{2.5pt}
\renewcommand{\arraystretch}{1.05}
\resizebox{\textwidth}{!}{%
\begin{tabular}{@{}llccccccc@{}}
\toprule
&
&
\multicolumn{4}{c}{\textbf{Predictive performance (\%)}} &
\multicolumn{3}{c}{\textbf{Training-graph topology}} \\
\cmidrule(lr){3-6}
\cmidrule(lr){7-9}
\textbf{Dataset} &
\textbf{Features} &
\textbf{Recall} &
\textbf{F1-score} &
\textbf{PR-AUC} &
\textbf{MCC} &
\shortstack{\textbf{Minority}\\\textbf{homophilic degree}} &
\shortstack{\textbf{Edge}\\\textbf{homophily}} &
\shortstack{\textbf{Connected}\\\textbf{components}} \\
\midrule

\multirow{4}{*}{DF-1}
& Top 5
& $97.42\pm0.48$
& $97.85\pm0.41$
& $98.31\pm0.35$
& $96.18\pm0.52$
& $6.84\pm0.42$
& $0.875\pm0.014$
& $24.6\pm2.8$ \\

& \textbf{Top 10}
& \textbf{$99.31\pm0.15$}
& \textbf{$99.39\pm0.08$}
& \textbf{$99.87\pm0.04$}
& \textbf{$98.78\pm0.14$}
& \textbf{$8.67\pm0.35$}
& \textbf{$0.912\pm0.009$}
& \textbf{$17.8\pm2.1$} \\

& Top 15
& $98.76\pm0.29$
& $98.92\pm0.24$
& $99.41\pm0.18$
& $97.85\pm0.31$
& $8.15\pm0.38$
& $0.897\pm0.011$
& $19.5\pm2.4$ \\

& All
& $98.15\pm0.37$
& $98.38\pm0.32$
& $98.94\pm0.26$
& $97.12\pm0.39$
& $7.42\pm0.45$
& $0.868\pm0.015$
& $22.9\pm2.9$ \\
\midrule

\multirow{4}{*}{DF-2}
& Top 5
& $96.18\pm0.58$
& $95.74\pm0.51$
& $96.42\pm0.44$
& $93.15\pm0.62$
& $6.51\pm0.46$
& $0.859\pm0.016$
& $15.8\pm2.1$ \\

& \textbf{Top 10}
& \textbf{$98.96\pm0.24$}
& \textbf{$98.39\pm0.17$}
& \textbf{$98.71\pm0.19$}
& \textbf{$96.42\pm0.38$}
& \textbf{$8.29\pm0.38$}
& \textbf{$0.897\pm0.011$}
& \textbf{$10.6\pm1.6$} \\

& Top 15
& $97.85\pm0.35$
& $97.52\pm0.31$
& $98.05\pm0.26$
& $95.28\pm0.42$
& $7.78\pm0.41$
& $0.882\pm0.013$
& $12.3\pm1.9$ \\

& All
& $97.12\pm0.43$
& $96.78\pm0.38$
& $97.41\pm0.33$
& $94.51\pm0.49$
& $7.05\pm0.48$
& $0.851\pm0.017$
& $14.7\pm2.3$ \\
\midrule

\multirow{4}{*}{DF-3}
& Top 5
& $96.45\pm0.54$
& $96.08\pm0.48$
& $96.75\pm0.41$
& $93.68\pm0.58$
& $6.63\pm0.44$
& $0.864\pm0.015$
& $16.4\pm2.2$ \\

& \textbf{Top 10}
& \textbf{$99.12\pm0.16$}
& \textbf{$98.52\pm0.41$}
& \textbf{$98.84\pm0.22$}
& \textbf{$96.91\pm0.51$}
& \textbf{$8.41\pm0.36$}
& \textbf{$0.902\pm0.010$}
& \textbf{$11.2\pm1.7$} \\

& Top 15
& $98.15\pm0.32$
& $97.81\pm0.29$
& $98.28\pm0.24$
& $95.74\pm0.39$
& $7.92\pm0.39$
& $0.888\pm0.012$
& $13.1\pm2.0$ \\

& All
& $97.48\pm0.41$
& $97.12\pm0.36$
& $97.65\pm0.31$
& $94.92\pm0.46$
& $7.18\pm0.46$
& $0.857\pm0.016$
& $15.3\pm2.4$ \\
\bottomrule
\end{tabular}%
}

\vspace{2pt}
\begin{minipage}{0.98\textwidth}
\footnotesize
\textit{Note:} Feature counts refer to transformed columns after encoding and
correlation filtering; ``All'' omits only ANOVA selection. Minority
homophilic degree is the mean number of same-class neighbors of a stroke node.
Edge homophily is the proportion of edges connecting nodes with the same
label. Graph statistics were calculated from post-DHRS cosine \(k\)-NN
training graphs with \(k=10\). Bold values identify the strongest mean within
each dataset.
\end{minipage}
\end{table*}
\endgroup

 {
Table~\ref{tab:feature_selection_joint_analysis} shows a consistent
association between the predictive and structural effects of feature
selection. The top-10 configuration produced the strongest Recall, F1-score,
PR-AUC, and MCC on all three datasets. It also yielded the highest minority
homophilic degree and edge homophily and the fewest connected components.
Reducing the representation to five features was associated with lower
predictive performance, weaker same-class connectivity, and greater graph
fragmentation. Conversely, retaining 15 or all candidate features did not
improve performance or topology relative to the validation-selected
configuration.
}

\begin{figure*}[!ht]
\centering
\includegraphics[width=0.98\textwidth]
{images/feature_selection_topology_analysis.png}
\caption{ {
Predictive and topological sensitivity to the retained feature count across
DF-1, DF-2, and DF-3: (a) outer-test F1-score and PR-AUC and
(b) post-DHRS training-graph edge homophily and connected-component count.
Error bars represent standard deviations across folds; the dashed line marks
the validation-selected top-10 configuration.
}}
\label{fig:feature_selection_topology}
\end{figure*}

 {
Figure~\ref{fig:feature_selection_topology} visually summarizes the joint
pattern reported in Table~\ref{tab:feature_selection_joint_analysis}. Across
the three datasets, predictive performance and edge homophily increased from
the top-5 to the top-10 configuration and subsequently declined when additional
features were retained, whereas graph fragmentation followed the opposite
pattern. The consistency of this profile across datasets suggests that the
selected feature count was not driven by one benchmark alone.
}

 {
Overall, the results suggest that ten transformed features provided an
effective balance between preserving discriminative information and limiting
dimensions that weakened the similarity structure. The concurrent predictive
and topological changes are consistent with the interpretation that feature
selection influences GNN-DHRS through both the patient representation and the
graph derived from that representation. Because these quantities change
together, however, the analysis does not establish that improved graph
homophily alone caused the predictive improvement.
}
%=========================================================
%----------------------------------------
%=========================================================
\subsection{Graph, Hyperparameter, and Stochastic Sensitivity Analysis}
\label{subsec:k_neighbor_sensitivity}

The graph neighborhood size controls the density of the patient-similarity
graph and the amount of relational information available during message
passing. Small neighborhoods may restrict information exchange, whereas
excessively large neighborhoods may introduce weak or uninformative
connections.

 {
The sensitivity analysis was extended to the principal architectural,
optimization, resampling, and graph-construction hyperparameters of
GNN-DHRS: attention heads \(H\), hidden dimension \(d_h\), dropout rate \(p\),
initial learning rate \(\eta\), weight decay \(\lambda\), DHRS oversampling
ratio \(\rho\), and graph neighborhood size \(k_g\). A one-factor-at-a-time
design was used: one parameter was varied while all remaining settings were
fixed at their validation-selected values.
}

 {
Hyperparameter selection was performed exclusively using inner-validation
predictions. The outer-test evaluations reported below were conducted only
after the selected configuration had been fixed and serve as a post-selection
robustness analysis. All configurations used identical outer folds,
training-only preprocessing, feature selection, strictly inductive graph
construction, and validation-selected decision thresholds.
}

\begingroup
 
\begin{table*}[!ht]
\centering
\caption{
One-factor-at-a-time sensitivity of GNN-DHRS to architectural, optimization,
resampling, and graph parameters. Values are cross-dataset means
\(\pm\) standard deviations (\%) from untouched outer-test evaluations.
}
\label{tab:hyperparameter_sensitivity}
\scriptsize
\setlength{\tabcolsep}{4.2pt}
\renewcommand{\arraystretch}{1.03}
\begin{tabular}{@{}llccc@{}}
\toprule
\textbf{Hyperparameter} &
\textbf{Value} &
\textbf{Recall} &
\textbf{F1-score} &
\textbf{PR-AUC} \\
\midrule

\multirow{4}{*}{Attention heads \(H\)}
& 2
& $98.42\pm0.38$ & $98.15\pm0.34$ & $98.68\pm0.29$ \\
& 4
& $98.87\pm0.29$ & $98.61\pm0.26$ & $99.12\pm0.22$ \\
& \textbf{8}
& \textbf{$99.13\pm0.21$}
& \textbf{$98.77\pm0.18$}
& \textbf{$99.14\pm0.15$} \\
& 12
& $99.05\pm0.24$ & $98.69\pm0.21$ & $99.28\pm0.18$ \\
\midrule

\multirow{4}{*}{Hidden dimension \(d_h\)}
& 32
& $98.15\pm0.45$ & $97.82\pm0.41$ & $98.35\pm0.36$ \\
& 64
& $98.68\pm0.33$ & $98.41\pm0.29$ & $98.95\pm0.25$ \\
& \textbf{128}
& \textbf{$99.13\pm0.21$}
& \textbf{$98.77\pm0.18$}
& \textbf{$99.14\pm0.15$} \\
& 256
& $98.94\pm0.27$ & $98.58\pm0.24$ & $99.18\pm0.19$ \\
\midrule

\multirow{4}{*}{Dropout \(p\)}
& 0.10
& $98.71\pm0.32$ & $98.45\pm0.28$ & $99.05\pm0.24$ \\
& \textbf{0.30}
& \textbf{$99.13\pm0.21$}
& \textbf{$98.77\pm0.18$}
& \textbf{$99.14\pm0.15$} \\
& 0.50
& $98.85\pm0.26$ & $98.52\pm0.23$ & $99.15\pm0.19$ \\
& 0.70
& $98.28\pm0.39$ & $97.95\pm0.35$ & $98.52\pm0.31$ \\
\midrule

\multirow{4}{*}{Learning rate \(\eta\)}
& \(10^{-4}\)
& $98.35\pm0.41$ & $98.08\pm0.37$ & $98.61\pm0.32$ \\
& \(5\times10^{-4}\)
& $98.79\pm0.28$ & $98.54\pm0.25$ & $99.12\pm0.21$ \\
& \(\mathbf{10^{-3}}\)
& \textbf{$99.13\pm0.21$}
& \textbf{$98.77\pm0.18$}
& \textbf{$99.14\pm0.15$} \\
& \(5\times10^{-3}\)
& $98.52\pm0.36$ & $98.19\pm0.32$ & $98.78\pm0.28$ \\
\midrule

\multirow{3}{*}{Weight decay \(\lambda\)}
& \(10^{-4}\)
& $98.84\pm0.29$ & $98.55\pm0.25$ & $99.06\pm0.22$ \\
& \(\mathbf{5\times10^{-4}}\)
& \textbf{$99.13\pm0.21$}
& \textbf{$98.77\pm0.18$}
& \textbf{$99.14\pm0.15$} \\
& \(10^{-3}\)
& $98.91\pm0.27$ & $98.61\pm0.23$ & $99.08\pm0.20$ \\
\midrule

\multirow{4}{*}{Oversampling ratio \(\rho\)}
& 0.50
& $97.68\pm0.47$ & $97.35\pm0.42$ & $97.95\pm0.38$ \\
& 0.75
& $98.51\pm0.31$ & $98.22\pm0.28$ & $98.85\pm0.24$ \\
& \textbf{1.00}
& \textbf{$99.13\pm0.21$}
& \textbf{$98.77\pm0.18$}
& \textbf{$99.14\pm0.15$} \\
& 1.25
& $98.74\pm0.27$ & $98.41\pm0.24$ & $99.05\pm0.20$ \\
\midrule

\multirow{5}{*}{Graph neighbors \(k_g\)}
& 5
& $97.86\pm0.35$ & $97.91\pm0.31$ & $98.25\pm0.28$ \\
& \textbf{10}
& \textbf{$99.13\pm0.21$}
& \textbf{$98.77\pm0.18$}
& \textbf{$99.14\pm0.15$} \\
& 15
& $98.72\pm0.26$ & $98.51\pm0.23$ & $99.12\pm0.19$ \\
& 20
& $98.41\pm0.29$ & $98.26\pm0.26$ & $98.85\pm0.22$ \\
& 25
& $98.08\pm0.33$ & $98.03\pm0.29$ & $98.51\pm0.25$ \\
\bottomrule
\end{tabular}

\vspace{2pt}
\begin{minipage}{0.97\textwidth}
\footnotesize
\textit{Note:} One parameter was varied at a time. Bold values identify the
configuration selected using inner-validation data, not necessarily the
largest individual outer-test metric. The ratio \(\rho\) is measured after
Borderline-SMOTE and before SMOTEENN cleaning.
\end{minipage}
\end{table*}
\endgroup

 {
Table~\ref{tab:hyperparameter_sensitivity} shows that the selected
architecture occupied a stable region of the evaluated parameter space.
Increasing the number of attention heads and hidden dimension improved Recall
and F1-score up to \(H=8\) and \(d_h=128\), whereas further increases produced
no consistent improvement. Similarly, moderate dropout and weight decay
provided a more favorable performance profile than either weak or excessive
regularization.
}

 {
The optimization results supported an initial learning rate of
\(\eta=10^{-3}\). Smaller values produced lower performance under the fixed
training budget, whereas \(5\times10^{-3}\) was associated with greater
variability and reduced performance. These results characterize empirical
sensitivity under the adopted optimization schedule and do not imply that the
same setting is universally optimal for other datasets or architectures.
}

 {
Sensitivity to \(\rho\) and \(k_g\) followed a similar pattern. Increasing
the oversampling ratio to \(1.00\) improved the minority-sensitive measures,
whereas increasing it further provided no additional benefit. For graph
construction, \(k_g=10\) produced the strongest joint Recall and F1-score.
Smaller neighborhoods restricted relational coverage, while larger
neighborhoods were associated with gradual performance reductions, consistent
with the inclusion of less informative relations. These interpretations
describe associations observed in the controlled experiment rather than
causal effects of graph density or synthetic-sample volume.
}

\begin{figure*}[!ht]
\centering
\includegraphics[width=0.98\textwidth]
{images/hyperparameter_sensitivity_analysis.png}
\caption{
Post-selection sensitivity of GNN-DHRS to (a) attention heads,
(b) hidden dimension, (c) dropout, (d) learning rate,
(e) oversampling ratio, and (f) graph neighborhood size. The marked
configurations were selected using inner-validation data; weight-decay
results are reported in Table~\ref{tab:hyperparameter_sensitivity}.
}
\label{fig:hyperparameter_sensitivity}
\end{figure*}

 {
Figure~\ref{fig:hyperparameter_sensitivity} provides a visual summary of the
one-factor-at-a-time analysis. Performance generally improved as each
parameter approached its selected setting and declined moderately beyond that
region. The absence of abrupt deterioration around the selected values
suggests that GNN-DHRS is not dependent on an isolated, narrowly tuned
configuration within the examined ranges.
}

%---------------------------------------------------------
\subsubsection{ {Stochastic Stability}}
\label{subsubsec:random_seed_stability}

{
Stochastic stability was evaluated by repeating the selected configuration
with the predefined seeds
\(\mathcal{S}=\{11,22,33,44,55\}\). The same outer-fold assignments were
retained across repetitions so that the analysis isolated variability arising
from parameter initialization, stochastic optimization, and DHRS sample
generation rather than variation in the test partitions.
\textcolor{red}{
No individual seed is treated as uniquely representative. The results in
Table~\ref{tab:random_seed_stability} provide the seed-aggregated performance
of the selected configuration and constitute the primary summary of
stochastic stability, whereas Table~\ref{tab:gnndhrs_overall_results}
provides the detailed ten-fold results from the reference evaluation run.
}
}

\begin{table*}[!ht]
\centering
\caption{
\textcolor{red}{
Seed-aggregated performance of the selected GNN-DHRS configuration across
the five predefined stochastic seeds
\(\{11,22,33,44,55\}\). Values are mean \(\pm\) standard deviation (\%).
}
}
\label{tab:random_seed_stability}
\small
\setlength{\tabcolsep}{6pt}
\renewcommand{\arraystretch}{1.08}
\begin{tabular}{@{}lcccc@{}}
\toprule
\textbf{Dataset} &
\textbf{Recall} &
\textbf{F1-score} &
\textbf{PR-AUC} &
\textbf{MCC} \\
\midrule
DF-1 &
 {$99.28\pm0.12$} &
 {$99.35\pm0.09$} &
 {$99.84\pm0.06$} &
 {$98.71\pm0.15$} \\
DF-2 &
 {$98.91\pm0.18$} &
 {$98.32\pm0.16$} &
 {$98.65\pm0.17$} &
 {$96.28\pm0.29$} \\
DF-3 &
 {$99.07\pm0.14$} &
 {$98.45\pm0.22$} &
 {$98.79\pm0.19$} &
 {$96.78\pm0.34$} \\
\bottomrule
\end{tabular}

\vspace{2pt}
\begin{minipage}{0.97\textwidth}
\footnotesize
\textit{Note:} Identical outer folds were used for all seeds; only stochastic
training and DHRS operations varied.
\end{minipage}
\end{table*}

\textcolor{red}{
Table~\ref{tab:random_seed_stability} shows limited stochastic variability
across the five predefined seeds. The between-seed standard deviations ranged
from \(0.06\) to \(0.34\) percentage points, and the seed-aggregated estimates
remained closely aligned with the corresponding ten-fold reference results in
Table~\ref{tab:gnndhrs_overall_results}. This agreement indicates that the
observed performance pattern was not driven by a single favorable
initialization or DHRS realization. Accordingly, the five-seed estimates are
used as the primary summary of stochastic stability, while the ten-fold
reference results are retained to characterize variation across outer-test
partitions.
}

 {
Overall, the hyperparameter and seed analyses provide complementary evidence.
The one-factor-at-a-time experiment evaluates local robustness around the
validation-selected configuration, whereas the seed experiment assesses
reproducibility under stochastic training variation. Both analyses remain
specific to the examined parameter ranges, datasets, and experimental
protocol.
}
%=========================================================
%-----------------------------
%---------------------------------------------------------
\subsection{Graph-Definition and DHRS Topology Analysis}
\label{subsec:graph_definition_comparison}

 {
The sensitivity of GNN-DHRS to the patient-graph definition was evaluated
using five constructions: the proposed cosine \(k\)-nearest-neighbor
(\(k\)-NN) graph, Euclidean \(k\)-NN, Gower \(k\)-NN, mutual cosine
\(k\)-NN, and a cosine similarity-threshold graph. These constructions and
their corresponding similarity rules are defined in
Section~\ref{subsec:graph_construction}. All configurations used identical
outer folds, inner-validation partitions, node features, DHRS outputs, model
architecture, optimization settings, random seeds, and validation-selected
classification thresholds; only the graph-construction rule was varied.
}

 {
The cosine, Euclidean, Gower, and mutual-cosine \(k\)-NN graphs used
\(k=10\). Gower neighborhoods were calculated from the mixed-type clinical
predictors using numerical ranges and categorical definitions fitted on the
training partition, whereas the node features supplied to the GNN remained
identical across all graph variants. For the similarity-threshold graph, the
threshold \(\tau\) was estimated from training data to approximate the mean
degree of the cosine \(k\)-NN reference graph. Neither validation nor
outer-test observations were used to determine the graph threshold.
}

 {
During strictly inductive inference, each outer-test patient was connected
only to original outer-training nodes according to the selected graph rule.
No test--test connections, test labels, or target-derived graph statistics
were used. The classification decision threshold was selected separately
from inner-validation predictions and was fixed before evaluating the
corresponding outer-test fold.
}

\begin{table*}[!ht]
\centering
\caption{ {
Predictive sensitivity of GNN-DHRS to the patient-graph definition.
Values are mean \(\pm\) standard deviation across ten untouched outer-test
folds.
}}
\label{tab:graph_definition_performance}
\scriptsize
\setlength{\tabcolsep}{2.8pt}
\renewcommand{\arraystretch}{1.05}
\resizebox{\textwidth}{!}{%
\begin{tabular}{@{}llcccccc@{}}
\toprule
\textbf{Dataset} &
\textbf{Graph definition} &
\textbf{Recall (\%)} &
\textbf{F1-score (\%)} &
\textbf{PR-AUC (\%)} &
\textbf{MCC (\%)} &
\textbf{BAcc. (\%)} &
\textbf{ROC-AUC (\%)} \\
\midrule

\multirow{5}{*}{DF-1}
& Euclidean \(k\)-NN
& $98.45\pm0.31$ & $98.62\pm0.27$ & $99.18\pm0.21$
& $97.42\pm0.35$ & $98.95\pm0.23$ & $99.84\pm0.09$ \\
& Gower \(k\)-NN
& $98.72\pm0.28$ & $98.85\pm0.23$ & $99.35\pm0.18$
& $97.68\pm0.31$ & $99.12\pm0.20$ & $99.89\pm0.07$ \\
& Mutual cosine \(k\)-NN
& $98.91\pm0.24$ & $99.02\pm0.19$ & $99.58\pm0.13$
& $98.05\pm0.27$ & $99.28\pm0.17$ & $99.94\pm0.04$ \\
& Cosine-threshold
& $98.63\pm0.29$ & $98.78\pm0.24$ & $99.42\pm0.16$
& $97.81\pm0.30$ & $99.08\pm0.21$ & $99.91\pm0.06$ \\
& \textbf{Proposed cosine \(k\)-NN}
& \textbf{$99.31\pm0.15$} & \textbf{$99.39\pm0.08$}
& \textbf{$99.87\pm0.04$} & \textbf{$98.78\pm0.14$}
& \textbf{$99.65\pm0.07$} & \textbf{$100.00\pm0.00$} \\
\midrule

\multirow{5}{*}{DF-2}
& Euclidean \(k\)-NN
& $97.25\pm0.42$ & $96.88\pm0.36$ & $97.45\pm0.31$
& $94.72\pm0.48$ & $97.81\pm0.33$ & $99.52\pm0.17$ \\
& Gower \(k\)-NN
& $97.58\pm0.38$ & $97.21\pm0.32$ & $97.78\pm0.27$
& $95.18\pm0.43$ & $98.05\pm0.29$ & $99.65\pm0.14$ \\
& Mutual cosine \(k\)-NN
& $98.15\pm0.31$ & $97.78\pm0.27$ & $98.28\pm0.23$
& $95.71\pm0.39$ & $98.42\pm0.25$ & $99.78\pm0.11$ \\
& Cosine-threshold
& $97.72\pm0.36$ & $97.38\pm0.31$ & $97.92\pm0.26$
& $95.35\pm0.42$ & $98.14\pm0.28$ & $99.69\pm0.13$ \\
& \textbf{Proposed cosine \(k\)-NN}
& \textbf{$98.96\pm0.24$} & \textbf{$98.39\pm0.17$}
& \textbf{$98.71\pm0.19$} & \textbf{$96.42\pm0.38$}
& \textbf{$98.90\pm0.15$} & \textbf{$99.92\pm0.08$} \\
\midrule

\multirow{5}{*}{DF-3}
& Euclidean \(k\)-NN
& $97.42\pm0.45$ & $97.08\pm0.39$ & $97.62\pm0.34$
& $94.95\pm0.51$ & $98.02\pm0.36$ & $99.58\pm0.16$ \\
& Gower \(k\)-NN
& $97.75\pm0.41$ & $97.42\pm0.35$ & $97.95\pm0.29$
& $95.38\pm0.46$ & $98.25\pm0.31$ & $99.69\pm0.13$ \\
& Mutual cosine \(k\)-NN
& $98.28\pm0.34$ & $97.95\pm0.30$ & $98.42\pm0.25$
& $95.92\pm0.41$ & $98.58\pm0.27$ & $99.82\pm0.10$ \\
& Cosine-threshold
& $97.88\pm0.39$ & $97.55\pm0.34$ & $98.08\pm0.28$
& $95.52\pm0.44$ & $98.34\pm0.30$ & $99.73\pm0.12$ \\
& \textbf{Proposed cosine \(k\)-NN}
& \textbf{$99.12\pm0.16$} & \textbf{$98.52\pm0.41$}
& \textbf{$98.84\pm0.22$} & \textbf{$96.91\pm0.51$}
& \textbf{$99.27\pm0.19$} & \textbf{$99.93\pm0.07$} \\
\bottomrule
\end{tabular}%
}

\vspace{2pt}

\begin{minipage}{0.98\textwidth}
\footnotesize
\textit{Note:} The four \(k\)-NN variants used \(k=10\). For the
cosine-threshold graph, \(\tau\) was fitted from training data to approximate
the mean degree of the cosine \(k\)-NN reference graph. The graph threshold
is distinct from the classification threshold, which was selected using
inner-validation predictions. BAcc. denotes balanced accuracy.
\end{minipage}
\end{table*}

 {
Table~\ref{tab:graph_definition_performance} shows that predictive performance
varied with the graph-construction rule. The proposed cosine \(k\)-NN graph
achieved the highest values for all six reported measures on DF-1, DF-2, and
DF-3. Mutual cosine \(k\)-NN was the strongest alternative, whereas the
Euclidean graph generally produced the lowest minority-sensitive performance.
These results indicate that both the similarity measure and the neighborhood
retention rule influence the information available for graph-attention
aggregation.
}

 {
The structural effects of DHRS were subsequently examined within each graph
definition. \textcolor{red}{
The topology analysis was performed using the same graph support employed by
the corresponding predictive configuration. For the primary cosine
\(k\)-NN graph, topology statistics were therefore calculated from the
directed edge set defined in Section~\ref{subsec:graph_construction}.
Reciprocal relations were counted as two directed edges when both directions
were independently retained. Self-loops introduced internally by the GAT
layers were excluded from these descriptive graph statistics.
}. For graph
\(G=(V,E)\), the minority homophilic degree was defined as
}

{ 
\begin{equation}
d_{\mathrm{hom,min}}
=
\frac{1}{|V_{\mathrm{min}}|}
\sum_{i\in V_{\mathrm{min}}}
\sum_{j:(i,j)\in E}
\mathbb{I}(y_j=y_i),
\label{eq:minority_homophilic_degree}
\end{equation}
}

 {
where \(V_{\mathrm{min}}\) is the set of stroke-class nodes. Edge homophily
was calculated as
}

{ 
\begin{equation}
h_E
=
\frac{1}{|E|}
\sum_{(i,j)\in E}
\mathbb{I}(y_i=y_j),
\label{eq:edge_homophily}
\end{equation}
}

 {
and graph fragmentation was summarized by the number of connected components.
All statistics were calculated exclusively from outer-training graphs;
validation and outer-test patients were excluded.
}

\begin{table*}[!ht]
\centering
\caption{ {
Descriptive topology of the outer-training graphs before and after DHRS.
Values are mean \(\pm\) standard deviation across the outer folds.
}}
\label{tab:graph_topology_dhrs}
\scriptsize
\setlength{\tabcolsep}{3.1pt}
\renewcommand{\arraystretch}{1.05}
\resizebox{\textwidth}{!}{%
\begin{tabular}{@{}lllccc@{}}
\toprule
\textbf{Dataset} &
\textbf{Graph definition} &
\textbf{Training graph} &
\shortstack{\textbf{Minority}\\\textbf{homophilic degree}} &
\shortstack{\textbf{Edge}\\\textbf{homophily}} &
\shortstack{\textbf{Connected}\\\textbf{components}} \\
\midrule

\multirow{10}{*}{DF-1}
& \multirow{2}{*}{Euclidean \(k\)-NN}
& Before DHRS & $1.92\pm0.21$ & $0.742\pm0.018$ & $38.7\pm3.4$ \\
&& After DHRS & $7.84\pm0.39$ & $0.884\pm0.012$ & $21.5\pm2.5$ \\
& \multirow{2}{*}{Gower \(k\)-NN}
& Before DHRS & $2.18\pm0.24$ & $0.758\pm0.017$ & $35.2\pm3.1$ \\
&& After DHRS & $8.02\pm0.37$ & $0.891\pm0.011$ & $20.1\pm2.3$ \\
& \multirow{2}{*}{Mutual cosine \(k\)-NN}
& Before DHRS & $1.74\pm0.19$ & $0.781\pm0.016$ & $46.8\pm4.2$ \\
&& After DHRS & $7.56\pm0.34$ & $0.923\pm0.008$ & $25.4\pm2.7$ \\
& \multirow{2}{*}{Cosine-threshold}
& Before DHRS & $2.35\pm0.26$ & $0.769\pm0.015$ & $31.6\pm2.9$ \\
&& After DHRS & $8.21\pm0.36$ & $0.905\pm0.010$ & $18.9\pm2.2$ \\
& \multirow{2}{*}{Proposed cosine \(k\)-NN}
& Before DHRS & $2.28\pm0.23$ & $0.773\pm0.014$ & $29.5\pm2.7$ \\
&& After DHRS & $8.67\pm0.35$ & $0.912\pm0.009$ & $17.8\pm2.1$ \\
\midrule

\multirow{10}{*}{DF-2}
& \multirow{2}{*}{Euclidean \(k\)-NN}
& Before DHRS & $1.68\pm0.22$ & $0.721\pm0.021$ & $26.4\pm2.8$ \\
&& After DHRS & $7.42\pm0.43$ & $0.867\pm0.014$ & $13.8\pm1.9$ \\
& \multirow{2}{*}{Gower \(k\)-NN}
& Before DHRS & $1.91\pm0.25$ & $0.739\pm0.019$ & $23.8\pm2.5$ \\
&& After DHRS & $7.68\pm0.41$ & $0.878\pm0.013$ & $12.7\pm1.8$ \\
& \multirow{2}{*}{Mutual cosine \(k\)-NN}
& Before DHRS & $1.52\pm0.20$ & $0.762\pm0.018$ & $31.2\pm3.3$ \\
&& After DHRS & $7.18\pm0.38$ & $0.908\pm0.010$ & $16.5\pm2.1$ \\
& \multirow{2}{*}{Cosine-threshold}
& Before DHRS & $2.05\pm0.27$ & $0.748\pm0.017$ & $21.5\pm2.4$ \\
&& After DHRS & $7.91\pm0.39$ & $0.889\pm0.012$ & $11.8\pm1.7$ \\
& \multirow{2}{*}{Proposed cosine \(k\)-NN}
& Before DHRS & $1.98\pm0.24$ & $0.753\pm0.016$ & $19.8\pm2.2$ \\
&& After DHRS & $8.29\pm0.38$ & $0.897\pm0.011$ & $10.6\pm1.6$ \\
\midrule

\multirow{10}{*}{DF-3}
& \multirow{2}{*}{Euclidean \(k\)-NN}
& Before DHRS & $1.72\pm0.23$ & $0.728\pm0.020$ & $27.8\pm2.9$ \\
&& After DHRS & $7.55\pm0.42$ & $0.872\pm0.014$ & $14.5\pm2.0$ \\
& \multirow{2}{*}{Gower \(k\)-NN}
& Before DHRS & $1.95\pm0.25$ & $0.744\pm0.019$ & $25.1\pm2.7$ \\
&& After DHRS & $7.81\pm0.40$ & $0.883\pm0.013$ & $13.4\pm1.9$ \\
& \multirow{2}{*}{Mutual cosine \(k\)-NN}
& Before DHRS & $1.57\pm0.21$ & $0.768\pm0.017$ & $32.6\pm3.5$ \\
&& After DHRS & $7.31\pm0.37$ & $0.913\pm0.010$ & $17.2\pm2.2$ \\
& \multirow{2}{*}{Cosine-threshold}
& Before DHRS & $2.09\pm0.26$ & $0.754\pm0.018$ & $22.7\pm2.5$ \\
&& After DHRS & $8.04\pm0.38$ & $0.894\pm0.012$ & $12.4\pm1.8$ \\
& \multirow{2}{*}{Proposed cosine \(k\)-NN}
& Before DHRS & $2.03\pm0.24$ & $0.759\pm0.016$ & $20.9\pm2.3$ \\
&& After DHRS & $8.41\pm0.36$ & $0.902\pm0.010$ & $11.2\pm1.7$ \\
\bottomrule
\end{tabular}%
}

\vspace{2pt}

\begin{minipage}{0.98\textwidth}
\footnotesize
\textcolor{red}{
\textit{Note:} Statistics were calculated from the same graph support used by
the corresponding predictive configuration. For the primary cosine
\(k\)-NN graph, this support is directed and is not automatically
symmetrized. GAT self-loops are excluded from the reported topology
statistics.
}
\end{minipage}
\end{table*}

 {
Across all graph definitions and datasets, the post-DHRS graphs exhibited
higher minority homophilic degree and edge homophily and fewer connected
components than their corresponding pre-DHRS graphs. The magnitude of these
changes differed by graph rule. Mutual cosine \(k\)-NN produced the highest
post-DHRS edge homophily but also retained more connected components, whereas
the proposed cosine \(k\)-NN graph provided a favorable combination of
minority connectivity, overall homophily, and lower fragmentation.
}

 {
Because DHRS changes both the number of nodes and the class distribution,
these before--after differences should not be interpreted as causal effects
of resampling on graph quality in isolation. Rather, they describe the
structural properties of the actual training graphs supplied to the model.
Together with the controlled predictive comparison in
Table~\ref{tab:graph_definition_performance}, the results identify cosine
\(k\)-NN as the best-performing graph definition among those evaluated under
the adopted protocol.
}
%---------------------------------------------------------
%---------------------------------------------------------
\subsection{Threshold Sensitivity and Probability Calibration}

\label{subsec:threshold_calibration}

 {
Threshold sensitivity was evaluated exclusively from inner-validation
predictions. Within each outer fold, predictions from the five inner-validation
folds were pooled so that every outer-training patient contributed once as an
inner-validation observation. These pooled predictions were used to select
one operating threshold for the corresponding outer fold.
}

The candidate threshold set was
\(\Theta=\{0.10,0.20,\ldots,0.90\}\). For outer fold \(f\), the threshold
maximizing the F1-score of the pooled inner-validation predictions was
selected as

\begin{equation}
\theta_f^{*}
=
\underset{\theta\in\Theta}{\operatorname{arg\,max}}
\;
\operatorname{F1}
\left(
\mathbf{y}_{\mathrm{val}}^{(f)},
\widehat{\mathbf{p}}_{\mathrm{val}}^{(f)};
\theta
\right),
\label{eq:validation_threshold_selection}
\end{equation}

where \(\mathbf{y}_{\mathrm{val}}^{(f)}\) and
\(\widehat{\mathbf{p}}_{\mathrm{val}}^{(f)}\) denote the pooled observed labels
and predicted probabilities obtained from the inner-validation folds,
respectively. The selected threshold was fixed before final model evaluation
and was applied without modification to the untouched outer-test patients.
No outer-test label, probability distribution, or test-derived performance
measure contributed to threshold selection. Unrounded validation results were
used in the optimization; when two thresholds produced exactly equal
F1-scores, the lower threshold was retained.

\begin{table*}[!ht]
\centering
\caption{
Threshold sensitivity of GNN-DHRS based on pooled inner-validation
predictions. Values are averages across the outer folds.
}
\label{tab:threshold_sensitivity}
\scriptsize
\setlength{\tabcolsep}{3.2pt}
\renewcommand{\arraystretch}{1.08}
\begin{tabular}{@{}cccccccccc@{}}
\toprule
&
\multicolumn{3}{c}{\textbf{DF-1}} &
\multicolumn{3}{c}{\textbf{DF-2}} &
\multicolumn{3}{c}{\textbf{DF-3}} \\
\cmidrule(lr){2-4}
\cmidrule(lr){5-7}
\cmidrule(lr){8-10}
\textbf{Threshold} &
\textbf{Prec.} & \textbf{Rec.} & \textbf{F1} &
\textbf{Prec.} & \textbf{Rec.} & \textbf{F1} &
\textbf{Prec.} & \textbf{Rec.} & \textbf{F1} \\
\midrule
\(0.10\) &
\(0.962\) & \(0.998\) & \(0.980\) &
\(0.941\) & \(0.997\) & \(0.968\) &
\(0.938\) & \(0.998\) & \(0.967\) \\
\(0.20\) &
\(0.974\) & \(0.997\) & \(0.985\) &
\(0.953\) & \(0.995\) & \(0.974\) &
\(0.951\) & \(0.996\) & \(0.973\) \\
\(0.30\) &
\(0.985\) & \(0.996\) & \(0.990\) &
\(0.964\) & \(0.993\) & \(0.978\) &
\(0.962\) & \(0.994\) & \(0.978\) \\
\(0.40\) &
\(0.991\) & \(0.995\) & \(0.993\) &
\(0.972\) & \(0.991\) & \(0.981\) &
\(0.971\) & \(0.992\) & \(0.981\) \\
\(\mathbf{0.50}\) &
\(\mathbf{0.994}\) & \(\mathbf{0.993}\) & \(\mathbf{0.993}\) &
\(\mathbf{0.978}\) & \(\mathbf{0.990}\) & \(\mathbf{0.984}\) &
\(\mathbf{0.979}\) & \(\mathbf{0.991}\) & \(\mathbf{0.985}\) \\
\(0.60\) &
\(0.996\) & \(0.988\) & \(0.992\) &
\(0.983\) & \(0.985\) & \(0.984\) &
\(0.984\) & \(0.986\) & \(0.985\) \\
\(0.70\) &
\(0.997\) & \(0.981\) & \(0.989\) &
\(0.987\) & \(0.978\) & \(0.982\) &
\(0.988\) & \(0.979\) & \(0.983\) \\
\(0.80\) &
\(0.998\) & \(0.968\) & \(0.983\) &
\(0.991\) & \(0.967\) & \(0.979\) &
\(0.991\) & \(0.968\) & \(0.979\) \\
\(0.90\) &
\(0.999\) & \(0.942\) & \(0.970\) &
\(0.994\) & \(0.948\) & \(0.970\) &
\(0.994\) & \(0.951\) & \(0.972\) \\
\bottomrule
\end{tabular}

\vspace{2pt}
\begin{minipage}{0.98\textwidth}
\footnotesize
\textit{Note:} Values were calculated from pooled inner-validation
predictions and are not outer-test results. Bold values identify the mean
selected operating region. Prec. and Rec. denote Precision and Recall,
respectively.
\end{minipage}
\end{table*}

 {
Table~\ref{tab:threshold_sensitivity} shows the expected operating-point
trade-off: increasing the threshold progressively increased Precision while
reducing Recall. F1-score was comparatively stable around thresholds
\(0.40\)--\(0.60\), with a mean selected threshold of approximately \(0.50\)
for each dataset. DF-1 selected \(0.50\) in every outer fold, whereas DF-2
and DF-3 showed limited fold-to-fold variation around this value.
}

\begin{figure*}[!ht]
\centering
\includegraphics[width=0.98\textwidth]
{images/threshold_sensitivity_all_datasets.png}
\caption{
Precision, Recall, and F1-score across candidate thresholds for
(a) DF-1, (b) DF-2, and (c) DF-3, calculated from pooled
inner-validation predictions. The dashed line denotes the mean
validation-selected threshold.
}
\label{fig:threshold_sensitivity}
\end{figure*}

Figure~\ref{fig:threshold_sensitivity} illustrates the threshold-dependent
trade-off summarized in Table~\ref{tab:threshold_sensitivity}. The operating
point was determined numerically using
Eq.~(\ref{eq:validation_threshold_selection}); curve intersections and visual
inspection were not used for threshold selection.

 {
Probability calibration was assessed, without post-hoc recalibration, using
the Brier score and Expected Calibration Error (ECE). These measures evaluate
the continuous predicted probabilities independently of the selected
classification threshold.
}

The Brier score was calculated as

\begin{equation}
\operatorname{Brier}
=
\frac{1}{N}
\sum_{i=1}^{N}
\left(
\widehat{p}_i-y_i
\right)^2,
\label{eq:brier_score_threshold_section}
\end{equation}

where \(\widehat{p}_i\) is the predicted stroke probability and
\(y_i\in\{0,1\}\) is the observed class label. ECE was calculated using
\(B=10\) equal-width probability bins:

\begin{equation}
\operatorname{ECE}
=
\sum_{b=1}^{B}
\frac{|S_b|}{N}
\left|
\operatorname{freq}(S_b)
-
\operatorname{conf}(S_b)
\right|,
\label{eq:ece_threshold_section}
\end{equation}

where \(S_b\) is the set of patients assigned to bin \(b\),
\(\operatorname{freq}(S_b)\) is the observed proportion of stroke-positive
patients in that bin, and \(\operatorname{conf}(S_b)\) is their mean predicted
stroke probability. Lower values indicate smaller probability-estimation and
calibration errors.

\begin{table*}[!ht]
\centering
\caption{
Validation-selected operating thresholds and outer-test probability
calibration of GNN-DHRS. Values are mean \(\pm\) standard deviation across
the outer folds.
}
\label{tab:threshold_calibration_results}
\small
\setlength{\tabcolsep}{6pt}
\renewcommand{\arraystretch}{1.10}
\begin{tabular}{@{}lccccc@{}}
\toprule
\textbf{Dataset} &
\shortstack{\textbf{Selected}\\\textbf{threshold}} &
\shortstack{\textbf{Test Precision}\\\textbf{(\%)}} &
\shortstack{\textbf{Test Recall}\\\textbf{(\%)}} &
\shortstack{\textbf{Brier}\\\textbf{score}} &
\shortstack{\textbf{ECE}\\\textbf{(\%)}} \\
\midrule
DF-1 &
\(0.50\pm0.00\) &
\(99.48\pm0.08\) &
\(99.31\pm0.15\) &
\(0.0068\pm0.0012\) &
\(1.24\pm0.31\) \\
DF-2 &
\(0.50\pm0.03\) &
\(97.84\pm0.18\) &
\(98.96\pm0.24\) &
\(0.0114\pm0.0021\) &
\(2.18\pm0.47\) \\
DF-3 &
\(0.50\pm0.04\) &
\(97.92\pm0.64\) &
\(99.12\pm0.16\) &
\(0.0107\pm0.0019\) &
\(1.97\pm0.42\) \\
\bottomrule
\end{tabular}

\vspace{2pt}
\begin{minipage}{0.96\textwidth}
\footnotesize
\textit{Note:} Thresholds were selected from inner-validation predictions
and fixed before outer-test evaluation. Test Precision and Recall were
calculated after applying the corresponding fixed fold-specific threshold.
Brier score and ECE were calculated from continuous outer-test probabilities
and were not used for model or threshold selection.
\end{minipage}
\end{table*}

\begin{figure*}[!ht]
\centering
\includegraphics[width=0.98\textwidth]
{images/calibration_curves_all_datasets.png}
\caption{
Reliability diagrams based on pooled untouched outer-test predictions for
(a) DF-1, (b) DF-2, and (c) DF-3. Empirical curves compare mean predicted
probability with observed stroke frequency in ten equal-width bins. The
diagonal denotes perfect calibration. Brier score and ECE were calculated
separately for each outer fold and are reported as mean \(\pm\) standard
deviation.
}
\label{fig:calibration_curves}
\end{figure*}

The mean Brier scores were \(0.0068\), \(0.0114\), and \(0.0107\), and the
corresponding ECE values were \(1.24\%\), \(2.18\%\), and \(1.97\%\) for
DF-1, DF-2, and DF-3, respectively. DF-1 showed the smallest probability
error and calibration discrepancy within the evaluated benchmarks. The
reliability diagrams in Figure~\ref{fig:calibration_curves} provide a
complementary descriptive assessment of agreement between predicted
probabilities and observed event frequencies.

 {
Because Brier score and ECE depend on event prevalence, sample size, and bin
occupancy, these results characterize within-benchmark calibration only.
They do not establish probability reliability under prospective deployment
or population shift. Evaluation on independent clinical cohorts and, where
necessary, training-only post-hoc recalibration would be required before the
probabilities could be interpreted as clinically calibrated risks.
}
%---------------------------------------------------------
\subsection{Confidence-Aware Selective Prediction}
\label{subsec:selective_prediction}

 {
Confidence-based selective prediction was evaluated to determine whether
uncertain cases could be deferred while preserving reliable predictions for
the remaining patients. For each target coverage, the confidence cutoff was
estimated exclusively from inner-validation predictions and fixed before
outer-test evaluation. Rejected cases were referred for clinical review rather
than being assigned automatically to either class.
}

\begingroup
 
\begin{table*}[!ht]
\centering
\caption{
Selective-prediction performance of GNN-DHRS. Values are mean \(\pm\)
standard deviation across the untouched outer-test folds.
}
\label{tab:selective_prediction_results}
\scriptsize
\setlength{\tabcolsep}{2.8pt}
\renewcommand{\arraystretch}{1.04}
\resizebox{\textwidth}{!}{%
\begin{tabular}{@{}lccccccc@{}}
\toprule
\textbf{Dataset} &
\shortstack{\textbf{Target}\\\textbf{coverage (\%)}} &
\shortstack{\textbf{Achieved}\\\textbf{coverage (\%)}} &
\shortstack{\textbf{Positive}\\\textbf{coverage (\%)}} &
\shortstack{\textbf{Selective}\\\textbf{Accuracy (\%)}} &
\shortstack{\textbf{Selective}\\\textbf{Recall (\%)}} &
\shortstack{\textbf{Selective}\\\textbf{F1-score (\%)}} &
\shortstack{\textbf{Selective}\\\textbf{risk (\%)}} \\
\midrule

\multirow{5}{*}{DF-1}
& 100 & $100.00\pm0.00$ & $100.00\pm0.00$
& $99.36\pm0.08$ & $99.31\pm0.15$ & $99.39\pm0.08$
& $0.64\pm0.08$ \\
& 95 & $95.12\pm0.41$ & $94.85\pm0.52$
& $99.68\pm0.06$ & $99.58\pm0.11$ & $99.71\pm0.05$
& $0.32\pm0.06$ \\
& 90 & $90.08\pm0.48$ & $89.71\pm0.61$
& $99.81\pm0.05$ & $99.72\pm0.09$ & $99.84\pm0.04$
& $0.19\pm0.05$ \\
& 80 & $80.15\pm0.55$ & $79.64\pm0.72$
& $99.91\pm0.03$ & $99.85\pm0.07$ & $99.92\pm0.03$
& $0.09\pm0.03$ \\
& 70 & $70.22\pm0.63$ & $69.58\pm0.81$
& $99.96\pm0.02$ & $99.92\pm0.05$ & $99.96\pm0.02$
& $0.04\pm0.02$ \\
\midrule

\multirow{5}{*}{DF-2}
& 100 & $100.00\pm0.00$ & $100.00\pm0.00$
& $98.85\pm0.16$ & $98.96\pm0.24$ & $98.39\pm0.17$
& $1.15\pm0.16$ \\
& 95 & $95.08\pm0.52$ & $94.62\pm0.68$
& $99.31\pm0.12$ & $99.28\pm0.18$ & $99.12\pm0.13$
& $0.69\pm0.12$ \\
& 90 & $90.11\pm0.59$ & $89.45\pm0.75$
& $99.52\pm0.10$ & $99.47\pm0.15$ & $99.38\pm0.11$
& $0.48\pm0.10$ \\
& 80 & $80.19\pm0.67$ & $79.28\pm0.89$
& $99.71\pm0.08$ & $99.65\pm0.12$ & $99.62\pm0.09$
& $0.29\pm0.08$ \\
& 70 & $70.25\pm0.74$ & $69.15\pm0.97$
& $99.83\pm0.06$ & $99.78\pm0.09$ & $99.79\pm0.07$
& $0.17\pm0.06$ \\
\midrule

\multirow{5}{*}{DF-3}
& 100 & $100.00\pm0.00$ & $100.00\pm0.00$
& $99.43\pm0.42$ & $99.12\pm0.16$ & $98.52\pm0.41$
& $0.57\pm0.42$ \\
& 95 & $95.15\pm0.49$ & $94.78\pm0.64$
& $99.72\pm0.11$ & $99.51\pm0.13$ & $99.21\pm0.15$
& $0.28\pm0.11$ \\
& 90 & $90.09\pm0.56$ & $89.62\pm0.71$
& $99.84\pm0.08$ & $99.68\pm0.11$ & $99.48\pm0.12$
& $0.16\pm0.08$ \\
& 80 & $80.17\pm0.64$ & $79.41\pm0.85$
& $99.92\pm0.05$ & $99.81\pm0.08$ & $99.71\pm0.09$
& $0.08\pm0.05$ \\
& 70 & $70.21\pm0.71$ & $69.32\pm0.93$
& $99.96\pm0.03$ & $99.89\pm0.06$ & $99.85\pm0.06$
& $0.04\pm0.03$ \\
\bottomrule
\end{tabular}%
}

\vspace{2pt}
\begin{minipage}{0.98\textwidth}
\footnotesize
\textit{Note:} Cutoffs were selected from inner-validation predictions.
Positive coverage is the proportion of stroke-positive test patients retained;
rejected cases were referred for review.
\end{minipage}
\end{table*}
\endgroup

 {
At \(100\%\) coverage, no prediction was rejected; therefore, the selective
results reproduce the primary GNN-DHRS performance. As coverage decreased,
patients with confidence values closest to the validation-selected decision
boundary were progressively deferred. At \(90\%\) target coverage, the
achieved coverage values were \(90.08\%\), \(90.11\%\), and \(90.09\%\) for
DF-1, DF-2, and DF-3, respectively. The corresponding selective risks declined
to \(0.19\%\), \(0.48\%\), and \(0.16\%\), compared with \(0.64\%\),
\(1.15\%\), and \(0.57\%\) when all patients received automatic predictions.
}

\begingroup
 
\begin{table}[!ht]
\centering
\caption{
Risk--coverage summary of GNN-DHRS selective prediction.
}
\label{tab:selective_aurc}
\small
\setlength{\tabcolsep}{5pt}
\renewcommand{\arraystretch}{1.06}
\begin{tabular}{@{}lccc@{}}
\toprule
\textbf{Dataset} &
\textbf{AURC} &
\shortstack{\textbf{Risk at 90\%}\\\textbf{coverage (\%)}} &
\shortstack{\textbf{Risk at 80\%}\\\textbf{coverage (\%)}} \\
\midrule
DF-1 & 0.0028 & $0.19\pm0.05$ & $0.09\pm0.03$ \\
DF-2 & 0.0051 & $0.48\pm0.10$ & $0.29\pm0.08$ \\
DF-3 & 0.0024 & $0.16\pm0.08$ & $0.08\pm0.05$ \\
\bottomrule
\end{tabular}

\vspace{2pt}
\begin{minipage}{0.94\linewidth}
\footnotesize
\textit{Note:} AURC was calculated from the complete pooled
risk--coverage curve using risk and coverage on the \([0,1]\) scale; lower
values indicate better selective performance.
\end{minipage}
\end{table}
\endgroup

 {
Table~\ref{tab:selective_aurc} summarizes selective performance across the
complete risk--coverage curve. AURC values of \(0.0028\), \(0.0051\), and
\(0.0024\) were obtained for DF-1, DF-2, and DF-3, respectively, with DF-3
showing the lowest aggregate selective risk. At \(80\%\) coverage, the risks
declined to \(0.09\%\), \(0.29\%\), and \(0.08\%\), respectively. These
reductions indicate that the deferred cases contained a comparatively greater
concentration of prediction errors than the accepted cases.
}

\begingroup
 
\begin{figure*}[!ht]
\centering
\includegraphics[width=0.98\textwidth]
{images/selective_prediction_risk_coverage.png}
\caption{
Selective risk versus achieved coverage for (a) DF-1, (b) DF-2, and
(c) DF-3. Cutoffs were selected from inner-validation predictions and applied
unchanged to untouched outer-test patients.
}
\label{fig:selective_risk_coverage}
\end{figure*}
\endgroup

 {
Figure~\ref{fig:selective_risk_coverage} shows a consistent decrease in
classification error as coverage is reduced. On DF-1, selective risk declined
from \(0.64\%\) at full coverage to \(0.04\%\) at \(70\%\) coverage. The
corresponding reductions were from \(1.15\%\) to \(0.17\%\) on DF-2 and from
\(0.57\%\) to \(0.04\%\) on DF-3. Thus, the confidence score provides useful
separation between reliable predictions and cases more likely to be
misclassified. DF-2 retained the highest risk across all coverage levels,
consistent with its comparatively lower full-coverage performance and higher
AURC.
}

 {
Positive-class coverage was examined to ensure that the reduction in risk was
not achieved by disproportionately rejecting stroke-positive patients. Across
the \(70\%\)--\(95\%\) target-coverage range, positive coverage varied from
\(69.58\%\) to \(94.85\%\) on DF-1, from \(69.15\%\) to \(94.62\%\) on DF-2,
and from \(69.32\%\) to \(94.78\%\) on DF-3. Positive coverage remained close
to achieved overall coverage, with a maximum difference of \(1.10\)
percentage points across all evaluated settings. The rejection mechanism
therefore did not exhibit substantial preferential deferral of stroke-positive
patients in these experiments.
}

 {
The selective-prediction mechanism is intended to complement clinical
assessment rather than replace it. Rejection does not indicate a non-stroke
outcome; instead, it identifies a case for which the model does not provide
sufficient confidence. Such patients should be referred for clinician review,
additional testing, or acquisition of further clinical information before an
automated prediction is acted upon.
}
%======================================================

%---------------------------------------------------------
\subsection{Graph-Based Explainability Experiments}
\label{subsec:explainability_results}

The explainability analysis evaluated the feature-level and relation-level
behavior of GNN-DHRS. Feature-level explanations were obtained using
\texttt{shap.KernelExplainer} with fold-specific, training-derived background
sets, whereas relation-level interpretations were derived from graph-attention
coefficients and evaluated through controlled edge-ablation analysis. The
feature-level analysis included 200 untouched outer-test patients per dataset,
corresponding to 600 patients across the complete evaluation.

 {
Each dataset contained ten raw modeling predictors after excluding the
identifier and outcome columns. Following categorical encoding and
training-fold feature selection, ten transformed features were supplied to
each fitted model. The explainability analysis therefore characterizes the
complete selected model input rather than a restricted subset of a larger
fitted representation. These analyses describe patterns learned by GNN-DHRS
and should not be interpreted as evidence of causal relationships or clinical
validity.
}

%=========================================================
% Combined feature-level explanation figure
%=========================================================
\begin{figure*}[!ht]
\centering

\begin{subfigure}[t]{0.48\textwidth}
\centering
\includegraphics[width=\linewidth]
{images/Figure_E1_Global_SHAP_Bar.png}
\caption{Global mean absolute Kernel SHAP importance.}
\label{fig:shap_global_panel}
\end{subfigure}
\hfill
\begin{subfigure}[t]{0.48\textwidth}
\centering
\includegraphics[width=\linewidth]
{images/Figure_E2_SHAP_Beeswarm.png}
\caption{SHAP contribution magnitude and direction.}
\label{fig:shap_beeswarm_panel}
\end{subfigure}

\vspace{4pt}

\begin{subfigure}[t]{0.48\textwidth}
\centering
\includegraphics[width=\linewidth]
{images/Figure_E3A_Local_SHAP_True_Positive.png}
\caption{Local attribution for a true-positive prediction.}
\label{fig:shap_local_tp_panel}
\end{subfigure}
\hfill
\begin{subfigure}[t]{0.48\textwidth}
\centering
\includegraphics[width=\linewidth]
{images/Figure_E3B_Local_SHAP_True_Negative.png}
\caption{Local attribution for a true-negative prediction.}
\label{fig:shap_local_tn_panel}
\end{subfigure}

\caption{
Kernel SHAP analysis of GNN-DHRS:
(a) global importance, (b) contribution distributions, and local explanations
of (c) true-positive and (d) true-negative predictions. Global results include
200 untouched outer-test patients per dataset.
}
\label{fig:shap_explainability_summary}
\end{figure*}

Figure~\ref{fig:shap_explainability_summary} summarizes the feature-level
explainability results. The global analysis in
Fig.~\ref{fig:shap_global_panel} shows that \texttt{age},
\texttt{avg\_glucose\_level}, and \texttt{bmi} received the largest mean
absolute attribution values, followed by hypertension, heart disease, and
smoking-related variables. The beeswarm distributions in
Fig.~\ref{fig:shap_beeswarm_panel} additionally show the magnitude and
direction of these contributions across the explained outer-test patients.

The local profiles in Figs.~\ref{fig:shap_local_tp_panel} and
\ref{fig:shap_local_tn_panel} illustrate how the global attribution patterns
were expressed in individual correct predictions. The true-positive example
exhibited stronger positive contributions from age, average glucose level,
BMI, hypertension, heart disease, and smoking-related variables, whereas the
true-negative example showed weaker positive or negative contributions from
these features.

 {
The local cases are included as illustrative examples of model behavior and
are not used to estimate global feature importance. The observed
feature-attribution patterns are broadly consistent with established stroke
risk factors; however, this agreement supports only the plausibility of the
learned associations and does not establish causality, clinical validity, or
suitability for patient-level decision-making.
}

%=========================================================
% Attention-relation examples
%=========================================================
 {
For relation-level interpretation, the attention coefficient associated with
each incoming edge was aggregated across the eight attention heads in the
final graph-attention layer. Edge rankings were then constructed separately
for each outer-test patient from its incoming training-to-test relations.
}

\begin{table}[!ht]
\centering
\caption{
Leading attention-weighted patient relations in a representative inference
graph.
}
\label{tab:top_attention_relations}
\small
\setlength{\tabcolsep}{4pt}
\renewcommand{\arraystretch}{1.06}
\begin{tabular}{@{}ccccc@{}}
\toprule
\textbf{Source} &
\textbf{Target} &
\textbf{Source class} &
\textbf{Target class} &
\textbf{Weight} \\
\midrule
314 & 110 & No stroke & No stroke & 0.9999 \\
110 & 314 & No stroke & No stroke & 0.9999 \\
421 & 37  & Stroke    & Stroke    & 0.9976 \\
37  & 421 & Stroke    & Stroke    & 0.9976 \\
58  & 421 & No stroke & Stroke    & 0.9948 \\
421 & 58  & Stroke    & No stroke & 0.9948 \\
\bottomrule
\end{tabular}

\vspace{2pt}
\begin{minipage}{0.94\linewidth}
\footnotesize
\textit{Note:} Source and target are anonymized node indices; weights are
final-layer attention coefficients averaged across heads.
\end{minipage}
\end{table}

%=========================================================
% Combined relational-interpretation figure
%=========================================================
\begin{figure*}[!ht]
\centering

\begin{subfigure}[t]{0.32\textwidth}
\centering
\includegraphics[width=\linewidth]
{images/Figure_E4A_Top_Attention_Relations.png}
\caption{Leading attention-weighted relations.}
\label{fig:attention_relations_panel}
\end{subfigure}
\hfill
\begin{subfigure}[t]{0.32\textwidth}
\centering
\includegraphics[width=\linewidth]
{images/Figure_E4B_Attention_Neighborhood_Subgraph.png}
\caption{Representative attention neighborhood.}
\label{fig:attention_subgraph_panel}
\end{subfigure}
\hfill
\begin{subfigure}[t]{0.32\textwidth}
\centering
\includegraphics[width=\linewidth]
{images/Figure_E5_Edge_Ablation_Bar.png}
\caption{Performance under controlled edge ablation.}
\label{fig:edge_ablation_panel}
\end{subfigure}

\caption{
Relational interpretation and edge-ablation analysis:
(a) leading attention-weighted relations, (b) a representative attention
neighborhood, and (c) performance after controlled edge removal. Ablation
results are mean \(\pm\) standard deviation.
}
\label{fig:relational_explainability_summary}
\end{figure*}

Table~\ref{tab:top_attention_relations} and
Figure~\ref{fig:relational_explainability_summary} provide complementary
illustrations of the relation-level behavior of GNN-DHRS.
Figure~\ref{fig:attention_relations_panel} displays highly weighted
patient-to-patient relations, while
Fig.~\ref{fig:attention_subgraph_panel} shows their organization within the
neighborhood of an individual outer-test patient. The displayed relations
include both same-class and cross-class connections, indicating that the
fitted model assigned unequal importance to neighboring training patients,
including relations located near class boundaries.

 {
Attention coefficients are model-internal relational weights. The displayed
relations illustrate how the fitted model distributes influence during
message passing but should not be interpreted as causal effects or
independently validated clinical relationships.
}

%=========================================================
% Dataset-wide controlled edge-ablation analysis
%=========================================================
\textcolor{red}{
For each outer-test patient, five incoming edges were removed under the
high-attention, low-attention, and random-edge strategies. Random-edge removal
was repeated ten independent times per patient, and the resulting prediction
measures were averaged at the patient level before fold-level aggregation.
The complete analysis was performed using frozen models without retraining or
accessing outer-test labels during edge ranking.
}

\begin{table*}[!ht]
\centering
\caption{
Edge-ablation validation using equal removal budgets. Values are mean
\(\pm\) standard deviation across outer-test folds.
}
\label{tab:edge_ablation_results}
\scriptsize
\setlength{\tabcolsep}{3.8pt}
\renewcommand{\arraystretch}{1.06}
\begin{tabular}{@{}llccc@{}}
\toprule
\textbf{Dataset} &
\textbf{Ablation strategy} &
\textbf{Recall} &
\textbf{F1-score} &
\textbf{ROC-AUC} \\
\midrule
\multirow{4}{*}{DF-1}
& Original graph
& \textbf{$0.993\pm0.002$}
& \textbf{$0.994\pm0.001$}
& \textbf{$1.000\pm0.000$} \\
& High-attention edges
& $0.958\pm0.008$
& $0.962\pm0.007$
& $0.984\pm0.005$ \\
& Low-attention edges
& $0.988\pm0.003$
& $0.990\pm0.002$
& $0.997\pm0.002$ \\
& Random edges
& $0.981\pm0.004$
& $0.985\pm0.003$
& $0.995\pm0.003$ \\
\midrule
\multirow{4}{*}{DF-2}
& Original graph
& \textbf{$0.990\pm0.002$}
& \textbf{$0.984\pm0.002$}
& \textbf{$0.999\pm0.001$} \\
& High-attention edges
& $0.942\pm0.011$
& $0.938\pm0.010$
& $0.972\pm0.008$ \\
& Low-attention edges
& $0.981\pm0.004$
& $0.976\pm0.004$
& $0.994\pm0.003$ \\
& Random edges
& $0.968\pm0.006$
& $0.962\pm0.006$
& $0.988\pm0.004$ \\
\midrule
\multirow{4}{*}{DF-3}
& Original graph
& \textbf{$0.991\pm0.002$}
& \textbf{$0.985\pm0.004$}
& \textbf{$0.999\pm0.001$} \\
& High-attention edges
& $0.947\pm0.010$
& $0.941\pm0.009$
& $0.975\pm0.007$ \\
& Low-attention edges
& $0.983\pm0.004$
& $0.978\pm0.005$
& $0.995\pm0.003$ \\
& Random edges
& $0.971\pm0.006$
& $0.965\pm0.006$
& $0.989\pm0.004$ \\
\bottomrule
\end{tabular}
\vspace{2pt}
\begin{minipage}{0.97\textwidth}
\footnotesize
\textcolor{red}{
\textit{Note:} Each strategy removed five incoming edges. Random-edge removal
was repeated ten independent times per patient and averaged before fold-level
aggregation. Edge ranking did not use outer-test labels.
}
\end{minipage}
\end{table*}

 {
As summarized in Table~\ref{tab:edge_ablation_results} and
Fig.~\ref{fig:edge_ablation_panel}, removing high-attention edges produced a
larger performance reduction than removing equal numbers of low-attention or
randomly selected edges. On DF-1, high-attention ablation reduced Recall from
\(0.993\) to \(0.958\) and ROC-AUC from \(1.000\) to \(0.984\). The same
ordered pattern—high-attention removal causing the largest degradation,
followed by random and then low-attention removal—was observed on DF-2 and
DF-3. These results indicate that high-attention edges carry greater
predictive influence within the fitted model and that the effect is
reproduced across all three datasets.
}

%=========================================================
% Comparison with established graph explainers
%=========================================================
To extend the analysis beyond attention-ranked ablation, attention-based edge
ranking was quantitatively compared with GNNExplainer, PGExplainer, and a
random-edge baseline across DF-1, DF-2, and DF-3. All methods used the same
frozen GNN-DHRS models, inductively constructed outer-test neighborhoods, and
an explanation budget of five incoming training-to-test edges per patient.

\begin{table*}[!ht]
\centering
\caption{
Fidelity and stability of the evaluated graph-explanation methods. Values are
mean \(\pm\) standard deviation across the untouched outer-test folds.
}
\label{tab:graph_explanation_comparison}
\small
\setlength{\tabcolsep}{5pt}
\renewcommand{\arraystretch}{1.08}
\begin{tabular}{@{}llccc@{}}
\toprule
\textbf{Dataset} &
\textbf{Explanation method} &
\shortstack{\(\mathbf{Fid^{+}}\)\\\textbf{(\%) \(\uparrow\)}} &
\shortstack{\(\mathbf{Fid^{-}}\)\\\textbf{(\%) \(\downarrow\)}} &
\shortstack{\textbf{Stability}\\\textbf{(\%) \(\uparrow\)}} \\
\midrule

\multirow{4}{*}{DF-1}
& Random edges
& $41.28\pm3.15$ & $38.64\pm2.87$ & $32.15\pm4.02$ \\
& GNNExplainer
& $68.74\pm2.41$ & $18.92\pm1.85$ & $71.48\pm2.67$ \\
& PGExplainer
& $71.35\pm2.18$ & $16.47\pm1.63$ & $74.82\pm2.35$ \\
& \textbf{Attention-based}
& \textbf{$79.62\pm1.74$}
& \textbf{$11.28\pm1.21$}
& \textbf{$86.35\pm1.58$} \\
\midrule

\multirow{4}{*}{DF-2}
& Random edges
& $38.95\pm3.42$ & $41.27\pm3.08$ & $29.84\pm4.25$ \\
& GNNExplainer
& $65.18\pm2.67$ & $21.45\pm2.04$ & $68.72\pm2.91$ \\
& PGExplainer
& $68.42\pm2.35$ & $18.93\pm1.79$ & $72.15\pm2.58$ \\
& \textbf{Attention-based}
& \textbf{$76.85\pm1.92$}
& \textbf{$13.64\pm1.45$}
& \textbf{$83.47\pm1.83$} \\
\midrule

\multirow{4}{*}{DF-3}
& Random edges
& $39.72\pm3.28$ & $40.15\pm2.95$ & $30.68\pm4.11$ \\
& GNNExplainer
& $66.51\pm2.54$ & $20.38\pm1.92$ & $69.85\pm2.78$ \\
& PGExplainer
& $69.87\pm2.27$ & $17.64\pm1.71$ & $73.41\pm2.46$ \\
& \textbf{Attention-based}
& \textbf{$78.24\pm1.81$}
& \textbf{$12.51\pm1.32$}
& \textbf{$85.12\pm1.69$} \\
\bottomrule
\end{tabular}

\vspace{2pt}
\begin{minipage}{0.97\textwidth}
\footnotesize
\textit{Note:} All methods used five edges. \(\operatorname{Fid}^{+}\),
\(\operatorname{Fid}^{-}\), and stability measure necessity, sufficiency, and
Jaccard agreement, respectively; arrows indicate preferred directions.
\end{minipage}
\end{table*}

As shown in Table~\ref{tab:graph_explanation_comparison}, attention-based edge
ranking achieved the highest explanation necessity, the lowest sufficiency
error, and the greatest stability across all three datasets. GNNExplainer and
PGExplainer consistently exceeded the random-edge baseline, with PGExplainer
producing stronger fidelity and stability than GNNExplainer.

 {
Under the common five-edge explanation budget, removing attention-selected
edges caused a larger reduction in predicted-class confidence than removing
edges selected by the comparison methods. Conversely, retaining only the
attention-selected edges reproduced the original prediction probabilities
more closely. These results support the functional relevance and
reproducibility of the selected relations with respect to the fitted
prediction function; they do not establish causal or clinically validated
patient relationships.
}

%=========================================================
% Cross-fold feature-ranking stability
%=========================================================
Cross-fold stability of the feature rankings was quantified using Jaccard
similarity:

\begin{equation}
J(a,b)
=
\frac{
\left|T_a\cap T_b\right|
}{
\left|T_a\cup T_b\right|
},
\label{eq:jaccard_stability}
\end{equation}

where \(T_a\) and \(T_b\) denote the sets of the five highest-ranked features
obtained from outer folds \(a\) and \(b\), respectively. Values approaching
one indicate greater agreement between the fold-specific ranking sets.

\begin{table}[!ht]
\centering
\caption{
Cross-fold stability of the leading Kernel SHAP feature rankings.
}
\label{tab:explanation_stability}
\small
\setlength{\tabcolsep}{10pt}
\renewcommand{\arraystretch}{1.06}
\begin{tabular}{@{}lc@{}}
\toprule
\textbf{Dataset} &
\textbf{Top-5 Jaccard similarity} \\
\midrule
DF-1 & 0.84 \\
DF-2 & 0.81 \\
DF-3 & 0.82 \\
\bottomrule
\end{tabular}

\vspace{2pt}
\begin{minipage}{0.94\linewidth}
\footnotesize
\textit{Note:} Values are mean pairwise similarities across outer folds for
the top five of ten selected features.
\end{minipage}
\end{table}

Table~\ref{tab:explanation_stability} indicates that the leading Kernel SHAP
rankings were relatively stable across the outer folds. Mean Top-5 Jaccard
similarity ranged from \(0.81\) to \(0.84\), with DF-1 exhibiting the highest
agreement. The repeated prominence of \texttt{age},
\texttt{avg\_glucose\_level}, and \texttt{bmi} further supports the
reproducibility of the observed feature-attribution patterns.

 {
Cross-fold stability and agreement with recognized stroke risk factors assess
reproducibility and plausibility, respectively; neither constitutes clinical
validation. Accordingly, the explainability findings should be interpreted as
descriptions of how GNN-DHRS uses the available benchmark variables and
patient-graph relations. Clinical validation would require independent review
by domain experts, evaluation on independently collected clinical cohorts,
and prospective investigation of whether the explanations remain reliable
and actionable across patient populations and care settings.
}
%=========================================================
%---------------------------------------------------------
\subsection{Statistical Validation}
\label{subsec:statistical_validation}

\textcolor{red}{
An exploratory paired statistical comparison was conducted between GNN-DHRS
and GAT-CBF under the same repeated nested cross-validation protocol. GAT-CBF
\cite{velickovic2018graph,lin2017focal,cui2019class} was selected after
inspection of the controlled baseline results because it exhibited the
highest observed performance among the evaluated imbalance-aware GAT
configurations. Because this comparator selection was informed by outer-test
performance, the subsequent inferential analysis is explicitly treated as
exploratory rather than confirmatory. Its purpose is to quantify the
magnitude, uncertainty, and consistency of the observed paired differences
under identical partitions and random seeds, rather than to represent a
prespecified confirmatory hypothesis test. Recall, F1-score, PR-AUC, and MCC
were examined because they characterize minority-class detection and
class-balanced performance under severe imbalance.
}

 {
Five repetitions of stratified 10-fold cross-validation were generated using
the predefined split seeds
\(\mathcal{S}=\{11,22,33,44,55\}\). Within each repetition, GNN-DHRS and
GAT-CBF used identical outer-training, inner-validation, and outer-test
partitions. Pairing was therefore performed by dataset, repetition, and outer
fold, yielding 50 paired outer-test differences for each metric and dataset.
}

 {
Because the training partitions overlap across folds and repetitions, the
paired differences cannot be treated as statistically independent. Inference
therefore used the corrected resampled \(t\)-test with the Nadeau--Bengio
variance correction \cite{nadeau2003inference}. For the adopted 10-fold
design, the correction incorporated an outer-test-to-training sample-size
ratio of approximately \(1/9\), thereby increasing the estimated uncertainty
to account for dependence among the cross-validation results.
}

 {
For each metric, Table~\ref{tab:statistical_comparison_results} reports four
complementary quantities. The mean difference is the average paired
performance of GNN-DHRS minus that of GAT-CBF and is expressed in percentage
points. The corrected 95\% confidence interval was calculated using the
Nadeau--Bengio dependence-adjusted standard error
\cite{nadeau2003inference}. The standardized paired effect size \(d_z\) was
calculated as the mean paired difference divided by the standard deviation of
the paired differences. Finally, the corrected-test \(p\)-values were adjusted
jointly across the 12 comparisons using Holm's sequential procedure
\cite{holm1979simple}.
}

\textcolor{red}{
The exploratory comparison family comprised four metrics across three
datasets, giving 12 paired statistical comparisons. Statistical significance
was defined as \(p_{\mathrm{Holm}}<0.05\). Interpretation considered the
direction and magnitude of the mean paired differences, dependence-adjusted
confidence intervals, standardized effect sizes, and multiplicity-adjusted
\(p\)-values jointly. Because GAT-CBF was selected after inspection of the
controlled baseline results, these inferential comparisons are interpreted as
exploratory rather than confirmatory.
}

 {
The corrected resampled test was selected instead of an ordinary paired
\(t\)-test because the latter would underestimate uncertainty by ignoring
overlap among the cross-validation training sets. A \(5\times2\)
cross-validation test was not adopted because it would replace the repeated
nested 10-fold protocol used throughout the study and train the models on
substantially smaller partitions. Directly bootstrapping the 50 fold-level
scores as independent observations was also avoided because such resampling
would not preserve their cross-validation dependence.
}

\begingroup
 
\begin{table*}[!ht]
\centering
\caption{
Corrected paired comparison of GNN-DHRS and GAT-CBF under repeated
stratified 10-fold cross-validation. Positive differences favor GNN-DHRS.
}
\label{tab:statistical_comparison_results}
\small
\setlength{\tabcolsep}{5pt}
\renewcommand{\arraystretch}{1.08}
\begin{tabular}{@{}llcccc@{}}
\toprule
\textbf{Dataset} &
\textbf{Metric} &
\shortstack{\textbf{Mean difference}\\\textbf{(pp)}} &
\shortstack{\textbf{Corrected 95\%}\\\textbf{CI (pp)}} &
\shortstack{\textbf{Effect size}\\\(\boldsymbol{d_z}\)} &
\shortstack{\textbf{Holm-adjusted}\\\(\boldsymbol{p}\)} \\
\midrule

\multirow{4}{*}{DF-1}
& Recall   & 5.42 & $[4.03,\;6.81]$ & 2.84 & $<0.001$ \\
& F1-score & 4.47 & $[3.22,\;5.72]$ & 2.61 & $<0.001$ \\
& PR-AUC   & 4.16 & $[2.96,\;5.36]$ & 2.53 & $<0.001$ \\
& MCC      & 4.50 & $[3.17,\;5.83]$ & 2.47 & $<0.001$ \\
\midrule

\multirow{4}{*}{DF-2}
& Recall   & 9.04 & $[6.93,\;11.15]$ & 3.12 & $<0.001$ \\
& F1-score & 7.15 & $[5.35,\;8.95]$ & 2.89 & $<0.001$ \\
& PR-AUC   & 6.78 & $[5.01,\;8.55]$ & 2.78 & $<0.001$ \\
& MCC      & 5.91 & $[4.22,\;7.60]$ & 2.54 & $<0.001$ \\
\midrule

\multirow{4}{*}{DF-3}
& Recall   & 9.61 & $[7.43,\;11.79]$ & 3.21 & $<0.001$ \\
& F1-score & 7.65 & $[5.76,\;9.54]$ & 2.95 & $<0.001$ \\
& PR-AUC   & 7.28 & $[5.41,\;9.15]$ & 2.84 & $<0.001$ \\
& MCC      & 6.72 & $[4.86,\;8.58]$ & 2.63 & $<0.001$ \\
\bottomrule
\end{tabular}

\vspace{2pt}
\begin{minipage}{0.96\textwidth}
\footnotesize
\textit{Note:} Results use 50 paired differences from five repeated
stratified 10-fold runs. CIs apply the Nadeau--Bengio correction, and Holm
adjustment covers all 12 comparisons. CI: confidence interval; pp: percentage
points.
\end{minipage}
\end{table*}
\endgroup

Table~\ref{tab:statistical_comparison_results} shows positive mean differences
for every evaluated metric and dataset. Compared with GAT-CBF, GNN-DHRS
increased Recall by \(5.42\), \(9.04\), and \(9.61\) percentage points on
DF-1, DF-2, and DF-3, respectively. Improvements ranged from \(4.47\) to
\(7.65\) percentage points for F1-score, from \(4.16\) to \(7.28\) percentage
points for PR-AUC, and from \(4.50\) to \(6.72\) percentage points for MCC.

 {
\textcolor{red}{
The consistently positive paired differences, dependence-adjusted confidence
intervals, standardized effect sizes, and multiplicity-adjusted \(p\)-values
indicate that GNN-DHRS produced higher observed values than GAT-CBF under the
evaluated repeated cross-validation protocol. Because GAT-CBF was selected
after inspection of the controlled baseline results, these inferential
findings are interpreted as exploratory and should not be considered evidence
from a prespecified confirmatory comparison.
}

 {
These results should nevertheless be interpreted within the scope of the
evaluated public benchmarks. Because DF-1, DF-2, and DF-3 are related
benchmark files rather than verified independent clinical cohorts, their
dataset-specific comparisons do not constitute independent external
replications. Evaluation on independently collected clinical cohorts remains
necessary to determine whether the observed differences generalize across
populations and data-collection settings.
}
%---------------------------------------------------------
\subsection{Computational Efficiency and Scalability}
\label{subsec:computational_efficiency}

 {
The computational profile of GNN-DHRS was evaluated in terms of DHRS
processing, graph construction, model fitting, peak GPU-memory utilization,
and strictly inductive inference. All measurements were obtained using the
common hardware and software environment documented in Supplementary
Table~S2. DF-1 provides the principal scalability assessment because its
outer-training partitions contain approximately \(39{,}000\) original
patients before DHRS.
}

\begingroup
 
\begin{table*}[!ht]
\centering
\caption{
Training-graph scale of GNN-DHRS across the evaluated datasets. Node and edge
counts are averages across the outer folds.
}
\label{tab:gnndhrs_scalability}
\small
\setlength{\tabcolsep}{7pt}
\renewcommand{\arraystretch}{1.08}
\begin{tabular}{@{}lrrrr@{}}
\toprule
\textbf{Dataset} &
\shortstack{\textbf{Original}\\\textbf{nodes}} &
\shortstack{\textbf{Post-DHRS}\\\textbf{nodes}} &
\shortstack{\textbf{Node}\\\textbf{increase (\%)}} &
\shortstack{\textbf{Directed \(k\)-NN}\\\textbf{links}} \\
\midrule
DF-1 & 39,060 & 60,800 & 55.66 & 608,000 \\
DF-2 & 4,483  & 6,808  & 51.86 & 68,080 \\
DF-3 & 4,599  & 6,981  & 51.79 & 69,810 \\
\bottomrule
\end{tabular}

\vspace{2pt}
\begin{minipage}{0.96\textwidth}
\footnotesize
\textcolor{red}{
\textit{Note:} Counts refer to outer-training graphs with directed
\(k\)-NN links (\(k=10\)). The graph is not symmetrized; self-loops added
within GAT layers are excluded from the reported link counts.
}
\end{minipage}
\end{table*}
\endgroup

 {
As shown in Table~\ref{tab:gnndhrs_scalability}, DHRS increased the mean
training-graph size by approximately \(52\%\)--\(56\%\). DF-1 consequently
produced the largest post-DHRS graph, containing approximately \(60{,}800\)
nodes and \(608{,}000\) directed \(k\)-NN links. The comparable proportional
increases across the three datasets indicate that absolute computational
demand is governed primarily by the original training-set size.
}

 {
Runtime was measured separately for DHRS, graph construction, and final
outer-fold model fitting. Model-fitting time excludes preprocessing,
resampling, graph construction, and inner hyperparameter search. GPU operations
were synchronized before and after timing, and peak memory was measured using
\texttt{torch.cuda.max\_memory\_allocated}. Inductive inference was evaluated
with a batch size of one after excluding warm-up runs. Reported values are
means \(\pm\) standard deviations across the outer folds.
}

\begingroup
 
\begin{table*}[!ht]
\centering
\caption{
Runtime, memory, and inductive-inference comparison under the common
experimental environment. Values are mean \(\pm\) standard deviation across
the outer folds.
}
\label{tab:runtime_baseline_comparison}
\scriptsize
\setlength{\tabcolsep}{2.4pt}
\renewcommand{\arraystretch}{1.05}
\resizebox{\textwidth}{!}{%
\begin{tabular}{@{}llcccccc@{}}
\toprule
\textbf{Dataset} &
\textbf{Model} &
\shortstack{\textbf{DHRS}\\\textbf{(s)}} &
\shortstack{\textbf{Graph}\\\textbf{(s)}} &
\shortstack{\textbf{Model fitting}\\\textbf{(s/fold)}} &
\shortstack{\textbf{Total offline}\\\textbf{(s/fold)}} &
\shortstack{\textbf{Peak GPU}\\\textbf{memory (GB)}} &
\shortstack{\textbf{Inference}\\\textbf{(ms/patient)}} \\
\midrule

\multirow{6}{*}{DF-1}
& LGBM-CW
& --- & --- & $18.4\pm1.2$ & $18.4\pm1.2$
& --- & $0.08\pm0.01$ \\
& CatB-CW
& --- & --- & $41.3\pm2.9$ & $41.3\pm2.9$
& $0.42\pm0.05$ & $0.15\pm0.02$ \\
& GCN-W
& --- & $9.8\pm0.7$ & $312.4\pm18.5$ & $322.2\pm18.5$
& $3.21\pm0.22$ & $1.85\pm0.12$ \\
& SAGE-W
& --- & $9.8\pm0.7$ & $348.7\pm21.3$ & $358.5\pm21.3$
& $3.54\pm0.25$ & $2.12\pm0.14$ \\
& GAT-CBF
& --- & $9.8\pm0.7$ & $395.2\pm24.1$ & $405.0\pm24.1$
& $4.12\pm0.28$ & $2.47\pm0.16$ \\
& \textbf{GNN-DHRS}
& $42.8\pm3.1$ & $14.6\pm1.1$ & $428.5\pm27.2$
& $485.9\pm27.4$ & $4.82\pm0.31$ & $2.61\pm0.18$ \\
\midrule

\multirow{6}{*}{DF-2}
& LGBM-CW
& --- & --- & $2.1\pm0.2$ & $2.1\pm0.2$
& --- & $0.07\pm0.01$ \\
& CatB-CW
& --- & --- & $4.8\pm0.4$ & $4.8\pm0.4$
& $0.18\pm0.03$ & $0.12\pm0.02$ \\
& GCN-W
& --- & $1.2\pm0.1$ & $41.6\pm3.2$ & $42.8\pm3.2$
& $0.85\pm0.07$ & $1.72\pm0.11$ \\
& SAGE-W
& --- & $1.2\pm0.1$ & $46.9\pm3.6$ & $48.1\pm3.6$
& $0.94\pm0.08$ & $1.98\pm0.13$ \\
& GAT-CBF
& --- & $1.2\pm0.1$ & $53.4\pm4.0$ & $54.6\pm4.0$
& $1.08\pm0.09$ & $2.31\pm0.15$ \\
& \textbf{GNN-DHRS}
& $4.9\pm0.4$ & $1.8\pm0.2$ & $58.3\pm4.1$
& $65.0\pm4.1$ & $1.15\pm0.09$ & $2.48\pm0.16$ \\
\midrule

\multirow{6}{*}{DF-3}
& LGBM-CW
& --- & --- & $2.2\pm0.2$ & $2.2\pm0.2$
& --- & $0.07\pm0.01$ \\
& CatB-CW
& --- & --- & $5.1\pm0.4$ & $5.1\pm0.4$
& $0.19\pm0.03$ & $0.13\pm0.02$ \\
& GCN-W
& --- & $1.3\pm0.1$ & $43.8\pm3.4$ & $45.1\pm3.4$
& $0.88\pm0.07$ & $1.75\pm0.11$ \\
& SAGE-W
& --- & $1.3\pm0.1$ & $49.2\pm3.7$ & $50.5\pm3.7$
& $0.97\pm0.08$ & $2.01\pm0.13$ \\
& GAT-CBF
& --- & $1.3\pm0.1$ & $56.1\pm4.2$ & $57.4\pm4.2$
& $1.11\pm0.09$ & $2.35\pm0.15$ \\
& \textbf{GNN-DHRS}
& $5.1\pm0.4$ & $1.9\pm0.2$ & $61.7\pm4.3$
& $68.7\pm4.3$ & $1.18\pm0.10$ & $2.52\pm0.17$ \\
\bottomrule
\end{tabular}%
}

\vspace{2pt}
\begin{minipage}{0.98\textwidth}
\footnotesize
\textit{Note:} CW: class weighting; W: weighted BCE; CBF: class-balanced focal
loss. Dashes indicate inapplicable operations. Model-fitting time excludes
DHRS, graph construction, and hyperparameter search.
\end{minipage}
\end{table*}
\endgroup

 {
Table~\ref{tab:runtime_baseline_comparison} shows that the tabular models
required the least computation because they did not perform graph construction
or message passing. On DF-1, GNN-DHRS increased model-fitting time by
approximately \(8.4\%\) relative to GAT-CBF. After including DHRS and graph
construction, total offline time increased from \(405.0\) to \(485.9\)
seconds per fold, while peak GPU-memory allocation increased from \(4.12\) to
\(4.82\) GB. Similar relative patterns were observed on DF-2 and DF-3.
}

 {
The additional computational demand was accompanied by higher observed
minority-sensitive performance. On DF-1, GNN-DHRS exceeded GAT-CBF by
\(5.42\), \(4.47\), and \(4.16\) percentage points in Recall, F1-score, and
PR-AUC, respectively. This comparison describes the empirical
efficiency--performance trade-off between the complete configurations and
should not be interpreted as isolating the causal contribution of an
individual component.
}

 {
Strictly inductive inference required no resampling, model refitting, or
test--test graph construction. Each unseen patient was connected only to its
\(k=10\) nearest original training nodes. Mean per-patient latency ranged from
\(2.48\) to \(2.61\) ms, indicating that most computational demand occurred
during offline resampling, graph preparation, and model fitting.
}

 {
Overall, GNN-DHRS remained executable on a single 16-GB GPU, including the
expanded DF-1 training graph. Its offline cost and memory requirements were
higher than those of the controlled graph baselines but remained within the
available hardware capacity. These measurements establish computational
feasibility in the reported experimental environment; deployment-specific
profiling remains necessary for larger cohorts, different graph densities,
concurrent requests, and alternative clinical infrastructure.
}%----------------------------------------
\subsection{Robustness Under Realistic Perturbations}
\label{subsec:robustness_perturbations}

 {
This section reports the four stress tests defined in
Section~\ref{subsec:experimental_setup_results}. All perturbations were
applied exclusively to the untouched outer-test patients under the strictly
inductive protocol; the training partitions, the trained model, and the
training graph remained frozen throughout.
}

\begingroup
 
\begin{table*}[!ht]
\centering
\caption{
Performance of GNN-DHRS under increasing outer-test missingness. Values are
mean $\pm$ standard deviation across the outer folds.
}
\label{tab:robustness_missingness}
\small
\setlength{\tabcolsep}{5pt}
\renewcommand{\arraystretch}{1.08}
\begin{tabular}{@{}lcccc@{}}
\toprule
\textbf{Dataset} & \textbf{Missingness} &
\shortstack{\textbf{Recall}\\\textbf{(\%)}} &
\shortstack{\textbf{F1-score}\\\textbf{(\%)}} &
\shortstack{\textbf{PR-AUC}\\\textbf{(\%)}} \\
\midrule
\multirow{4}{*}{DF-1}
& 0\%  & $99.31\pm0.15$ & $99.39\pm0.08$ & $99.87\pm0.04$ \\
& 10\% & $99.05\pm0.19$ & $99.16\pm0.15$ & $99.62\pm0.11$ \\
& 20\% & $98.42\pm0.28$ & $98.61\pm0.24$ & $99.08\pm0.19$ \\
& 30\% & $97.38\pm0.41$ & $97.72\pm0.35$ & $98.21\pm0.29$ \\
\midrule
\multirow{4}{*}{DF-2}
& 0\%  & $98.96\pm0.24$ & $98.39\pm0.17$ & $98.71\pm0.19$ \\
& 10\% & $98.31\pm0.31$ & $97.68\pm0.27$ & $97.95\pm0.24$ \\
& 20\% & $97.02\pm0.44$ & $96.35\pm0.39$ & $96.48\pm0.35$ \\
& 30\% & $95.18\pm0.61$ & $94.29\pm0.55$ & $94.12\pm0.49$ \\
\midrule
\multirow{4}{*}{DF-3}
& 0\%  & $99.12\pm0.16$ & $98.52\pm0.41$ & $98.84\pm0.22$ \\
& 10\% & $98.47\pm0.29$ & $97.85\pm0.26$ & $98.05\pm0.22$ \\
& 20\% & $97.15\pm0.42$ & $96.48\pm0.37$ & $96.61\pm0.33$ \\
& 30\% & $95.34\pm0.58$ & $94.41\pm0.52$ & $94.28\pm0.47$ \\
\bottomrule
\end{tabular}

\vspace{2pt}
\begin{minipage}{0.94\linewidth}
\footnotesize
\textit{Note:} Missing values were re-imputed using training-fold statistics
prior to inference; the trained model and training graph were unchanged.
\end{minipage}
\end{table*}
\endgroup

 {
As reported in Table~\ref{tab:robustness_missingness}, Recall declined
gradually with missingness. At 10\% missingness, Recall decreased by 0.26,
0.65, and 0.65 percentage points on DF-1, DF-2, and DF-3, respectively; at
30\% missingness, the reductions were 1.93, 3.78, and 3.78 percentage points.
No abrupt performance collapse was observed at any evaluated rate.
}

\begingroup
 
\begin{table*}[!ht]
\centering
\caption{
Performance of GNN-DHRS under additive Gaussian feature noise applied to
standardized continuous predictors of outer-test patients. Values are mean
$\pm$ standard deviation across the outer folds.
}
\label{tab:robustness_noise}
\small
\setlength{\tabcolsep}{5pt}
\renewcommand{\arraystretch}{1.08}
\begin{tabular}{@{}lcccc@{}}
\toprule
\textbf{Dataset} & \textbf{Noise level $\sigma$} &
\shortstack{\textbf{Recall}\\\textbf{(\%)}} &
\shortstack{\textbf{F1-score}\\\textbf{(\%)}} &
\shortstack{\textbf{PR-AUC}\\\textbf{(\%)}} \\
\midrule
\multirow{4}{*}{DF-1}
& 0.00 & $99.31\pm0.15$ & $99.39\pm0.08$ & $99.87\pm0.04$ \\
& 0.05 & $99.02\pm0.17$ & $99.08\pm0.14$ & $99.61\pm0.10$ \\
& 0.10 & $98.35\pm0.26$ & $98.47\pm0.22$ & $99.02\pm0.17$ \\
& 0.20 & $96.78\pm0.39$ & $97.02\pm0.34$ & $97.56\pm0.28$ \\
\midrule
\multirow{4}{*}{DF-2}
& 0.00 & $98.96\pm0.24$ & $98.39\pm0.17$ & $98.71\pm0.19$ \\
& 0.05 & $98.42\pm0.28$ & $97.81\pm0.25$ & $98.15\pm0.22$ \\
& 0.10 & $97.15\pm0.38$ & $96.42\pm0.34$ & $96.58\pm0.31$ \\
& 0.20 & $94.62\pm0.55$ & $93.71\pm0.50$ & $93.85\pm0.46$ \\
\midrule
\multirow{4}{*}{DF-3}
& 0.00 & $99.12\pm0.16$ & $98.52\pm0.41$ & $98.84\pm0.22$ \\
& 0.05 & $98.61\pm0.26$ & $97.98\pm0.23$ & $98.24\pm0.21$ \\
& 0.10 & $97.38\pm0.36$ & $96.65\pm0.32$ & $96.82\pm0.29$ \\
& 0.20 & $94.85\pm0.52$ & $93.94\pm0.48$ & $94.08\pm0.44$ \\
\bottomrule
\end{tabular}
\end{table*}
\endgroup

 {
Table~\ref{tab:robustness_noise} shows the expected monotonic decline as the
noise level increases. At $\sigma=0.05$, Recall decreased by no more than
0.58 percentage points across the three datasets; at $\sigma=0.20$, Recall
decreased by 2.53, 4.34, and 4.27 percentage points on DF-1, DF-2, and DF-3,
respectively.
}

\begingroup
 
\begin{table*}[!ht]
\centering
\caption{
Performance of GNN-DHRS under simulated stroke-prevalence shift, obtained by
stratified subsampling of the untouched outer-test fold. Values are mean
$\pm$ standard deviation across the outer folds.
}
\label{tab:robustness_prevalence}
\small
\setlength{\tabcolsep}{4.5pt}
\renewcommand{\arraystretch}{1.08}
\begin{tabular}{@{}llccccc@{}}
\toprule
\textbf{Dataset} & \textbf{Prevalence} &
\shortstack{\textbf{Precision}\\\textbf{(\%)}} &
\shortstack{\textbf{Recall}\\\textbf{(\%)}} &
\shortstack{\textbf{F1-score}\\\textbf{(\%)}} &
\shortstack{\textbf{PR-AUC}\\\textbf{(\%)}} &
\shortstack{\textbf{ROC-AUC}\\\textbf{(\%)}} \\
\midrule
\multirow{3}{*}{DF-1}
& Half (0.9\%)     & $98.86\pm0.21$ & $99.24\pm0.18$ & $99.05\pm0.16$ & $99.71\pm0.09$ & $99.98\pm0.02$ \\
& Original (1.8\%) & $99.48\pm0.08$ & $99.31\pm0.15$ & $99.39\pm0.08$ & $99.87\pm0.04$ & $100.00\pm0.00$ \\
& Double (3.6\%)   & $99.72\pm0.09$ & $99.28\pm0.14$ & $99.50\pm0.08$ & $99.91\pm0.04$ & $100.00\pm0.00$ \\
\midrule
\multirow{3}{*}{DF-2}
& Half (2.5\%)     & $95.62\pm0.38$ & $98.81\pm0.25$ & $97.19\pm0.29$ & $97.85\pm0.24$ & $99.87\pm0.06$ \\
& Original (5.0\%) & $97.84\pm0.18$ & $98.96\pm0.24$ & $98.39\pm0.17$ & $98.71\pm0.19$ & $99.92\pm0.08$ \\
& Double (10.0\%)  & $98.95\pm0.16$ & $98.87\pm0.21$ & $98.91\pm0.14$ & $99.24\pm0.12$ & $99.94\pm0.05$ \\
\midrule
\multirow{3}{*}{DF-3}
& Half (2.45\%)    & $95.84\pm0.41$ & $98.95\pm0.24$ & $97.37\pm0.31$ & $98.02\pm0.27$ & $99.89\pm0.06$ \\
& Original (4.9\%) & $97.92\pm0.64$ & $99.12\pm0.16$ & $98.52\pm0.41$ & $98.84\pm0.22$ & $99.93\pm0.07$ \\
& Double (9.8\%)   & $99.02\pm0.15$ & $99.05\pm0.19$ & $99.03\pm0.13$ & $99.31\pm0.11$ & $99.95\pm0.04$ \\
\bottomrule
\end{tabular}

\vspace{2pt}
\begin{minipage}{0.96\textwidth}
\footnotesize
\textit{Note:} Prevalence was varied only by subsampling the outer-test fold;
training, graph construction, and the decision threshold were unchanged.
\end{minipage}
\end{table*}
\endgroup

 {
Table~\ref{tab:robustness_prevalence} confirms the expected prevalence
sensitivity of Precision and PR-AUC, while ROC-AUC remained nearly
unchanged. At half prevalence, Precision and F1-score decreased (for example,
from 97.84\% to 95.62\% and from 98.39\% to 97.19\% on DF-2), whereas at
double prevalence, Precision, F1-score, and PR-AUC increased above their
original-prevalence values. Recall remained comparatively stable across all
evaluated prevalence levels, consistent with the theoretical prevalence
invariance of threshold-independent ranking measures.
}

\begingroup
 
\begin{table*}[!ht]
\centering
\caption{
Performance of GNN-DHRS under corrupted inductive test-to-training edges.
Values are mean $\pm$ standard deviation across the outer folds.
}
\label{tab:robustness_edge_corruption}
\small
\setlength{\tabcolsep}{5pt}
\renewcommand{\arraystretch}{1.08}
\begin{tabular}{@{}lcccc@{}}
\toprule
\textbf{Dataset} & \textbf{Edge corruption $\rho$} &
\shortstack{\textbf{Recall}\\\textbf{(\%)}} &
\shortstack{\textbf{F1-score}\\\textbf{(\%)}} &
\shortstack{\textbf{PR-AUC}\\\textbf{(\%)}} \\
\midrule
\multirow{4}{*}{DF-1}
& 0\%  & $99.31\pm0.15$ & $99.39\pm0.08$ & $99.87\pm0.04$ \\
& 10\% & $98.94\pm0.21$ & $99.02\pm0.18$ & $99.54\pm0.13$ \\
& 20\% & $98.12\pm0.33$ & $98.28\pm0.29$ & $98.95\pm0.22$ \\
& 30\% & $96.85\pm0.47$ & $97.15\pm0.42$ & $97.68\pm0.35$ \\
\midrule
\multirow{4}{*}{DF-2}
& 0\%  & $98.96\pm0.24$ & $98.39\pm0.17$ & $98.71\pm0.19$ \\
& 10\% & $98.05\pm0.32$ & $97.42\pm0.28$ & $97.81\pm0.26$ \\
& 20\% & $96.48\pm0.46$ & $95.71\pm0.41$ & $96.02\pm0.38$ \\
& 30\% & $93.72\pm0.62$ & $92.85\pm0.57$ & $93.15\pm0.52$ \\
\midrule
\multirow{4}{*}{DF-3}
& 0\%  & $99.12\pm0.16$ & $98.52\pm0.41$ & $98.84\pm0.22$ \\
& 10\% & $98.24\pm0.30$ & $97.58\pm0.27$ & $97.94\pm0.24$ \\
& 20\% & $96.71\pm0.44$ & $95.94\pm0.39$ & $96.28\pm0.35$ \\
& 30\% & $94.02\pm0.59$ & $93.15\pm0.54$ & $93.48\pm0.49$ \\
\bottomrule
\end{tabular}
\end{table*}
\endgroup

 {
Table~\ref{tab:robustness_edge_corruption} shows that Recall decreased by
0.37 to 0.91 percentage points at 10\% edge corruption and by 2.46, 5.24, and
5.10 percentage points on DF-1, DF-2, and DF-3, respectively, at 30\% edge
corruption, with F1-score and PR-AUC degrading proportionally.
}

\begin{figure*}[!ht]
\centering
\includegraphics[width=0.98\textwidth]
{images/robustness_stress_tests_summary.png}
\caption{ {
Degradation curves of GNN-DHRS under the four robustness stress tests across
DF-1, DF-2, and DF-3: (a) missingness rate, (b) Gaussian feature-noise level,
(c) simulated stroke prevalence, and (d) inductive edge-corruption rate. All
perturbations were applied only to the untouched outer-test patients under
the frozen trained model and training graph.
}}
\label{fig:robustness_summary}
\end{figure*}

 {
Considered jointly, Figure~\ref{fig:robustness_summary} and
Tables~\ref{tab:robustness_missingness}--\ref{tab:robustness_edge_corruption}
show a graceful, monotonic decline in performance under every stress test
rather than an abrupt collapse. DF-1 was consistently the most stable
dataset, plausibly reflecting its larger sample size, whereas DF-2 and DF-3
were comparatively more sensitive to missingness, feature noise, and edge
corruption, consistent with their smaller sample sizes and higher baseline
imbalance ratios. These results support the robustness of GNN-DHRS under
realistic data-quality degradation, while indicating that performance
guarantees should be interpreted relative to the severity of the
perturbation encountered in a given deployment setting.
}




%==============================================================
%--------------------------------
% Discussion
%--------------------------------
\section{Discussion}
\label{sec:discussion}

 {
The results indicate that GNN-DHRS achieves strong within-benchmark
performance for imbalanced stroke-status classification by combining
training-only hybrid resampling with patient-similarity graph modeling and
multi-head graph attention. Unlike conventional tabular models that process
patients independently, GNN-DHRS incorporates relationships among similar
patient profiles. The controlled comparisons suggest that this relational
representation provides information complementary to individual patient
features under the adopted evaluation protocol, although they do not
establish a causal explanation for the observed performance differences.
}

 {
The component-wise ablations suggest complementary contributions from
minority enhancement, neighborhood cleaning, patient-graph construction,
attention aggregation, and cosine-similarity edge weighting. Improvements in
minority-sensitive and class-balanced measures were observed on DF-1, DF-2,
and DF-3 rather than being confined to the largest dataset. These findings
support the combined framework but should not be interpreted as proving a
definitive causal mechanism for its performance.
}

{\color{red}
The strong within-benchmark performance should be considered in relation to
the scope and characteristics of the evaluated public datasets. Although the
adopted leakage-control procedures strengthen the reliability of the reported
estimates, these datasets provide only a limited representation of the
variability that may be encountered across different patient populations and
clinical settings. Factors that are not captured by the available predictors
may also influence stroke status and affect the reproducibility of the
observed performance beyond the evaluated benchmarks.

The lower cross-benchmark performance, particularly for
DF-2$\rightarrow$DF-1 and DF-3$\rightarrow$DF-1, further indicates that model
behavior is affected by differences in feature distributions and by the
patient-similarity structure induced within each dataset. These findings do
not by themselves demonstrate statistical overfitting, but they emphasize
that the reported within-benchmark results should be interpreted within the
boundaries of the evaluated data. Validation on independently collected,
multi-center, and prospective cohorts is therefore necessary before broader
clinical applicability can be established.
}

 {
The explainability results provide two complementary descriptions of model
behavior. Kernel SHAP produced feature-attribution patterns broadly consistent
with established stroke risk factors, including age, average glucose level,
BMI, hypertension, heart disease, and smoking-related variables.
Attention-based relational analysis further identified a restricted set of
patient connections that strongly influenced the fitted predictions. The
larger degradation following high-attention edge removal suggests that these
relations are functionally relevant to the model. Nevertheless, agreement with
recognized risk factors and attention-based fidelity indicate plausibility and
model-level relevance only; they do not establish causal effects, clinical
validity, or clinically verified relationships among patients.
} \textcolor{red}{
Moreover, because GNN-DHRS predictions are conditioned on graph-based
neighborhood aggregation, Kernel-SHAP feature attributions should not be
interpreted as isolated or causal effects of individual clinical variables.
They describe feature contributions within a predictive function in which
patient-level information and relational context are jointly represented.
}

From a clinical perspective, sensitivity to stroke-positive cases is important
because false-negative classifications may reduce the usefulness of a
screening-oriented system.
 {
Recall was therefore interpreted together with Precision, PR-AUC, specificity,
NPV, balanced accuracy, and MCC. Moreover, the available datasets contain
retrospective stroke-status labels rather than verified longitudinal outcomes.
Accordingly, the results do not demonstrate prediction of future stroke events
or improvement in patient outcomes; instead, they motivate further evaluation
of GNN-DHRS for potential decision-support applications.
}

 {
Several data-related limitations affect the interpretation of the results.
DF-1, DF-2, and DF-3 are related public benchmark files rather than
independently collected external clinical cohorts and may not represent the
diversity of real clinical populations. \textcolor{red}{
The within-file record-integrity audit did not identify exact duplicate
predictor--target profiles or redundant records in the final benchmark files.
However, the repositories do not provide verified longitudinal patient
identities or a common identifier across datasets. Consequently, latent
patient-level overlap within or between benchmark files cannot be excluded,
and the cross-benchmark analyses should not be interpreted as independent
external clinical replication.
}
}

 {
DHRS introduces additional limitations associated with synthetic data.
Borderline-SMOTE interpolates minority observations in the transformed feature
space and may generate combinations that are mathematically valid but
clinically uncommon, particularly for sparse categorical patterns or
correlated predictors. SMOTEENN may remove informative boundary observations
together with noisy samples, thereby altering both class composition and graph
topology. Because synthetic and original training nodes are treated
identically during graph learning, synthetic observations may influence the
model through their training contributions and incident graph edges. The
mixed-data sensitivity and component-ablation experiments reduce concern about
substantial distortion under the evaluated conditions but cannot verify the
clinical plausibility of every generated patient profile.
}

\textcolor{red}{
The controlled synthetic-node ablation additionally distinguishes
data-level resampling from synthetic-node message propagation. Because
generated observations were retained during optimization while their incident
graph edges were selectively removed, differences between the complete and
synthetic-edge-disabled configurations quantify the additional contribution
associated with synthetic-node relational connectivity. This analysis reduces
ambiguity concerning the source of the performance gain, although it does not
establish that generated observations or their graph relations correspond to
clinically valid patient states.
}

 {
The patient graph introduces a related source of potential failure. Cosine
similarity may connect patients who are close within the selected feature
space but differ in unmeasured clinical characteristics, allowing misleading
neighborhood information to propagate during message passing. Performance may
also decline when important predictors are unavailable, unseen categories
occur, missingness or measurement error exceeds the tested levels, or the
deployment population differs substantially in prevalence, demographic
composition, clinical practice, or data-acquisition procedures. Population
shift may additionally alter probability calibration and invalidate a
threshold selected from the development data.
}

 {
The robustness experiments showed gradual degradation under the evaluated
levels of missingness, feature noise, prevalence shift, and edge corruption.
These findings indicate stability within the specified perturbation ranges but
do not establish robustness to more severe or naturally occurring failures.
Simulated perturbations cannot reproduce all dependencies, coding practices,
recording errors, and measurement processes encountered in clinical systems.
Evaluation using independently collected, naturally noisy data is therefore
necessary.
}

 {
The computational evaluation showed that most additional cost arises from
offline resampling, graph construction, and model fitting. The expanded
training graphs remained executable on the reported 16-GB GPU, with peak
memory below \(5\)~GB and inductive inference of approximately
\(2.5\)~ms per patient. These measurements demonstrate feasibility only in
the reported experimental environment and exclude data integration,
concurrent workloads, monitoring, and clinical infrastructure costs.
Deployment-specific profiling remains necessary, particularly for larger or
denser patient graphs.
}

Future work should evaluate GNN-DHRS using longitudinal, independently
collected multi-center cohorts and richer clinical information, including
electronic health records, imaging biomarkers, and multimodal variables.
Dynamic or heterogeneous graph structures may also provide more informative
representations than static feature-space similarity.
 {
Before clinical use is considered, further studies should include prospective
validation, expert assessment of synthetic samples and explanations, subgroup
and fairness analyses, calibration monitoring, site-specific threshold
assessment, and evaluation of the framework's effects on clinical decisions
and patient outcomes.
}
%==========================================================
\section{Conclusion and Future Work}
\label{sec:conclusion}

This paper presented GNN-DHRS, an explainable graph-learning framework for
imbalanced stroke-status classification. The framework combines training-only
dual hybrid resampling, patient-similarity graph construction, multi-head
graph attention, and feature- and relation-level interpretation. Evaluation
on three public benchmark datasets produced consistently strong
minority-sensitive and class-balanced performance. The controlled ablation,
sensitivity, robustness, and statistical analyses suggest that DHRS and
weighted graph-attention learning provide complementary contributions under
the adopted experimental protocol. Kernel SHAP and attention-based analyses
additionally produced feature-attribution and relational patterns that were
reproducible across folds and broadly consistent with recognized stroke risk
factors.

 {
These results represent within-benchmark evidence obtained using leakage-safe
nested cross-validation and strictly inductive outer-test inference. Because
the datasets contain retrospective stroke-status labels and are not verified
independent prospective cohorts, the findings do not establish prediction of
future stroke events, clinical benefit, or deployment readiness. They instead
support further investigation of GNN-DHRS for potential decision-support
applications.
}

Future work will focus on two key fronts. First, to improve representation, we plan to extend the model to multi-center longitudinal cohorts by incorporating multimodal EHR data and dynamic graph architectures. Second, before clinical translation, prospective validation must establish external generalizability, probability calibration, and fairness across patient subgroups, while evaluating the framework's real-world impact on clinical decision-making.
%==========================================================
\section*{Declarations}
\noindent \textbf{Data Availability:}\\
\textcolor{red}{
The DF-1, DF-2, and DF-3 datasets used in this study are publicly available
from the repositories cited in
\cite{dataset_df1,dataset_df2,dataset_df3}.
The corresponding references provide the repository URLs and access dates.
No new patient data were collected or generated in this study.
} \\

\noindent \textbf{Code Availability:}\\
 {
The source code developed for the proposed GNN-DHRS framework, including
data preprocessing, DHRS resampling, graph construction, model training,
evaluation, and explainability analysis, is available from the corresponding
author upon reasonable request. https://github.com/HanyEjust/GNN-DHRS
} \\


% \noindent \textbf{Funding Declaration}: Not applicable.\\

\noindent \textbf{Competing Interests}:  \\
The authors declare that they have no known competing financial interests or personal relationships that could have appeared to influence the work reported in this paper. \\

\noindent \textbf{Ethics declaration}: Not applicable.
%============================================================
% \section*{Declarations}
% ============================================================

% \subsection*{Funding}

% The authors would like to thank the Research Chair of Prince Dr. Faisal bin Mishaal bin Saud bin Abdulaziz for Artificial Intelligence (CPFAI) at Qassim University for Supporting this research work (QU-CPFAI-25HC213).

% \subsection*{Conflict of Interest / Competing Interests}

% The authors declare that they have no known competing financial interests or personal relationships that could have appeared to influence the work reported in this paper.

% \subsection*{Ethics Approval and Consent to Participate}

% This study was conducted using publicly available and/or anonymized secondary datasets. No direct interaction with human participants was performed by the authors, and no personally identifiable information was used. Therefore, formal ethical approval and informed consent were not required. If required by the target journal, this statement should be revised according to the ethical approval status of the original data source.

% \subsection*{Consent for Publication}

% Not applicable. The study does not contain identifiable individual-level personal data, images, or clinical information requiring consent for publication.

% \subsection*{Data Availability}

% The datasets used in this study are publicly available from their original repositories/sources. The processed datasets and experimental outputs can be made available from the corresponding author upon reasonable request, subject to the data-use conditions of the original dataset providers.

% \subsection*{Materials Availability}

% Not applicable. This study did not generate or use physical materials.

% \subsection*{Code Availability}

% The source code developed for the proposed GNN-DHRS framework, including preprocessing, DHRS-based graph-guided resampling, graph construction, model training, evaluation, and explainability analysis, can be made available from the corresponding author upon reasonable request. A public repository may be provided after manuscript acceptance.

% \subsection*{Author Contributions}

% All authors contributed to the conception and design of the study. The methodology, model development, implementation, experiments, statistical analysis, and visualization were performed by the authors. The first draft of the manuscript was prepared by the authors, and all authors reviewed, edited, and approved the final version of the manuscript. 



\bibliography{sn-bibliography}% common bib file
%% if required, the content of .bbl file can be included here once bbl is generated
%%\input sn-article.bbl

\end{document}
